\section{The mapping-cone construction and its scope}\label{sec:B}
This section checks the supplied version of Ben-Bassat's manuscript
\cite{BenBassat2026B}.  Its
numbering is kept distinct from the numbering of the present paper.  The
symbols $\mathbf L$, $N(L)$, and $\operatorname{Kar}(L)$ have the conventions
stated in that manuscript: composition in $\mathbf L$ is multiplication,
whereas $N(L)=N(\mathbf L^{\mathrm{op}})$ has the usual ordered bar labels.
All maps between these realizations use reversal of simplices explicitly.

\subsection{A necessary correction to the resolution request}
\begin{proposition}\label{prop:B-finite-resolution-obstruction}
For either eleven-element monoid of Example~3.4 in the supplied manuscript,
the trivial integral module has no projective resolution of finite length.
Assuming Theorem~A, their positive-degree integral homology is as follows:
\[
\begin{array}{c|ccc}
 &H_1&H_2&H_q\ (q\ge3)\\ \hline
 M(S^1_\Delta,f_1)&\mathbb Z/2&\mathbb Z&
 \mathbb Z/2\text{ for odd }q,\quad0\text{ for even }q\\
 M(S^1_\Delta,f_g)&0&\mathbb Z&
 \mathbb Z/2\text{ for odd }q,\quad0\text{ for even }q
\end{array}
\]
In both rows $H_0=\mathbb Z$.  Thus ``finite projective resolution'' in
Part~B2 cannot mean a resolution of finite length in these two cases.
\end{proposition}
\begin{proof}
The constant labelled map has homotopy cofiber $BC_2\vee S^2$.
For the nonconstant label the cofiber is $\mathbb{RP}^{\infty}/
\mathbb{RP}^{1}$, since $\mathbb{RP}^{1}\hookrightarrow
\mathbb{RP}^{\infty}$ is a CW-subcomplex inclusion.  The cellular boundary
of real projective space is multiplication by $2$ in positive even degree
and zero in odd degree.  Passing to the quotient removes degrees $0$ and
$1$ from the reduced chain complex and changes the differential out of
degree $2$ to zero.  The displayed groups follow in every degree.  A
projective resolution of length $d$ makes
$\operatorname{Tor}^{\mathbb ZM}_q(\mathbb Z,\mathbb Z)=0$ for every $q>d$,
in contradiction with the nonzero odd-degree groups.  The constructions
and the proof of Theorem~A are checked below; this argument does not infer
infinite tails from any finite bar calculation.
\end{proof}

\subsection{Step-by-step check of the cofiber theorem}
The prompt's ``Lemma~2.3'' is Definition~2.3 in the supplied PDF, and its
``Example~3.2'' is Example~3.4 there; Lemma~3.2 is the normal-form lemma.
We check the intended statements with these numbering corrections.

\begin{lemma}[Checks of Lemmas 2.1--2.2 and Definition 2.3]
\label{lem:B-elementary}
The cofiber-invariance, gluing, and collage statements in the supplied
manuscript hold with the stated hypotheses.
\end{lemma}
\begin{proof}
\emph{Lemma 2.1.}  Use the double mapping cylinder
$Z\amalg (X\times[0,1])\amalg *$, identifying $(x,0)$ with $f(x)$
and $(x,1)$ with the cone vertex.  A specified homotopy
$vf\simeq f'u$ supplies a map to the analogous cylinder for $f'$: on the
first half of each cylinder traverse this homotopy, and on the second
half use $u$.  Choosing homotopy inverses of $u$ and $v$ gives an inverse
up to homotopy.  The homotopies are defined on the cylinders by inserting
or deleting traversals of the chosen homotopies, so they respect their
end identifications.  This proves homotopy invariance, including the
compatibility with the map out of $Z$.

\emph{Lemma 2.2.}  The comparison from the double mapping cylinder to
$X\cup_A B$ is a homotopy equivalence when one attaching leg is a
cofibration: the homotopy extension property straightens the cylinder
along that leg.  Replacing $B$ by $A$ through a homotopy inverse of $g$
therefore identifies this double mapping cylinder with that for
$X\leftarrow A\xrightarrow{\mathrm{id}}A$, which retracts to $X$.
This proves that the actual canonical map $X\to X\cup_A B$ is an
equivalence.  For the applications, every space is a CW realization and
every cited cofibration comes from a simplicial monomorphism, so the
homotopy extension hypothesis holds.

\emph{Definition 2.3.}  In a collage every composable string either lies
entirely in one side, or crosses from $\mathcal A$ to $\mathcal B$ once.
The first two cases inherit associativity.  In a string with one crossing,
associativity on either side is functoriality of $\Pi$, and compatibility
of the left and right actions is the equality of the two commuting
actions of a functor on $\mathcal A^{\mathrm{op}}\times\mathcal B$.
The identity laws follow in the same way.  The equation defining a
morphism of $\operatorname{El}(\Pi)$, namely
$x'\alpha=\beta x$, is preserved under composition because
$x''\alpha'\alpha=\beta'\beta x$; identities also satisfy it.
Thus the category of elements and both projection functors are defined
as claimed, with no invertibility assumption on the actions.
\end{proof}

\begin{lemma}[Checks of Lemmas 2.4--2.6]\label{lem:B-prelim}
The retract-generator, collage-pushout, and truncation lemmas in the
supplied manuscript hold.  The collage equivalence preserves both
structure maps.
\end{lemma}
\begin{proof}
\emph{Retract-generator lemma.}  For the inclusion
$i:\operatorname{End}_{\mathcal D}(c_0)\to\mathcal D$, the objects of
$(i\downarrow c)$ are maps $c_0\to c$.  If
$c\xrightarrow{s}c_0\xrightarrow{r}c$ has $rs=\mathrm{id}_c$,
postcomposition with $s$ and $r$ exhibits this comma category as a
retract of $(i\downarrow c_0)$.  The latter has terminal object
$\mathrm{id}_{c_0}$.  After realization, composing its contraction with
the retraction gives a contraction of the former.  Quillen's Theorem~A
therefore applies to every $c$.  This also proves the Karoubi-envelope
comparison used subsequently; it does not treat a nonunital corner
inclusion as a monoid homomorphism.

\emph{Collage-pushout lemma.}  Put $E=\operatorname{El}(\Pi)$ and form
the Grothendieck construction $T$ of the span
$\mathcal A\leftarrow E\to\mathcal B$.  Its objects are the disjoint
union of these three sets of objects.  In addition to the arrows internal
to the three categories, there are arrows $e\to a'$ represented by
$s(e)\to a'$ and arrows $e\to b'$ represented by $t(e)\to b'$.
The functor $F:T\to\mathcal C(\Pi)$ fixes $\mathcal A$ and
$\mathcal B$, sends $(a,x,b)$ to $a$, and sends its arrow to $b'$
represented by $\beta$ to $\beta x$.  The equation
$x'\alpha=\beta x$ proves preservation of every mixed composite.
For $b\in\mathcal B$, $(b\downarrow F)$ has an initial object.
For $a\in\mathcal A$, the comma category is another such span
construction, with left category $(a\downarrow\mathcal A)$,
right category the elements $(b,y)$ of $\Pi(a,-)$, and middle category
$((a',x,b),\alpha:a\to a')$.  The functor from the middle to the right
sends this object to $(b,x\alpha)$.  Its section chooses
$((a,y,b),\mathrm{id}_a)$, and $(\alpha,\mathrm{id}_b)$ gives a
natural transformation from section-after-projection to the identity.
The right leg is therefore a homotopy equivalence and the left vertex is
contractible.  Their homotopy pushout is contractible.

Here the instance of the homotopy-colimit theorem being used is that
for any small indexing category $I$ and functor $D:I\to\mathbf{Cat}$,
$B(\int_I D)$ is naturally weakly equivalent to
$\operatorname{hocolim}_I BD$.  This is Thomason's Theorem~1.2
\cite{Thomason1979B}; the independently checked primary proof is
Theorem~26.8 and Corollary~26.9 of \cite{ChacholskiScherer2002B}.
For a span this is the double mapping
cylinder; the equivalence is natural on its two endpoint inclusions.
Applying this once to the comma span and then to $T$, Quillen's
Theorem~A gives the claimed equivalence, with both endpoint maps
preserved.  No flatness of a biset is a hypothesis of this statement.

\emph{Truncation lemma.}  Fix an arbitrary simplex $\sigma\in Y_m$.
In the comma category from $([m],\sigma)$ to the truncated simplex
category, an object is a factorization $\sigma=\theta^*\tau$ with
$\theta:[m]\to[k]$, $k\le n$.  Replace $\theta$ by its surjection
onto its ordered image, pulling $\tau$ back along the inclusion of that
image.  This is functorial: a commuting triangle of ordinal maps induces
the corresponding map of their ordered images.  The image inclusions
give a natural transformation from this retraction to the identity.
In the subcategory of surjections, the unique expression of $\sigma$
as a degeneracy of a nondegenerate simplex gives a terminal object.
The uniqueness is essential: if $\rho:[m]\twoheadrightarrow[k]$
and $\sigma=\rho^*\tau$, the degeneracy of $\tau$ to its
nondegenerate ancestor is the unique arrow to that terminal object,
since $\rho$ is surjective.  Its dimension is at most $n$.
Quillen's Theorem~A gives the truncated-to-full equivalence.
The full last-vertex comparison is the usual realization equivalence
from the simplicial replacement of the diagram of simplices of $Y$:
on each simplex its source contracts to that simplex's last vertex,
and these contractions agree with the simplicial replacement maps.
Equivalently, attach the nondegenerate simplices in increasing dimension;
the comparison on a simplex and on its already attached boundary is an
equivalence, so the gluing lemma proves the assertion by induction.
Its restriction factors as the just-proved inclusion followed by this
comparison, and is consequently an equivalence as well.
\end{proof}

\begin{lemma}[Checks of the algebraic construction and Lemma 3.6]
\label{lem:B-monoid}
Definition~3.3 defines an associative unital finite monoid of the stated
order.  Its corners $e_iM(Y,f)e_i$ are isomorphic to $L$ as unital
monoids, and every object of the retraction category is a retract of
$[n]$.
\end{lemma}
\begin{proof}
For a monotone map $\theta:[j]\to[k]$, the product of labels along an
interval in $\theta^*y$ is the product along the corresponding interval
in $y$.  Splitting the latter interval at $\theta(j)$ gives
\[
 \operatorname{hol}_{\psi(i)}(\theta^*y)
 \operatorname{hol}_{\theta(j)}(y)
 =\operatorname{hol}_{\theta\psi(i)}(y).
\]
This proves the simplicial identities for
$\theta^*(y,l)=(\theta^*y,\operatorname{hol}_{\theta(j)}(y)l)$.
Right multiplication by $L$ commutes with them by associativity in $L$.

To check the normal-form assertion, first compose all consecutive old
collage arrows.  Between any two occurrences of $\rho$ lies a map
$\bullet_L\to[0]$, necessarily and uniquely $s_0a$ for $a\in L$.
Replace $\rho s_0a$ by $a$.  This shortens the number of $\rho$'s;
disjoint reductions commute, and a reduction next to an old composite
is unchanged because $s_0(ab)=(s_0a)b$.  Thus reduction terminates and
is independent of order.  A surviving $\rho$ has to have a vertex
map $i^{\mathrm{op}}$ immediately to its right; on its left is either
$a\in L$ or a unique old arrow $(y,l)$.  These are precisely the
normal forms in Lemma~3.2.  In particular, no new relation identifies
old arrows with arrows containing $\rho$.

Multiplying these forms gives, without a choice of representatives,
\[
\begin{split}
 (y,l,i)(z,m,j)&=(y,l\operatorname{hol}_i(z)m,j),\\
 \theta^{\mathrm{op}}(y,l,i)&=
 (\theta^*y,\operatorname{hol}_{\theta(n)}(y)l,i),\\
 (y,l,i)\theta^{\mathrm{op}}&=(y,l,\theta(i)).
\end{split}
\]
Together with multiplication in $O_{n+1}^{\mathrm{op}}$, these exhaust
all pairs.  The identity is $\mathrm{id}_{[n]}^{\mathrm{op}}$, not any
of the displayed triples, and associativity follows from the established
normal-form category (or from the interval identity above in the mixed
cases).  Counting the disjoint forms gives
$\binom{2n+1}{n}+(n+1)|L||Y_n|$.

For $k\le n$, choose a monotone inclusion $[k]\to[n]$ and a monotone
left inverse.  Opposites give the required retraction of $[k]$ from
$[n]$.  The remaining retraction is
\[
 \bullet_L\xrightarrow{s_n=(*_n,1)}[n]
 \xrightarrow{r_i=\rho i^{\mathrm{op}}}\bullet_L,
 \qquad r_i s_n=1.
\]
Consequently $e_i=s_nr_i$ is idempotent and
$l\mapsto s_nlr_i=(*_n,l,i)$ is the corner isomorphism, inverse to
$x\mapsto r_i x s_n$.  These equations prove surjectivity onto the
whole corner as well as injectivity.  The retract-generator lemma
identifies the classifying space of this endomorphism monoid with that
of the retraction category.
\end{proof}

\begin{proposition}[Check of Proposition 4.1 and both cases of Lemma 4.2]
\label{prop:B-cofiber}
Theorem~A of the supplied manuscript is correct, including its assertion
under $BL$ and its independence of the chosen vertex $i$.
\end{proposition}
\begin{proof}
\emph{Identification of the collage.}  Projection
$\pi_Y:\operatorname{El}(E_f)\to
(\Delta_{\le n}\downarrow Y)^{\mathrm{op}}$ forgets $l$.
Its section sends $y$ to $(y,1)$ and an opposite simplex operator to
$(\operatorname{hol}_{\theta(j)}(y),\theta)$.
The interval identity proves functoriality.  At $(y,l)$, the arrow
$(l,\mathrm{id})$ gives a natural transformation from the identity to
section-after-projection.  Thus this is a homotopy equivalence.
After reversing simplices, the projection to $B\mathbf L$ becomes
exactly $f$ after the last-vertex map: the arrow selected from a simplex
runs from vertex $\theta(j)$ to its last vertex and has label
$\operatorname{hol}_{\theta(j)}(y)$.  Equality of these labels on each
edge determines equality of the maps of nerves in every dimension.
The other collage vertex $B\Delta_{\le n}^{\mathrm{op}}$ contracts
to its initial object $[0]$.  The preceding collage lemma therefore
identifies $B\mathcal C(Y,f)$ with $\operatorname{Cof}(|f|)$ under
$BL$.

\emph{First comma-category case in Lemma 4.2.}  We check it for the
more general right $L$-simplicial set $E$ framed by $E_0\cong L$.
In $(J_E\downarrow\bullet_L)$, retain the objects $(\bullet_L,a)$
and $v_l=([0],l\rho)$.  Send $([k],l\rho i^{\mathrm{op}})$ to
$v_l$.  The arrow $i^{\mathrm{op}}$ is the component of a natural
transformation to this retraction.  For a mixed arrow $u$, naturality
is $i^{\mathrm{op}}u=i^*u$, and for a simplex operator it is
$i'^{\mathrm{op}}\theta^{\mathrm{op}}=(\theta(i'))^{\mathrm{op}}$.
The full subcategory on $(\bullet_L,a)$ has terminal object
$(\bullet_L,1)$.  Its comma category over $v_l$ has objects indexed by
$a\in L$ (the arrow labelled $a$ has source $(\bullet_L,la)$), and
arrows $a\to a'$ indexed by $m$ with $a'm=a$.  It has terminal object
$1$.  Quillen's Theorem~A twice gives contractibility of the original
comma category.  This argument never cancels $l$ in $L$.

\emph{Second comma-category case in Lemma 4.2.}  For the target $[j]$,
retain $(\bullet_L,u)$ for $u\in E_j$, the object
$t=([j],\mathrm{id})$, and $v_u=([0],u\rho)$.
An old arrow $\theta^{\mathrm{op}}:[k]\to[j]$ retracts to $t$;
an added arrow $u\rho i^{\mathrm{op}}$ retracts to $v_u$.
The components are respectively $\theta^{\mathrm{op}}$ and
$i^{\mathrm{op}}$.  There are no arrows between the two kinds of
non-$\bullet_L$ object: an old composite stays old and a composite
containing $\rho$ stays in its normal-form family.  On mixed arrows the
retraction is pullback along $\theta$ or the vertex $i$.  Simplicial
functoriality and $L$-equivariance prove all naturality squares.
The full subcategory $T$ on $t$ and the $(\bullet_L,u)$ has terminal
object $t$.  For each $v_u$, $(T\downarrow v_u)$ again has objects
$a\in L$ and arrows $a'm=a$, hence terminal object $1$.
Thus $BT\to BR\to B(J_E\downarrow[j])$ are equivalences and all
three spaces are contractible.  This checks the previously unverified
second case with its possibly nonfree $L$-action intact.

\emph{Assembling Theorem A.}  Quillen's Theorem~A now makes
$BJ_E:B\mathcal C(E)\to B\mathcal C^{\mathrm{ret}}(E)$ an
equivalence.  Combining it with the retract-generator comparison and
the collage identification gives $BM(Y,f)\simeq\operatorname{Cof}(|f|)$.
In the Karoubi envelope of the retraction category, the object
$([n],e_i)$ is isomorphic to $(\bullet_L,1)$ via $s_n$ and $r_i$.
The corresponding endomorphism isomorphism is precisely $\lambda_i$.
Therefore the canonical corner map followed by this equivalence is the
original inclusion of $BL$ into the collage and then into its cofiber.
The three comparisons defining the equivalence did not use $i$; the
same argument applies to every vertex simultaneously.
\end{proof}

\begin{proposition}[Checks of Section 5 and Theorems B--C]
\label{prop:B-finite-realization}
Theorem~B and Theorem~C of the supplied manuscript hold: every finite
simply connected CW complex is the classifying-space homotopy type of a
finite aperiodic monoid.
\end{proposition}
\begin{proof}
\emph{Lemma 5.1, simplicial approximation.}  Jardine's Corollary~4.8
\cite{Jardine2004B}
applies to a simplicial cofibration $A\hookrightarrow B$ with $B$
finite, a simplicial map $A\to X$, and a continuous extension
$|B|\to|X|$.  It supplies a simplicial approximation from
$\mathrm{sd}^m\mathrm{sd}_*\mathrm{sd}\mathrm{sd}_* B$, agreeing
with the prescribed map on the similarly subdivided $A$, and a
homotopy relative to that subspace.  The comparison to $B$ is a
composite of first- and last-vertex subdivision comparisons, whose
realizations are homotopy equivalences.  In the application take
$A=\varnothing$, $B=K$, and $X=N(L)$.  The inclusion is a cofibration,
$K$ is finite, and there is no condition on the target simplicial set.
Each subdivision preserves finiteness.  This proves exactly the
finite-source approximation used in Lemma~5.1, without a Kan hypothesis.

\emph{Lemma 5.2, stationary powers.}  An order-preserving self-map of a
finite chain has eventually stationary iterates on each vertex: once
$\theta(i)\ge i$ the successive values are nondecreasing, and the
reverse inequality gives a nonincreasing sequence.  A maximum of the
finitely many stopping times works for the whole map.  For an ideal
element $x=(y,l,i)$, put $c=\operatorname{hol}_i(y)l$; the
multiplication formula gives $x^r=(y,lc^{r-1},i)$.  Stationary powers
in $L$ therefore imply stationary powers in the whole monoid.
In any group image, $a^r=a^{r+1}$ forces $a=1$, so its universal group
is trivial; this checks simple connectivity as well.

\emph{Lemma 5.3, reducing vertices without changing the map.}  For
$a\in L$, choose $e=a^r$ with $a^r=a^{r+1}$.  Then $e^2=e$ and
$eae=e$.  The two horn attachments described in the manuscript are
well typed: the first triangle has edges $e,ye,e$ and the second
$y,e,ye$, where $y=a$ or $1$; the sole additional relation is
$eye=e$.  Each horn inclusion is a cofibration and a homotopy
equivalence, so each resulting two-triangle complex is contractible.
In the $y=1$ copy, collapse the edge labelled $1$; it has distinct
vertices before collapse, so this is a simplicial pushout along a
monomorphism and a homotopy equivalence.  Glue the $e$ edge of the
$y=a$ copy to the resulting $e$ loop.  The first leg is again a
monomorphism and both gluing spaces are contractible.  The result is a
finite reduced contractible simplicial set with a loop labelled $a$.
Replacing an edge between distinct vertices of $Y'$ by this complex
is a homotopy equivalence by Lemma~\ref{lem:B-elementary}, preserves
the map to $N(L)$ strictly, and lowers the vertex count by exactly one.
Connectivity ensures an edge with distinct endpoints whenever more
than one vertex remains.  Finite iteration proves the lemma, including
the claimed equality of realized maps.

\emph{Lemma 5.4, removal of one-cells.}  Geoghegan's Proposition~4.2.1
\cite{Geoghegan2008B}
states that for an $n$-connected CW pair $(X,A)$ with relative cells
only in dimensions at least $n$, the relative $n$-cells can be traded
for relative $(n+2)$-cells, preserving the number of $(n+1)$-cells
and cells in larger dimensions.  Its proof adds an $(n+1,n+2)$ cell
pair for each $n$-cell using a null-homotopy relative to $A$, then
pushes out a contraction of the subcomplex consisting of $A$, the old
$n$-cell and the new $(n+1)$-cell.  These are homotopy equivalences of
pairs.  Collapse a maximal tree of the connected finite $X^{(1)}$ to
obtain one vertex.  The resulting pair with that vertex is
$1$-connected, so take $n=1$, $A=*$; only finitely many cell pairs
are added.  The result is a finite model without one-cells.

\emph{Theorem B.}  Choose a finite simplicial model of the connected
finite CW source.  Apply the verified approximation, and then the
vertex reduction.  A homotopy inverse of the latter comparison gives
the map $h_{Y,f}$ with $|f_S|\simeq f h_{Y,f}$.  Theorem~A and
cofiber invariance identify the desired cofiber with the constructed
monoid.  Finiteness follows from its displayed underlying set, and
aperiodicity from the stationary-powers check.

\emph{Theorem C.}  In the one-vertex model with no one-cells, every
attaching sphere is $S^m$ with $m\ge1$, so it meets the connected-source
hypothesis of Theorem~B.  Begin with the trivial monoid and attach the
finitely many cells in nondecreasing dimension.  At each step compose
the attaching map with the already constructed equivalence to the
current classifying space.  The cofiber-invariance map is an
equivalence because it is the identity on the sphere and the current
equivalence on its target.  Theorem~B supplies the next finite aperiodic
monoid.  This finite induction proves the claim, with no assumption
that equivalences of spaces are induced by monoid homomorphisms.
\end{proof}

\subsection{Integral projective certificates}
The computational certificates specify the multiplication tables, every
projective summand, and every integral differential used here.  Their
source is \texttt{b\_projective\_resolutions.py}; outputs use the prefix
\texttt{b\_resolution\_}. The convention is left modules over $R=\mathbb ZM$.
The following elementary
description makes the certificates independently checkable.

\begin{lemma}\label{lem:B-projective-bases}
For an idempotent $e\in M$, $P(e)=Re$ has basis $Me$.  If $e,f$ are
idempotents satisfying $ef=fe=f$, then $P(e,f)=R(e-f)$ is projective,
has basis
\[
 \{v-vf:v\in Me\setminus Mf\},
\]
and left multiplication by $a$ sends this basis vector to
$av-(av)f$ if $av\notin Mf$, and to zero otherwise.  Tensoring with
the trivial right module sends $P(e)$ to $\mathbb Z$ and $P(e,f)$ to zero.
\end{lemma}
\begin{proof}
The first assertion follows from the idempotent endomorphism $r\mapsto re$
of the free module $R$.  The identities on $e,f$ imply
$(e-f)^2=e-f$, so the second module is also a direct summand of $R$.
Right multiplication by $f$ retracts the basis set $Me$ onto its subset
$Mf$, fixing the latter pointwise.  The displayed differences form a
basis of the kernel of this retraction: independence follows by looking
at their coefficients on $Me\setminus Mf$, and spanning follows by
subtracting the coefficient of each such basis element.  This kernel
equals $R(e-f)$.  The action formula follows from associativity.
On tensoring, the idempotent endomorphisms become multiplication by
$\eps(e)=1$ and $\eps(e-f)=0$, respectively.
\end{proof}

The construction of the certificates uses Smith transformations over
$\mathbb Z$, not rational nullspaces.  If $UDV$ is the Smith form of a
differential of rank $r$, the last columns of $V$ give an integral basis
$K$ of its kernel.  For the next displayed differential $D'=KW$,
the Hermite normal form of $W$ is the identity matrix.  Hence its image
is the whole integral kernel.  Equivariance is checked on every monoid
element and basis vector.  The augmentation is checked in exactly the
same manner.  These conditions prove exactness as $R$-modules, because
exactness of $R$-module maps is exactness of their underlying abelian
groups.  The trivial module is therefore resolved, rather than merely
approximated through a bar-complex degree bound.

\begin{proposition}\label{prop:B-resolutions}
The integral projective certificates give the following resolutions and
homology groups.  The tuple in the third column records the underlying
abelian ranks of $P_0,P_1,\ldots$; a repeated tail is marked explicitly.
\[
\begin{array}{c|r|l|l}
M&|M|&\operatorname{rank}_{\mathbb Z}P_*&\widetilde H_*(BM;\mathbb Z)\\ \hline
A_2&7&(2,4,3)&\mathbb Z\text{ in degree }2\\
A_3&16&(2,8,10,3)&\mathbb Z\text{ in degree }3\\
M(\operatorname{sk}_2NC_2,*)&22&(4,10,12,5)&\mathbb Z/2\text{ in degree }2\\
M(Y_z,f_z)&175&(22,244,290,67)&0\\
M(Y_z,*)&175&(22,277,312,56)&\mathbb Z\text{ in degrees }2,3\\
M(S^1_\Delta,f_1)&11&(4,10,9,4,4,\ldots)&\text{as in Proposition~\ref{prop:B-finite-resolution-obstruction}}\\
M(S^1_\Delta,f_g)&11&(4,6,5,4,4,\ldots)&\text{as in Proposition~\ref{prop:B-finite-resolution-obstruction}}\\
M(\operatorname{sk}_2NC_2,\iota)&34&(8,14,12,9,8,8,\ldots)&\mathbb Z/2\text{ in odd degrees }\ge3 .
\end{array}
\]
Every group not specified in the last column is zero.  In each infinite
resolution the tail is two-periodic as a complex of projective
$\mathbb ZM$-modules.
\end{proposition}
\begin{proof}
For the first three finite resolutions the exact-kernel computation ends with
zero.  After tensoring, their positive-degree chain groups are,
respectively, a single $\mathbb Z$ in degree $2$, a single
$\mathbb Z$ in degree $3$, and $\mathbb Z$ in degrees $2,3$ with
$d_3=2$.  This gives the first three rows.

For $Y_z$, the four generating triangles are
$(x_{00},x_{00})$, $(x_{00},x_{10})$,
$(x_{01},x_{00})$, $(x_{01},x_{10})$; their alternating cycle
has coefficients $(1,-1,-1,1)$ and pairs to one with $u$.
The three edges are killed by the first three triangular relations,
and the normalized boundary matrix has rank three with primitive
image, so $|Y_z|\simeq S^2$.  The inclusion into $N(L_0)$ has
degree one.  Both order-$175$ resolutions terminate after $P_3$.
For the inclusion all positive-degree projectives have zero
coinvariants, whereas for the constant map the only positive
coinvariants are one $\mathbb Z$ in each of degrees two and three,
with zero differential.  The certificates therefore give exactly the
two asserted homology lists.  Theorem~A identifies the homotopy types
as a point and $S^2\vee S^3$, respectively, and the resolution check
provides an independent all-degree homology verification.

For either circle example the repeated projective is $P(e)$ of abelian
rank four.  In the basis stored in the certificate, the two maps are
two identical blocks of
\[
 U=\begin{pmatrix}1&1\\1&1\end{pmatrix},\qquad
 V=\begin{pmatrix}1&-1\\-1&1\end{pmatrix},
\]
with a harmless common sign on $V$ in one example.  These maps are
$R$-linear, since this is checked against all the action matrices.
Moreover $UV=VU=0$, $\ker U=\mathbb Z(1,-1)=\im V$, and
$\ker V=\mathbb Z(1,1)=\im U$.  Thus repetition gives an exact
infinite resolution, with no extrapolation from finitely many homology
calculations.  The finite prefix meets this periodic complex exactly,
as verified by the kernel certificates.  Tensoring sends the maps to
$2$ and $0$.  For $f_1$, the tensor complex starts with ranks
$(1,1,2,1,1,\ldots)$ and maps $d_1=0$, $d_2=(2\;0)$,
$d_3=0$, $d_4=2$.  For $f_g$ its ranks are
$(1,0,1,1,1,\ldots)$ and maps $d_3=0$, $d_4=2$.
This proves both rows in all degrees independently of Theorem~A.

For the last example the repeated $P(e)$ has rank eight, and the same
two matrices occur in four blocks.  The periodic argument is unchanged.
Its tensor complex has one generator in degree zero and in every
degree at least three, none in degrees one and two, and differentials
$d_{2k}=2$, $d_{2k+1}=0$ in the tail.  It has the indicated homology.
Here the simplex $(g,g)$ has vertex holonomies $(1,g,1)$, so this is
an explicit two-dimensional case with nonconstant holonomy.  Theorem~A
also identifies its space as
$\mathbb{RP}^{\infty}/\mathbb{RP}^{2}$, consistently with the
projective calculation.
\end{proof}

\subsection{The two explicit projective-plane models}
Write $L_0=\{1,x_{00},x_{01},x_{10},x_{11}\}$, where
$x_{ij}x_{kl}=x_{il}$.  The proof of
Proposition~\ref{prop:B-band-wedge} below gives $BL_0\simeq S^2$.
Use normalized integral bar cochains, zero whenever an argument is $1$.
Define
\[
 u(x_{ij},x_{kl})=[j=0][k=0],\qquad
 w(a,b,c)=u(a,b)u(b,c).
\]
The bracket denotes the integer indicator of the stated condition.
Direct substitution gives $du=0$ and $dw=u\smile u$.  For clarity,
on four band elements the latter identity reduces to
\[
\begin{split}
 &u(b,c)u(c,d)-u(ab,c)u(c,d)+u(a,bc)u(bc,d)\\
 &\hspace{18mm}-u(a,b)u(b,cd)+u(a,b)u(b,c)
   =u(a,b)u(c,d),
\end{split}
\]
using $x_{ij}x_{kl}=x_{il}$; when an argument is $1$, adjacent terms
cancel.  The supplied program checks all $5^3$ and $5^4$ possibilities
as well.  The two-cycle
\[
 z=(x_{11},x_{11})+(x_{00},x_{00})
       -(x_{10},x_{10})-(x_{11},x_{00})
\]
has $\langle u,z\rangle=1$, so $u$ represents a generator of
$H^2(BL_0;\mathbb Z)$.

\begin{proposition}\label{prop:B-CP1795}
The twelve tetrahedra from the public supplement
\cite{BenBassatSupplementB}, specified below, define a finite reduced simplicial
set $Y\simeq S^3$ and a map $Y\to N(L_0)$ of Hopf invariant $1$.
Consequently Definition~3.3 gives an associative $1795$-element monoid
whose classifying space is $\mathbb{CP}^2$.
\end{proposition}
\begin{proof}
Use the central extension with multiplication
$(a,r)(b,s)=(ab,r+s+u(a,b))$.  Number the elements of $L_0$ by
$0=1$, $1=x_{00}$, $2=x_{01}$, $3=x_{10}$, $4=x_{11}$.
Generate $Y$ by the following ordered list of extension-labelled
tetrahedra and all their faces and degeneracies:
\[
\begin{array}{lll}
((2,-1),(0,1),(2,-1)),&((2,-1),(0,1),(3,-1)),&((2,-1),(2,0),(2,0)),\\
((2,0),(2,0),(2,-1)),&((2,0),(3,0),(3,-1)),&((1,-1),(0,1),(2,-1)),\\
((1,-1),(0,1),(3,-1)),&((1,-1),(2,0),(3,-1)),&((1,0),(2,-1),(3,0)),\\
((1,0),(1,-1),(1,-1)),&((1,0),(1,-1),(3,0)),&((1,0),(3,0),(3,-1)).
\end{array}
\]
The map to $N(L_0)$ forgets the integer labels.  The normalized chain
groups have ranks $(1,7,18,12)$, with $d_2$ of Smith rank seven and
all seven invariant factors $1$, and $d_3$ of Smith rank eleven with
all eleven invariant factors $1$.  The stored matrices give
$\ker d_2=\im d_3$ over $\mathbb Z$ and the generator
\[
 c=(-1,1,-1,-1,1,1,-1,1,-1,-1,1,-1)\in\ker d_3.
\]
These are exact integral assertions.  For an explicit fundamental-group
check, name the seven edge generators
$a=(1,-1)$, $b=(3,0)$, $c=(2,0)$, $d=(2,-1)$,
$e=(3,-1)$, $f=(1,0)$, $g=(0,1)$.
Among the eighteen triangular relations are
$a^2=a$, $ab=a$, $dc=d$, $db=a$, $ce=a$, $cb=f$, $gd=c$.
Successive cancellation gives $a=b=c=d=e=f=g=1$.
Thus $Y$ is simply
connected and has the integral homology of $S^3$.  Hurewicz supplies
a map $S^3\to |Y|$ representing $c$; it is a homology equivalence
between simply connected CW complexes and hence a homotopy equivalence.

Let $t(a,r)=r$.  The normalized coboundary convention gives
$dt=-f^*u$.  Consequently the cocycle
\[
 \Theta=-t\smile f^*u-f^*w,
 \qquad
 \Theta((a,r),(b,s),(c,t))=-(r+u(a,b))u(b,c)
\]
is the usual cochain definition of the Hopf invariant of $f$, since
$d(-t)=f^*u$ and $dw=u^2$.  Its values on the listed tetrahedra
are zero except for the tenth, where its value is $-1$.
Since the tenth coefficient of $c$ is $-1$,
$\langle\Theta,c\rangle=1$.  A generator of
$\pi_3(S^2)\cong\mathbb Z$ has Hopf invariant $\pm1$ and the
invariant multiplies by the degree on $S^3$; the attaching map is
therefore a Hopf map.  Its cofiber is $\mathbb{CP}^2$.

There are $88$ degree-three simplices including degeneracies.  Hence
the order formula is $\binom73+4\cdot5\cdot88=1795$.
The independent script reconstructs all faces, integral matrices,
cochains and products; it also checks all $1795^3$ associativity
triples.  Theorem~A then identifies the monoid's classifying space
with the computed cofiber.
\end{proof}

For the smaller model it is useful to state a slight extension of the
framed-collage argument explicitly.
\begin{lemma}\label{lem:B-framed-collage}
Let $K$ be a finite category, let $P:K\to\mathbf{Set}$ be a finite
covariant functor, and suppose $P(d_0)$ is a singleton and every object
of $K$ is a retract of $c_0$.  Then there is a finite monoid
\[
 \operatorname{End}_K(c_0)\ \amalg\
       \big(P(c_0)\times K(c_0,d_0)\big)
\]
whose classifying space is the homotopy cofiber of
$B\operatorname{El}(P)\to BK$.  The product of two elements in the
second summand is $(p,\alpha)(p',\alpha')=(p,\alpha')$;
the old monoid acts on the left through $P$ and on the right by
precomposition on $\alpha$.
\end{lemma}
\begin{proof}
Form the collage $G^-$ by adjoining an object $*$ with
$G^-(*,d)=P(d)$ and no arrows $d\to *$; its classifying space is
the stated cofiber by Lemma~\ref{lem:B-prelim}.  Let $s_0:*\to d_0$
be the unique such arrow.  Adjoin $\rho:d_0\to *$ and the relation
$\rho s_0=1_*$.  The same terminating normal-form calculation as
in Lemma~\ref{lem:B-monoid} gives
\[
 G(d,*)=K(d,d_0),\qquad
 G(d,d')=K(d,d')\amalg\big(P(d')\times K(d,d_0)\big).
\]
For the inclusion $J:G^-\to G$, the comma category
$(J\downarrow *)$ has terminal object $(d_0,\rho)$: for an
object $(d,\rho\alpha)$ the unique arrow to it is $\alpha$,
and for $(*,1)$ it is $s_0$.
For a target $d_1$, retract $(J\downarrow d_1)$ to the full
subcategory on $(*,p)$, $t=(d_1,1)$, and
$v_p=(d_0,p\rho)$, using an old arrow to $t$ and its last factor
$\alpha:d\to d_0$ for an added arrow.  The normal forms verify
naturality exactly as in Proposition~\ref{prop:B-cofiber}.
The subcategory on $t$ and $(*,p)$ has terminal object $t$;
its comma category over $v_p$ consists of the single arrow
$(*,p)\to v_p$ labelled $s_0$.  Quillen's Theorem~A makes this
comma category, and then $(J\downarrow d_1)$, contractible.
Thus $BJ$ is an equivalence.  The object $*$ retracts from $d_0$,
which retracts from $c_0$, so the retract-generator lemma identifies
$BG$ with $B\operatorname{End}_G(c_0)$.  Its multiplication is the
one asserted; the equality $P(d_0)=\{s_0\}$ is exactly what
removes the two middle factors in the product of new arrows.
\end{proof}

\begin{proposition}\label{prop:B-CP122}
The annulus construction of \cite{BenBassatSupplementB} gives an associative monoid of order $122$
with classifying space $\mathbb{CP}^2$.
\end{proposition}
\begin{proof}
Define a two-object category $D$ with objects $a,b$,
$D(a,a)=L_0$, $D(b,b)=\{1_b\}$,
$D(a,b)=\{r_0,r_1\}$ and $D(b,a)=\{i_0,i_1\}$,
where $r_j i_k=1_b$ and $i_k r_j=x_{kj}$.
The other composites follow from these identities.  In particular,
$b$ retracts from $a$, so $BD\simeq BL_0\simeq S^2$.

Let $\widetilde A$ have as simplices the nondecreasing integer tuples
$n_0\le\cdots\le n_q\le n_0+2$, and let
$\widetilde C$ use $n_q\le n_0+1$.  Put
$A=\widetilde A/(n\mapsto n+2)$ and
$C=\widetilde C/(n\mapsto n+1)$.  These simplicial sets have
dimensions two and one.  The maps
\[
 i_k(c)=2c+k,\qquad r_j(n)=\lfloor(n+j)/2\rfloor
\]
induce simplicial maps on the quotients and satisfy the displayed
relations of $D$.  They therefore give a diagram $E:D\to\mathbf{sSet}$
with $E(a)=A$, $E(b)=C$.  Both covers are contractible: $\widetilde C$
is the triangulated line, and $\widetilde A$ is a union of the
triangles $[n,n+1,n+2]$ with consecutive overlap an edge.  Each finite
consecutive union collapses successively onto its first triangle.
Every sphere and homotopy has compact image in a finite such union,
which proves contractibility of the locally finite union.

Take $K=D\times\Delta_{\le2}^{\mathrm{op}}$ and
$P(d,[q])=E(d)_q$.  Let $c_0=(a,[2])$, $d_0=(b,[0])$.
Every object retracts from $c_0$, and $P(d_0)$ is a singleton.
Now $|A_2|=12$, $|D(a,b)|=2$, and there are three vertices of
$[2]$.  Lemma~\ref{lem:B-framed-collage} gives order
\[
 |L_0|\,|O_3|+|A_2|\,|D(a,b)|\,3
       =5\cdot10+12\cdot2\cdot3=122.
\]

It remains to identify the attaching map, rather than just the two
spaces' homology.  The chosen lifts of arrows of $D$ commute with
deck translations, and their composites differ from the chosen lift
of the composite by a deck translation.  Its integer exponent is a
normalized two-cocycle $\omega$.  Direct substitution in the floor
functions gives
\[
 \omega(u,v)=1\quad\Longleftrightarrow\quad
 u\in\{r_1,x_{01},x_{11}\},\quad
 v\in\{i_1,x_{10},x_{11}\},
\]
for composable pairs, and zero otherwise.  The cocycle identity
expresses associativity of three lifted composites.  The explicit cycle
\[
 z' =-(x_{11},x_{11})-(x_{00},x_{00})
          +(x_{10},x_{10})+(x_{11},x_{00})
\]
has boundary zero and $\langle\omega,z'\rangle=-1$.
Since $H_2(BD)=\mathbb Z$, this proves that $[\omega]$ is a
generator, without assuming that a nonzero cocycle is primitive.

More explicitly, form the central extension $D_\omega$ with arrows
$(u,k)$ and product $(u,k)(v,l)=(uv,k+l+\omega(u,v))$.
Its action by lifted arrows on the contractible covers defines a
diagram over $D_\omega$.  The category of its truncated simplices
maps to $\operatorname{El}(P)$ by taking deck orbits.  This is an
equivalence of categories: every simplex has a lift, and a morphism
between two chosen lifts has one and only one required deck exponent.
Projection of this lifted category onto $D_\omega$ is a homotopy
equivalence by the homotopy-colimit theorem, since its fibers are the
truncated simplex categories of the contractible covers.  The
truncation lemma applies because their dimensions are at most two.
These comparisons commute with projection to $D$.

Finally $BD_\omega\to BD$ is the homotopy circle bundle classified
by $[\omega]$.  One can see this directly on nerves: over an
$n$-simplex of $ND$, the choices of integer labels on its $n$ arrows
form $(\mathbb Z)^n$, the $n$-simplices of $N\mathbb Z$.
Composition adds the cocycle, so this is the twisted principal
$N\mathbb Z$-bundle associated to $\omega$.  The resulting
classifying map to $B^2\mathbb Z\simeq K(\mathbb Z,2)$ has
cellular two-cocycle $\omega$.  Pulling back to the equivalent
$S^2$ gives Euler number $\pm1$, hence the Hopf bundle
$S^3\to S^2$, up to orientation.  Thus
$B\operatorname{El}(P)\to BK$ is a Hopf map, and its cofiber is
$\mathbb{CP}^2$.

The independent certificate verifies all composites in $D$, all
truncated simplicial naturality identities, the deck cocycle, its
pairing, and all $122^3$ associativity triples of the final table.
In addition, the separately generated integral projective resolution
\texttt{b\_resolution\_CP2\_annulus122.json} has length four and
abelian ranks $(12,108,182,99,14)$.  Its integral kernels and all
equivariance identities are verified as in
Proposition~\ref{prop:B-resolutions}.  After tensoring, only degrees
zero, two and four have a generator, giving a further all-degree
check of the additive homology used for this upper bound.
\end{proof}

\subsection{Exact minima and certified bounds}
\begin{lemma}\label{lem:B-band-Tor}
Suppose the minimal ideal of $M$ is a rectangular band $K=I\times J$.
Put $R=\mathbb ZM$, $A=R/\mathbb ZK$,
$V=\ker(\mathbb ZI\to\mathbb Z)$ as a left module, and
$W=\ker(\mathbb ZJ\to\mathbb Z)$ as a right module.  Then
\[
 H_n(BM;\mathbb Z)\cong\operatorname{Tor}^{A}_{n-2}(W,V)
 \qquad(n\ge2).
\]
The same assertion holds after replacing $\mathbb Z$ by any field.
In particular $H_2$ is the quotient of $W\otimes_{\mathbb Z}V$
by $wm\otimes v-w\otimes mv$ for $m\in M$.
\end{lemma}
\begin{proof}
Choose $e\in K$.  The multiplication map
$Re\otimes_{\mathbb Z}eR\to\mathbb ZK$ is an isomorphism:
on bases it is the bijection $I\times J\to K$.
The ideal $\mathbb ZK$ is therefore projective as both a left and
a right $R$-module.  The exact sequence
$0\to\mathbb ZK\to R\to A\to0$, tensored with either of
the modules annihilated by $K$, gives
$\operatorname{Tor}^{R}_q(A,V)=0$ for $q>0$ and
$A\otimes_R V=V$.  The change-of-rings spectral sequence, obtained
by filtering the two-sided double bar resolution, has
\[
 E^2_{p,q}=\operatorname{Tor}^{A}_p
       (W,\operatorname{Tor}^{R}_q(A,V))
 \Longrightarrow\operatorname{Tor}^{R}_{p+q}(W,V).
\]
Only its $q=0$ row is nonzero, giving
$\operatorname{Tor}^{R}_p(W,V)=\operatorname{Tor}^{A}_p(W,V)$.
Now use the two exact sequences
$0\to V\to Re\to\mathbb Z\to0$ and
$0\to W\to eR\to\mathbb Z\to0$.
Their middle terms are projective, and $eR\otimes_R V=eV=0$.
The two long exact Tor sequences give the asserted shift, including
$n=2$.  The proof is unchanged over a field.  The last formula is
the definition of the balanced tensor product.
\end{proof}

\begin{lemma}[Third homology under a band-preserving quotient]
\label{B:band-quotient-H3-surjection}
Let $q:M\twoheadrightarrow N$ be a surjective homomorphism of
finite monoids which restricts to a bijection between their
rectangular-band minimal ideals.  After identifying those ideals
by $q$, the induced maps satisfy
\[
 H_2(BM;k)\xrightarrow{\ \cong\ }H_2(BN;k),\qquad
 H_3(BM;k)\twoheadrightarrow H_3(BN;k)
\]
for every commutative coefficient ring $k$, including $\mathbb Z$.
In particular an actor image with nonzero third homology over a
field cannot be the image of a projective-plane monoid of this
kind, even when the action map is not faithful.
\end{lemma}
\begin{proof}
Write the common band as $I\times J$, and put
$V=\widetilde{kI}$ and $W=\widetilde{kJ}$.
Both action modules are unchanged by $q$: multiplication against
a band element is preserved, and $q$ is injective on the band.
The two augmentation sequences used in the band--Tor argument
are therefore natural under $q$, with identity maps on $V,W$.
Their connecting homomorphisms identify the actual homology map
with
\[
 H_{n+2}(Bq;k)=
 \operatorname{Tor}^{kM}_n(W,V)\longrightarrow
 \operatorname{Tor}^{kN}_n(W,V),\qquad n\geq0.
\]
This is also seen directly from the natural full two-sided bar
comparison in the proof of Lemma~\ref{lem:B-band-Tor}.

In degree zero, the tensor-product relations are identical:
the actions of the source elements on $W,V$ are exactly the
actions of their images, and $q$ is surjective.  Thus the Tor
map in degree zero is an isomorphism.

For degree one, use the ordinary two-sided bar over $k$.
Its first two terms are
\[
 C_1(M)=W\otimes_k kM\otimes_k V,
 \qquad C_0(M)=W\otimes_k V,
\]
with differential
\[
 d(w\otimes m\otimes v)=wm\otimes v-w\otimes mv.
\]
The induced map on $C_0$ is the identity and the map on $C_1$
is surjective.  If $z$ is a one-cycle for $N$, choose any lift
$\widetilde z\in C_1(M)$.  Commutation with the differential
gives $d\widetilde z=dz=0$ in the identical degree-zero module.
Every target cycle therefore lifts to a source cycle, and
boundaries map to boundaries.  This proves the Tor-one
surjection.  All modules in this bar construction are free over
$k$ in the requisite outer factors, so it computes the stated
ordinary Tor groups over any $k$.  Naturality of the two
connecting maps gives the asserted statement for $H_3$.
\end{proof}


\begin{proposition}\label{prop:B-K22-obstruction}
If the minimal ideal of $M$ is a $2\times2$ rectangular band and
$H_2(BM;\mathbb Z)$ has nonzero torsion, then
$H_2(BM;\mathbb Z)=\mathbb Z/2$ and
$H_n(BM;\mathbb F_2)\ne0$ for arbitrarily large $n$.
Thus no such monoid is a finite-homotopy-type model of either
torsion Moore space in Part~B4.
\end{proposition}
\begin{proof}
Both $V$ and $W$ in Lemma~\ref{lem:B-band-Tor} have rank one.
Their action characters $\alpha,\beta$ take values in
$\{0,1,-1\}$: a constant transformation acts by zero and the two
permutations act by their signs.  Therefore
\[
 H_2(BM;\mathbb Z)=
 \mathbb Z/(\alpha(m)-\beta(m):m\in M).
\]
Nonzero torsion is possible exactly when all differences are even
and some difference is $\pm2$.  This gives $\mathbb Z/2$;
it also says that $\alpha$ and $\beta$ have the same nonzero support.
Let $X$ be that support.  It is a submonoid containing the identity,
and its complement is a two-sided ideal consisting of elements acting
constantly on both two-element sets.  The ring $A$ consequently
retracts onto $\mathbb ZX$, by sending the other basis elements to
zero.  Both $V$ and $W$ factor through this retraction.

Apply the retraction and its section to the middle factors of the
two-sided bar complexes for $\operatorname{Tor}$.  This gives an
actual chain retraction, so
$\operatorname{Tor}^{\mathbb F_2X}_*(W\otimes\mathbb F_2,
V\otimes\mathbb F_2)$ is a direct summand of
$\operatorname{Tor}^{A\otimes\mathbb F_2}_*(W\otimes\mathbb F_2,
V\otimes\mathbb F_2)$.  The former is $H_*(BX;\mathbb F_2)$,
because both signs become the trivial action modulo two.
The character $\alpha\beta:X\to\{\pm1\}$ is onto, so
$U(X)$ has even order.  Theorem~\ref{A:bounded} makes
$H_*(BX;\mathbb F_2)$ unbounded.  The direct summand and the
two-degree shift in Lemma~\ref{lem:B-band-Tor} give the result.
\end{proof}

\begin{proposition}\label{prop:B-two-row-products}
Suppose the minimal ideal of a finite monoid $M$ is a rectangular
band with two rows or with two columns.  If $BM$ has finite homotopy
type, then every product of two integral degree-two cohomology classes
is zero.  In particular such a monoid cannot have classifying space
$\CP^2$ or $S^2\times S^2$.
\end{proposition}
\begin{proof}
Taking the opposite monoid exchanges rows and columns and preserves
the classifying-space homotopy type, so suppose that the rows are
$\{0,1\}$ and the columns form $J$ with distinguished column $0$.
Write the left row action of $m$ as
$i\mapsto a_m+\chi_m i$, where
$a_m\in\{0,1\}$ and $\chi_m\in\{0,1,-1\}$.
The right column transformation is denoted $\tau_m$.
The augmentation module $V$ is the rank-one character $\chi$.
By Lemma~\ref{lem:B-band-Tor}, an integral degree-two cohomology
class is specified by an integral functional $\lambda:W\to\mathbb Z$
satisfying
\[
 \lambda(wm)=\chi_m\lambda(w).
\]
Indeed the minimal group is trivial, so $H_1(BM;\mathbb Z)=0$;
the universal coefficient theorem identifies $H^2$ with the dual of
$W\otimes_A V$.  Set
\[
 \ell_j=\lambda(e_j-e_0),\qquad
 b_m=\ell_{\tau_m(0)}.
\]
The balanced-functional condition is precisely
\[
 \ell_{\tau_m(j)}=b_m+\chi_m\ell_j.
\]
Composition of the row and column actions gives
\[
 a_{mn}=a_m+\chi_m a_n,\qquad
 b_{mn}=b_n+\chi_n b_m.
\]
Consequently $u_\lambda(m,n)=b_m a_n$ is a normalized integral
two-cocycle.  It represents the indicated functional, up to a fixed
sign: on the rectangular-band two-cycle
\[
 -(x_{1j},x_{1j})-(x_{00},x_{00})
 +(x_{10},x_{10})+(x_{1j},x_{00})
\]
its value is $-\ell_j$.  These cycles give the natural generators
of $W\otimes V$ before imposing the balanced relations, as is also
seen from the two augmentation sequences in the preceding lemma.
Thus these cocycles represent all integral degree-two classes.

The finite set $L=\{\ell_j:j\in J\}\subset\mathbb Z$ has the
following elementary consequence.  If $\chi_m=1$, translation by
$b_m$ maps $L$ into itself, so $b_m=0$.  If $\chi_m=-1$, reflection
$x\mapsto b_m-x$ maps $L$ into itself and hence onto itself, so
\[
 b_m=c:=\min L+\max L.
\]
If no negative character occurs, $\chi_m b_m=0$ for every $m$.
Suppose a negative character does occur.  Its support
$X=\{m:\chi_m\ne0\}$ is a finite submonoid containing $1$, and
$\chi:X\to C_2$ is onto.  On $X$ one has
$a_m=(1-\chi_m)/2$ and $b_m=ca_m$.
Thus $u_\lambda|_X$ is $c$ times the integral cocycle representing
the first Chern class of the sign line bundle pulled back from $BC_2$.

We claim that $c$ is even.  Choose an idempotent in the minimal ideal
of $X$ and its finite maximal subgroup $H$.  Lemma~\ref{A:minimal-group}
shows that $H\to C_2$ is onto; its map to $BM$ is understood through
the Karoubi envelope, as in that lemma.  If $c$ were odd, the restriction
of a complex line bundle with first Chern class $[u_\lambda]$ would
be the nontrivial sign line over $BH$.  Such line bundles are classified
by $H^2$, so this assertion follows from the displayed cocycle, not
only from a rational characteristic class.  The reduced class of any
line bundle on a finite-dimensional CW complex is nilpotent in complex
$K$-theory, by the multiplicative skeletal filtration used in
Lemma~\ref{A:filtration}.  Its pullback would therefore be nilpotent.
On the other hand choose a $2$-power-order element of $H$ with negative
sign and restrict to its cyclic subgroup $C_{2^r}$.  Atiyah's completion
theorem, in the form checked in Proposition~\ref{A:atiyah}, identifies
its completed representation ring with $K^0(BC_{2^r})$.
Evaluation of the defining character at a primitive $2^r$-th root
of unity extends to the $2$-adic cyclotomic integers: the augmentation
ideal maps into their maximal ideal.  The reduced sign representation
maps to $-2$, whose every power is nonzero.  This contradicts
nilpotence and proves that $c$ is even.

Replace $b_m$ by
$b'_m=b_m+(\chi_m-1)c/2$ when negative characters occur, and leave
$b_m$ unchanged otherwise.  It still obeys the same affine composition
law, and now $\chi_m b'_m=0$ for every $m$.
This changes $u_\lambda$ by the coboundary
$-(c/2)\delta a$, since
$\delta a(m,n)=(1-\chi_m)a_n$.
Make this choice separately for any two functionals $\lambda,\mu$.
For their resulting cocycles $u_\lambda,u_\mu$, define
\[
 w(m,n,p)=b'_\lambda(m)\,a_n\,b'_\mu(n)\,a_p.
\]
Expanding the five terms of its coboundary gives
\[
 \begin{split}
 (\delta w)(m,n,p,q)
 &=u_\lambda(m,n)u_\mu(p,q)\\
 &\quad+b'_\lambda(m)a_q\,
       \chi_n b'_\mu(n)\,\chi_p a_p
  =u_\lambda(m,n)u_\mu(p,q).
 \end{split}
\]
Thus every mixed cup product is a coboundary, as asserted.
\end{proof}

\begin{lemma}\label{lem:B-collapse-minimal-group}
Let the minimal ideal of a finite monoid $M$ be the Rees matrix
semigroup $K=I\times H\times J$.  Collapse each of its $H$-classes
to one element, leaving $M\setminus K$ unchanged, and call the resulting
monoid $N$.  The quotient homomorphism $q:M\to N$ induces a rational
homology equivalence, hence an isomorphism of rational cohomology rings.
The minimal ideal of $N$ is the rectangular band $I\times J$.
\end{lemma}
\begin{proof}
Left multiplication by an element of $M$ acts on $K$ in the form
$(i,h,j)\mapsto(\sigma(i),a_i h,j)$, and right multiplication
in the form $(i,h,j)\mapsto(i,hb_j,\tau(j))$; these forms follow
by commuting with all translations on the opposite side.  Thus the
specified equivalence relation is a two-sided congruence.

Put $R=\mathbb QM$ and $R'=\mathbb QN$.  The kernel $J_0$ of
$R\to R'$ consists of the differences within each $H$-class of $K$.
Choose an idempotent $e\in K$ and identify its maximal group with $H$.
Multiplication gives an isomorphism
\[
 \mathbb QK\cong Re\otimes_{\mathbb QH}eR.
\]
Since $Re$ is free as a right $\mathbb QH$-module, this is a
projective right $R$-module.  Averaging the group coordinate in each
$H$-class gives an $R$-bimodule idempotent on $\mathbb QK$:
the displayed left and right actions preserve the averages.
Its kernel is $J_0$, so $J_0$ is projective as a right $R$-module.
Furthermore averaging acts as the identity on
\[
 \mathbb QK\otimes_R\mathbb Q
 \cong Re\otimes_{\mathbb QH}\mathbb Q,
\]
because the group coordinate is already trivial there.  It follows
that $J_0\otimes_R\mathbb Q=0$.
The length-one right-projective resolution
$0\to J_0\to R\to R'\to0$ now gives
$\operatorname{Tor}^R_i(R',\mathbb Q)=0$ for $i>0$ and
$R'\otimes_R\mathbb Q=\mathbb Q$.
The natural change-of-rings spectral sequence therefore identifies
$\operatorname{Tor}^R_*(\mathbb Q,\mathbb Q)$ with
$\operatorname{Tor}^{R'}_*(\mathbb Q,\mathbb Q)$ through the map
induced by $q$.  This proves the homology assertion.  Over a field,
the universal coefficient theorem makes $q^*$ a cohomology
isomorphism, and it preserves products because it is induced by a map
of classifying spaces.
\end{proof}

\begin{corollary}\label{cor:B-two-row-general}
If the minimal ideal of a finite monoid has two rows or two columns,
every product
$H^2(BM;\mathbb Q)\otimes H^2(BM;\mathbb Q)\to H^4(BM;\mathbb Q)$
is zero.  Neither finite homotopy type nor a trivial minimal maximal
group is required.
\end{corollary}
\begin{proof}
Apply Lemma~\ref{lem:B-collapse-minimal-group}.  Integral balanced
functionals for the rectangular-band quotient $N$ span its rational
degree-two cohomology, and hence their pulled-back affine cocycles
span $H^2(BM;\mathbb Q)$.
Repeat the affine calculation of
Proposition~\ref{prop:B-two-row-products} over $\mathbb Q$.
For each balanced functional the finite set of rational column
values has no nonzero translation symmetry, and every reflection
has the same center $c/2$.  Consequently the rational gauge
$b'_m=b_m+(\chi_m-1)c/2$ always satisfies $\chi_m b'_m=0$.
If there is no reflection, take $c=0$.
Unlike the integral assertion, this requires no parity restriction
on $c$ and no finite-homotopy hypothesis.  The same displayed
three-cochain therefore kills each mixed product over $\mathbb Q$.
The rational spanning assertion and naturality under the
minimal-group quotient prove the result.
\end{proof}

\begin{proposition}[Two rows permit only elementary two-torsion]
\label{B:two-row-torsion}
Let $M$ be a finite monoid whose minimal ideal is a rectangular band
with two rows or two columns. If $BM$ has finite CW homotopy type, then
the torsion subgroup of $H_2(BM;\mathbb Z)$ is a direct sum of copies
of $\mathbb Z/2$. In particular it has no odd-primary torsion and no
element of order four. If every row action fixes the two rows or is
constant, then $H_2(BM;\mathbb Z)$ is torsion-free.
\end{proposition}
\begin{proof}
Pass to the opposite monoid if necessary, and write the minimal ideal
as $\{0,1\}\times J$. Let $\chi_m\in\{0,1,-1\}$ be the action of $m$
on the row augmentation module, and let $\tau_m:J\to J$ be its right
column action. Fix $j_0\in J$. The natural tensor-product description
in Lemma~\ref{lem:B-band-Tor} presents the abelian group
$A=H_2(BM;\mathbb Z)$ by generators $\ell_j$, with $\ell_{j_0}=0$,
and relations
\begin{equation}\label{B:torsion-affine}
 \ell_{\tau_m(j)}=\chi_m\ell_j+b_m,
 \qquad b_m=\ell_{\tau_m(j_0)}.
\end{equation}
The generator $\ell_j$ is the class of
$(e_j-e_{j_0})\otimes(e_1-e_0)$. Composition gives
\[
 \chi_{mn}=\chi_m\chi_n,\qquad
 b_{mn}=b_n+\chi_n b_m.
\]

We first prove that $b_m=0$ whenever $\chi_m=1$. Let
$X=\{m\in M:\chi_m=1\}$. This is a finite submonoid, and
$b:X\to A$ is an additive monoid homomorphism. Its finite image is a
group. The local maximal group of the minimal ideal of $X$ maps onto
this image, by the same finite-semigroup argument as in
Lemma~\ref{A:minimal-group}. Thus it suffices to prove $b_h=0$ for elements
of any local finite subgroup $H\subset X$.

Let $e$ be the identity of $H$. Its row action is the identity, and its
column action is an idempotent projection. Choose a column orbit
$J_0\subseteq\operatorname{im}\tau_e$ for $H$, and let $K\leq H$ be
the stabilizer of a column $j\in J_0$. Inside the corner $eMe$ there is
the finite monoid
\[
 S=H\ \amalg\ (\{0,1\}\times J_0).
\]
The group fixes the two rows and acts transitively on the columns.
The cofiber description in Theorem~\ref{E:group-band-cofiber}
therefore gives
\begin{equation}\label{B:torsion-group-orbit}
 BS\simeq\Sigma\operatorname{hocofib}(BK\longrightarrow BH).
\end{equation}
For completeness, the Borel construction of the join of the trivial
two-point row set and the orbit $H/K$ is
$BH\cup^h_{BK}BH$. Its map to $BH$ is the fold map. The homotopy
cofiber of that fold is the suspension of the homotopy cofiber of
$BK\to BH$, which proves the displayed formula.

Every based map from the right side of
\eqref{B:torsion-group-orbit} to a finite connected CW complex is
null-homotopic. Indeed Miller's full mapping-space theorem says that
the based mapping spaces from $BH$ and $BK$ to such a target are
weakly contractible. Applying based mapping spaces to the cofiber
sequence and then looping gives the assertion. Apply it to a finite
CW model of $BM$. The corner inclusion $BS\to BM$ is consequently
null-homotopic.

The degree-two band-Tor formula is natural for the inclusion of $S$
into the corner and then into $M$. It identifies
\[
 H_2(BS;\mathbb Z)
   \cong \bigl(\ker(\mathbb ZJ_0\to\mathbb Z)\bigr)_H.
\]
To interpret this naturality using only unital monoid maps, adjoin a
fresh identity to $S$. The new monoid $S^+$ maps unitally to $M$,
and its old-identity corner is $S$ with the usual one-sided corner
equivalence $BS\simeq BS^+$. Its minimal band and balanced tensor
presentation are unchanged. Thus the ordinary unital naturality of
Lemma~\ref{lem:B-band-Tor} gives exactly the stated corner map.
The coinvariant class of $e_{\tau_h(j)}-e_j$ maps to
$\ell_{\tau_h(j)}-\ell_j=b_h$ in $A$, by
\eqref{B:torsion-affine}. Since the inclusion is null-homotopic,
$b_h=0$. Surjectivity from the local group proves $b_m=0$ for every
$m\in X$.

If there are no elements with $\chi_m=-1$, the remaining relations
in \eqref{B:torsion-affine} only identify pairs of column values:
positive-character maps identify $\ell_{\tau_m(j)}$ with $\ell_j$,
and zero-character maps identify all values in their image. Together
with $\ell_{j_0}=0$, this presents a free abelian group.

Otherwise all negative-character constants have the same value $c$.
Indeed for $\chi_m=\chi_n=-1$ the product has positive character, so
$0=b_{mn}=b_n-b_m$. The presentation is now a signed graph on the
column variables: positive edges impose $\ell_v=\ell_w$, negative
edges impose $\ell_v+\ell_w=c$, and $\ell_{j_0}=0$ anchors one vertex.
An extra variable $c$ with its defining relation can be used without
changing the group. Each negative action supplies the negative edge
from $j_0$ to $\tau_m(j_0)$, so this defining relation is already one
of the signed-graph relations.

In a connected component with no cycle containing an odd number of
negative edges, all its variables are either $x$ or $c-x$, for one
free variable $x$. If there is an odd signed cycle, the additional
relation is exactly $2x=c$. The anchored component has values $0,c$;
an odd signed cycle there forces $c=0$. These conclusions follow by
choosing a spanning tree in each component and then examining each
remaining edge, so they describe the complete presentation.
If $c=0$, the only torsion relations are $2x=0$. If $c\ne0$ and there
are odd components, eliminate $c$ using one relation $2x_1=c$; every
other such relation becomes $2(x_i-x_1)=0$. With no odd components
the presentation is free. Thus the only possible torsion summands
are $\mathbb Z/2$, as claimed.
\end{proof}

\begin{lemma}[Finite targets kill the torsion from a group above a band]
\label{B:group-band-torsion-killed}
Let $H$ be a finite group acting on finite nonempty row and column
sets $I,J$, and let $S=H\amalg(I\times J)$ be the corresponding
group-band monoid. Every map $BS\to X$ to a connected finite CW
complex kills the torsion subgroup of $H_2(BS;\Z)$.
\end{lemma}
\begin{proof}
Theorem~\ref{E:group-band-cofiber} identifies $BS$ with the homotopy
cofiber of $(I*J)_{hH}\to BH$. The Borel construction is a finite
graph of classifying spaces of finite stabilizer groups: its vertex
groups are stabilizers of row and column orbits, and its edge groups
are stabilizers of pairs. Miller's full mapping-space theorem
\cite[Theorem~A]{miller1984} makes evaluation
$\operatorname{Map}(BL,X)\to X$ a weak equivalence for each of
these finite groups $L$. Naturality and the finite homotopy-colimit
description show that collapse to the orbit graph induces a weak
equivalence on mapping spaces into $X$. The same holds for
$BH\to *$. Applying the cofiber mapping-space sequence shows that
\[
 BS\longrightarrow\Sigma\bigl((I*J)/H\bigr)
\]
induces a weak equivalence on mapping spaces into $X$. The orbit
graph is finite and connected, so its suspension is a finite wedge
of two-spheres. Every map from $BS$ to $X$ therefore factors up to
homotopy through a space with torsion-free second homology, which
proves the assertion.
\end{proof}

\begin{theorem}[A three-by-three minimal band has no three-primary torsion]
\label{B:three-by-three-no-three}
Let $M$ be a finite monoid whose minimal ideal is a $3\times3$
rectangular band. If $BM$ has finite CW homotopy type, then
$H_2(BM;\Z)$ has no nonzero three-primary torsion.
\end{theorem}
\begin{proof}
Let $V,W$ be the row and column augmentation modules, each free
of rank two. Lemma~\ref{lem:B-band-Tor} identifies
$A=H_2(BM;\Z)$ with the balanced quotient of $W\otimes V$.
Suppose that $A$ has three-primary torsion. Choose
$\lambda:A\to\F_3$ nonzero on that torsion. In augmentation
bases, write its nonzero bilinear matrix as $L\in M_2(\F_3)$.
For the row and column transformations $\sigma_m,\tau_m$, let
$A_m,B_m$ be their integral augmentation matrices. Balance says
\begin{equation}\label{B:three-balance}
 B_m^{\mathsf T}L=LA_m\pmod3.
\end{equation}

We first record the complete small linear classification. Up to
relabeling the two three-point sets, transposition, and multiplying
$L$ by $-1$, its possibilities are the following five types:
\[
\begin{array}{c|c|c}
\text{type}&L&\text{number of projective matrices}\\ \hline
\text{two levels, two levels}&
 \left(\begin{smallmatrix}0&0\\0&1\end{smallmatrix}\right)&9\\
\text{two levels, three levels}&
 \left(\begin{smallmatrix}0&0\\1&2\end{smallmatrix}\right)&6\\
\text{three levels, three levels}&
 \left(\begin{smallmatrix}1&2\\2&1\end{smallmatrix}\right)&1\\
\text{rank two, first orbit}&
 \left(\begin{smallmatrix}0&1\\1&0\end{smallmatrix}\right)&18\\
\text{rank two, second orbit}&
 \left(\begin{smallmatrix}0&1\\1&1\end{smallmatrix}\right)&6.
\end{array}
\]
Here is a proof of completeness. The projective dual of the
three-point augmentation module has four points. One is represented
by values $(0,1,2)$ and is fixed by all permutations; the other three
are the two-level functions and are permuted transitively. Rank-one
matrices therefore give the first three orbits, of sizes $9,6,1$.
For rank two, the projective permutation action is the order-six
stabilizer of the invariant point in
$\operatorname{PGL}_2(\F_3)$, which has order $24$ and acts on four
points. Its left and right actions on projective isomorphisms have
two double cosets, according as the distinguished point maps to the
distinguished point or to one of the other three. Their sizes are
$6$ and $18$, represented by the two displayed invertible matrices.
The orbit sizes sum to $(3^4-1)/2=40$.

In the first, fourth and fifth types choose the displayed integral
matrix, whose entries are zero or one. Each side of a single entry
of \eqref{B:three-balance} is a difference of two entries of the
extended zero-one grid, with row and column zero adjoined. Their
difference lies between $-2$ and $2$. Divisibility by three therefore
forces integral equality. Hence $L$ defines an integral balanced
functional $A\to\Z$ lifting $\lambda$. Integral functionals kill
torsion, contrary to the choice of $\lambda$. Relabelings lift to
integral permutation maps, so the conclusion holds throughout each
of these three orbits.

Consider the mixed type. Write $L=b\otimes a$, with row values
$a=(0,1,2)$ and column values $b=(0,0,1)$. The rank-one balance
equation means that every element has a common slope
$\chi_m\in\{0,1,-1\}\subset\F_3$:
\[
 a_{\sigma_m(i)}=\chi_ma_i+c_m,\qquad
 b_{\tau_m(j)}=\chi_mb_j+d_m.
\]
Suppose an element with $\chi_m=1$ acts nontrivially on rows.
Its row action is a three-cycle. Its column action preserves both
$b$-fibers, because a nonzero translation cannot preserve a two-point
subset of $\F_3$. In the eventual cyclic local group of powers of
$m$, take its three-Sylow subgroup $H$, with identity $e$.
The row action remains transitive on all three rows. The column
permutation image has order at most two, so $H$ fixes every column
in $J'=\operatorname{im}\tau_e$. This image retains both $b$-values.
The corner contains
$S=H\amalg(\{0,1,2\}\times J')$. Its balanced presentation is
\[
 H_2(BS;\Z)
 \cong\ker(\Z J'\to\Z)\otimes
       \bigl(\ker(\Z^3\to\Z)\bigr)_{C_3}
 \cong(\Z/3)^{|J'|-1}.
\]
The functional $\lambda$ is nonzero on this group: take the
difference of columns with distinct $b$-values and the difference
of rows with distinct $a$-values. This contradicts
Lemma~\ref{B:group-band-torsion-killed} for the map to a finite CW
model of $BM$. For the nonunital corner inclusion, adjoin a fresh
identity to $S$; its old-identity corner has the same classifying
space and balanced presentation, exactly as in
Proposition~\ref{B:two-row-torsion}.

Consequently every slope-one row action is the identity. Any two
slope-minus-one row reflections have the same fixed point, since
two distinct reflections compose to a nontrivial slope-one
translation. Assign integral row values $-1,0,1$, placing zero at
that common fixed point and choosing the signs to recover $a$
modulo three up to an additive constant. If no reflection occurs,
choose these three lifts in any order with the same residues.
Keep the integral column values $(0,0,1)$. A slope-zero row action
is constant and a slope-zero column action has constant $b$-value;
slope-one actions preserve both functions; and a slope-minus-one
action negates the row function and sends $b$ to $1-b$.
Thus their integral linear parts have the same slope, and their
tensor product is an integral balanced lift of $\lambda$.
This again contradicts torsion detection.

It remains to treat the three-level, three-level type. Both value
functions are bijections to $\F_3$. A slope-zero pair consists of
two constant transformations. Every other pair consists of two
affine permutations with the same slope. The permutation pairs in
the image of $M$ form a finite group $P$. The preimage of $P$ in
$M$ is a finite submonoid mapping onto $P$, so a local maximal
group $H$ maps onto $P$ by Lemma~\ref{A:minimal-group}. Its identity
acts identically on the whole minimal band $K$. Hence the corner
contains $S=H\amalg K$.

The map $H_2(BS;\Z)\to H_2(BM;\Z)$ is an isomorphism. Indeed
the balanced presentations have exactly the same relations from
$P$, and every remaining action pair is constant on both sides,
so both augmentation actions are zero and it contributes no
relation. Naturality of Lemma~\ref{lem:B-band-Tor}, with the same
fresh-identity treatment of the corner, identifies this isomorphism
with the actual map of classifying spaces. But
Lemma~\ref{B:group-band-torsion-killed} says that this map kills
torsion. Thus $A$ is torsion-free in this last case, the final
contradiction.
\end{proof}

\begin{lemma}[Rank-one balanced forms]
\label{B:rank-one-square}
Let the minimal ideal of a finite monoid be a rectangular band.
A rational degree-two class whose balanced matrix has rank at most one
has square zero.  No finite-homotopy hypothesis is needed.
\end{lemma}
\begin{proof}
There is nothing to prove for the zero matrix.  Write a rank-one
matrix as $L=b\otimes a$.  The balancing equation
$B_m^tL=LA_m$ implies that the scalar row and column augmentation
functionals have the same character $\chi_m\in\mathbb Q$.
Let their finite sets of values be $\{a_i\}$ and $\{b_j\}$,
with the distinguished values zero.  The corresponding affine actions
have the form
\[
 a_{\sigma_m(i)}=a_m+\chi_m a_i,\qquad
 b_{\tau_m(j)}=b_m+\chi_m b_j.
\]
If $\chi_m\ne0$, each of these maps is an injection of a finite
nonconstant set into itself and hence a permutation.  Its rational
slope is therefore $1$ or $-1$.  Positive-slope translations vanish,
and every negative-slope map has the same midpoint on each value set.
Choose rational midpoints $A,B$, taking zero if no negative map occurs,
and replace
\[
 a'_m=a_m+A(\chi_m-1),\qquad
 b'_m=b_m+B(\chi_m-1).
\]
Then $\chi_m a'_m=\chi_m b'_m=0$.  The original representative is
$u(m,n)=b_m a_n$; these affine representatives are identified with
balanced functionals exactly as in
Proposition~\ref{prop:B-two-row-products}.  The new representative
$u'(m,n)=b'_m a'_n$ has the same cohomology class, because
\[
 u'-u=-\delta\bigl(A b+B a+AB(\chi-1)\bigr).
\]
For $w(m,n,p)=b'_m a'_n b'_n a'_p$, direct expansion gives
\[
 \delta w=u'\smile u'
   +b'_m a'_q\,\chi_n b'_n\,\chi_p a'_p
 =u'\smile u'.
\]
This proves the assertion.
\end{proof}

\begin{theorem}[The finite-type three-by-three obstruction]
\label{B:three-by-three-cups}
Let $M$ be a finite monoid whose minimal ideal has three rows and
three columns, with any finite maximal group.  If $BM$ has finite
CW homotopy type, then every product
\[
 H^2(BM;\mathbb Q)\otimes H^2(BM;\mathbb Q)
       \longrightarrow H^4(BM;\mathbb Q)
\]
is zero.  In particular, $BM$ is neither $\mathbb CP^2$ nor
$S^2\times S^2$.
\end{theorem}
\begin{proof}
First suppose the minimal ideal is the band $I\times J$ with
$I=J=\{0,1,2\}$.  A transformation $f$ of a three-point set acts
on its augmentation lattice by the matrix
\[
 A_f(i,j)=\mathbf1_{f(j)=i}-\mathbf1_{f(0)=i},
 \qquad i,j\in\{1,2\}.
\]
An actor is a pair $(\sigma,\tau)$, with multiplication
$(\sigma,\tau)(\alpha,\beta)=(\sigma\alpha,\beta\tau)$.
Its balanced forms satisfy
\begin{equation}\label{B:3x3-balance}
 B_\tau^tL=LA_\sigma.
\end{equation}
Write $R_M\subset\mathbb Q^4$ for the space spanned by all
relations in this equation.  The maximal action hull $H(R_M)$ consists
of all pairs whose relation space is contained in $R_M$.
It is a monoid: the balancing equation is preserved by composition.
The actual action homomorphism $M\to H(R_M)$ fixes the minimal band.
The hull has exactly the relation space $R_M$, since it contains
all actual actors and introduces no relation outside $R_M$.
The natural balanced presentation of Lemma~\ref{lem:B-band-Tor}
therefore says that the induced map on $H_2(-;\mathbb Q)$ is an
isomorphism.  Consequently all rational degree-two classes on $BM$
are pulled back from its maximal hull.

Here the finite classification of possible proper relation spaces
can be made very small.  There are $27^2=729$ transformation pairs.
Their distinct proper relation spaces number $81$: one has rank
zero, $53$ have rank two and $27$ have rank three.  Closing these
under sums of rank less than four gives $111$ spaces, distributed
as $1,53,57$ in ranks $0,2,3$.  Row and column permutations and
transposition reduce them to the following eleven orbits.
A tuple denotes the four entries of a balanced matrix, in row order.
\[
\begin{array}{c|r|r|r|l}
\text{orbit}&\text{size}&\dim H^2&|H|&\text{basis of balanced forms}\\\hline
1&9&1&101&(0,0,0,1)\\
2&18&1&33&(0,0,-1,1)\\
3&6&2&31&(0,0,1,0),\ (0,0,0,1)\\
4&18&1&19&(0,-1,1,1)\\
5&6&1&30&(0,-1,1,0)\\
6&18&2&12&(0,-1,1,0),\ (0,0,0,1)\\
7&18&2&14&(0,1,0,0),\ (0,0,1,0)\\
8&6&1&15&(-1,1,1,2)\\
9&9&2&11&(1,-1,1,0),\ (-1,1,0,1)\\
10&2&2&12&(1,1,1,0),\ (-1,0,0,1)\\
11&1&4&10&\text{the four matrix units}.
\end{array}
\]
This classification is an exact finite computation, not a census
assumption.  The independent checker
\path{code/b_verify_rational_3x3_cups.py} recomputes the individual
relation spaces using rational fractions, verifies closure of the
claimed list under every proper sum, and checks that each listed
space is generated by the actors in its hull.  These two inclusions
prove exhaustiveness.  It separately checks all permutation orbits.
The complete integer bases and actor lists are in
\path{code/outputs/b_rational_action_hulls_3x3.json} and
\path{code/outputs/b_rational_action_hull_orbits_3x3.json}.

Orbits 1 and 2 are covered by Lemma~\ref{B:rank-one-square}.
For orbits $3,4,7,8,9,10,11$, the same independent checker proves
$H_4(-;\mathbb F_5)=0$, hence $H^4(-;\mathbb Q)=0$.
More explicitly, in the two-sided bar complex of
Lemma~\ref{lem:B-band-Tor}, let $q=|H|-10$ be the number of
nonidentity elements outside the band.  Its degree-two term has
rank $4q^2$.  The independently checked lower bounds for the ranks
of the maps into and out of this term are
\[
\begin{array}{c|rrrrrrr}
 \text{orbit}&3&4&7&8&9&10&11\\\hline
 \operatorname{rank}d_2&82&33&14&17&2&6&0\\
 \operatorname{rank}d_3&1682&291&50&83&2&10&0.
\end{array}
\]
In each column their sum is $4q^2$.  The certificate supplies
independent column indices; the checker reconstructs those matrices
and uses separate finite-field elimination.  Associativity and the
bar differential identity ensure that these lower bounds imply
vanishing, without a numerical rank approximation.

Orbit 6 has nonzero fourth homology, but all its products vanish.
The two affine cocycles for the displayed forms are
\[
 u_L(m,n)=L\bigl(\tau_m(0),\sigma_n(0)\bigr),
\]
where row or column zero has value zero.  The certificate
\path{code/outputs/a_rational_hull_orbit6_cups.json}
gives integral three-cochains for their two squares and mixed
product.  Their nonzero supports have sizes $110,260,39$ and
maximum absolute coefficients $1,2,1$.  Each of the three
coboundary identities is checked on all $11^4=14641$ normalized
quadruples.  Independent modular rank bounds and two integral band
cycles with pairing determinant $-1$ also verify that these classes
form the entire rational degree-two cohomology.

It remains to treat orbit 5, where the finite-type hypothesis is
essential.  Its form is
\[
 L=\begin{pmatrix}0&-1\\1&0\end{pmatrix}.
\]
Let $u=u_L$ be its integral affine cocycle.  Number the thirty
transformation pairs satisfying \eqref{B:3x3-balance}
lexicographically, moving the identity to index zero.
For any actor $m$, let $G_m$ be its eventual cyclic group, let $e$
be that group's identity, and put $g=em$.  The restriction of $[u]$
to $BG_m$ is represented by the residue
\begin{equation}\label{B:3x3-cyclic-character}
 c_m=\sum_{h\in G_m}u(h,g)\pmod{|G_m|}.
\end{equation}
Indeed $u(e,h)=u(h,e)=u(e,e)$; subtracting this constant normalizes
the restricted cocycle and changes the sum by a multiple of $|G_m|$.
The elementary cyclic-group cocycle formula then identifies the
residue with its class in $H^2(BG_m;\mathbb Z)$.
Exactly eight actors have nonzero residue:
\[
 10,11,17,18,20,21,25,26.
\]
Six have eventual group of order two and two have group of order three.
The checker verifies these characters directly from the full table.

No such actor can lie in the image of a finite monoid with finite
CW classifying space.  If $m\in M$ mapped to it, the eventual cyclic
group of $m$ would map onto $G_m$.  The nonzero character remains
nonzero after inflation: for a finite group $G$,
$H^2(BG;\mathbb Z)=\operatorname{Hom}(G,\mathbb Q/\mathbb Z)$,
and precomposition with a surjection is injective.  But the resulting
integral class is pulled back from $BM$, and Miller's mapping-space
theorem makes the map from this finite local group's classifying
space to a finite CW model of $BM$ null-homotopic.  This is a
contradiction.  A local identity is interpreted by adjoining a fresh
identity before mapping to $M$, as in
Proposition~\ref{B:two-row-torsion}.

Removing those eight actors leaves twelve optional actors besides
the band and the identity.  The independent checker tests all
$2^{12}=4096$ subsets.  Exactly $154$ are closed submonoids, and
every one lies in one of nine maximal submonoids of order sixteen.
Their outside-band actor sets are
\[
\begin{gathered}
\{0,4,5,7,9,16,19\},\quad
\{0,4,5,6,7,8,9\},\quad
\{0,5,6,12,16,22,24\},\\
\{0,5,9,16,19,22,24\},\quad
\{0,6,7,8,9,23,24\},\quad
\{0,6,7,12,16,23,24\},\\
\{0,4,5,6,7,12,16\},\quad
\{0,5,6,8,9,22,24\},\quad
\{0,7,9,16,19,23,24\}.
\end{gathered}
\]
For every one, the two-sided bar ranks over $\mathbb F_5$ are
$\operatorname{rank}d_2=21$ and
$\operatorname{rank}d_3=123$.  Their sum is the degree-two chain
rank $144$, so all nine have $H_4(-;\mathbb F_5)=0$.
The actual action map factors through one of these nine targets;
therefore its pulled-back square is zero over $\mathbb Q$.
The safe-submonoid lists, tables and independent rank witnesses are
in \path{code/outputs/b_rational_hull5_finite_type_pruning.json},
\path{code/outputs/b_rational_hull5_safe_models.json}, and
\path{code/outputs/b_rational_3x3_rank_certificates.json}.
The complete independent verification result is
\path{code/outputs/b_rational_3x3_cups_verified.json}.
This finishes all eleven cases for a minimal band.

Finally let the minimal ideal of the original monoid be
$I\times H\times J$.  Apply the actual minimal-group quotient of
Lemma~\ref{lem:B-collapse-minimal-group}; it is a rational cohomology
ring isomorphism.  Its band actions and their rational balanced forms
are subject to the same classification.  In the exceptional case,
pull the integral hull cocycle back to the original $M$ before using
\eqref{B:3x3-cyclic-character}.  The eventual local group in $M$
still surjects onto the forbidden cyclic group, so the same Miller
argument applies.  Thus no finite-type assertion about the quotient
is needed, and the result follows for arbitrary $H$.
\end{proof}

\begin{remark}\label{B:three-by-three-relative}
The last proof also applies when a homomorphism from the original
finite monoid has an action image containing a three-by-three band,
and the degree-two classes under consideration factor through that
image.  The finite-CW hypothesis is imposed on the original monoid:
its local groups supply the same forbidden-character obstruction.
This relative form permits quotienting duplicate row or column
values in a balanced-form computation.
\end{remark}

\begin{proposition}[A finite model with square two]
\label{B:even-hopf23}
There is a $23$-element monoid with classifying space
$S^2\cup_{2\eta}e^4$, where $\eta$ generates $\pi_3(S^2)$.
Its minimal ideal is a three-by-four rectangular band.
\end{proposition}
\begin{proof}
Let $I=\{0,1,2\}$, $J=\{0,1,2,3\}$, and put
\[
 L=\begin{pmatrix}0&-1\\2&1\\0&1\end{pmatrix}.
\]
Define $H$ to be the monoid of all pairs of transformations
$(\sigma:I\to I,\tau:J\to J)$ satisfying
$B_\tau^tL=LA_\sigma$, with the augmentation matrices and
multiplication used in Theorem~\ref{B:three-by-three-cups}.
The equation is closed under multiplication.  The full list has
$23$ elements and includes the twelve constant pairs, which form its
minimal band.  These assertions, including all $23^3$ associativity
triples, are checked in
\path{code/b_verify_rational_3x4_cups.py}.
Number the identity first, then the twelve constant pairs
lexicographically, and then the remaining pairs lexicographically;
write $u_i$ for element $i$ and $R=\mathbb ZH$.

Put $P=Ru_1$, $E=R(u_{14}-u_1)$, $F=R(1-u_{14})$ and
$G=R(u_{15}-u_3)$.  The projective resolution has the following terms:
\[
\begin{array}{c|c|r}
 n&P_n&\operatorname{rank}_{\mathbb Z}P_n\\\hline
0&P&3\\
1&E&12\\
2&P\oplus F\oplus E&23\\
3&G\oplus G\oplus E&22\\
4&P\oplus E&15\\
n\geq5&E&12.
\end{array}
\]
All displayed ring elements are idempotents.  The complete integral
matrices are supplied by
\path{code/outputs/b_resolution_rational_hull_3x4_orbit28.json}.
The independent resolution verifier checks every module action,
saturated image, kernel rank, and zero composite.  From degree six
the maps alternate between $1+S$ and $1-S$, where $S$ is the
involution of the twelve-element displayed basis with transpositions
\[
 (0\ 1)(2\ 3)(4\ 5)(6\ 8)(7\ 10)(9\ 11).
\]
This is an exact two-periodic tail over $\mathbb Z$, and the checker
also verifies its $R$-linearity and its exact attachment to the prefix.
Only $P_0,P_2,P_4$ contain an augmentation-one projective summand.
Tensoring therefore gives one copy of $\mathbb Z$ in degrees zero,
two and four, with zero differentials, and zero in all other degrees.

The integral affine cocycle
$u(m,n)=L(\tau_m(0),\sigma_n(0))$, with zero row and column zero,
has primitive degree-two class and square twice a primitive
class in degree four, up to sign.  This assertion is established by
an actual comparison to the normalized free left bar resolution,
not inferred from the homology groups.  Its integral contraction is
\[
 s(a;m_1,\ldots,m_n)=
 \begin{cases}
 (1;a,m_1,\ldots,m_n),&a\ne1,\\
 0,&a=1.
 \end{cases}
\]
For an idempotent projective generator $\varepsilon$, the comparison
is defined recursively by
$F_n(\varepsilon)=\varepsilon s(F_{n-1}(d\varepsilon))$.
The stored comparison in
\path{code/outputs/b_resolution_rational_hull_3x4_orbit28_bar_comparison.json.gz}
has respectively $3,48,332,2008,3636$ nonzero integral terms in
degrees zero through four.  A separate checker verifies every chain
square and every $H$-equivariance equation.  Evaluation of $u$ on the
degree-two trivial-coefficient generator is $-1$, and evaluation of
$u\smile u$ on the degree-four generator is $-2$.
Thus, after orienting the generators, the ring has $x^2=2y$.
The complete verification result is
\path{code/outputs/b_rational_3x4_cups_verified.json}.

A monoid with a rectangular minimal band has trivial universal group:
every band element becomes the identity, and multiplication of any
monoid element by a band element then kills that element too.
Hence $BH$ is simply connected.  Choose a sphere representing its
$H_2$ generator.  The mapping-cylinder pair $(BH,S^2)$ is
three-connected: relative Hurewicz first identifies its zero relative
$H_3$ with relative $\pi_3$.  Relative Hurewicz in degree four then
allows a single four-cell to be attached to $S^2$, producing a map
to $BH$ that is an integral homology equivalence.  Homological
Whitehead makes it a homotopy equivalence.  The attaching map has
Hopf invariant two by the computed cup square, and
$\pi_3(S^2)=\mathbb Z\eta$, so its sign can be chosen as $2\eta$.
This proves the claimed homotopy type.
\end{proof}

\begin{theorem}[Three-by-four target obstructions]
\label{B:three-by-four-targets}
A finite monoid with a three-by-four rectangular minimal band cannot
have classifying space $\mathbb CP^2$.
A finite monoid whose minimal ideal has three rows and four columns,
with any maximal group, cannot have classifying space $S^2\times S^2$.
Both assertions also hold after interchanging rows and columns.
\end{theorem}
\begin{proof}
Use the maximal action-hull construction of
Theorem~\ref{B:three-by-three-cups}, now with balanced matrices of
size $3\times2$.  There are $27\cdot256=6912$ transformation pairs.
Their distinct proper rational relation spaces number $853$, with
ranks $0,2,3,4,5$ occurring $1,6,282,504,60$ times.
Closing under proper sums gives $2767$ spaces.  Their balanced-form
spaces have dimensions $6,4,3,2,1$, occurring respectively
$1,6,282,1278,1200$ times.  Row and column permutations give $53$
orbits, of which $44$ have degree-two rank one or two.
Only those forty-four can occur for the two target spaces.

The classification is independently certified in
\path{code/b_verify_rational_3x4_lattice.py}.
It uses finite-field canonical keys only to propose a containing
space; exact integer annihilator equations and the reverse modular
rank inequality then prove each equality over $\mathbb Q$.
It checks every proper sum, every maximal actor list, all rational
ranks of individual relation spaces, and all permutation orbits.
The full certificate is
\path{code/outputs/b_rational_3x4_lattice_verified.json}.
The integer bases and explicit transformation pairs for the fifty-three
numbered orbits are in
\path{code/outputs/b_rational_action_hull_orbits_3x4.json}.
Thus no source-enumeration completeness is used here.

Twenty-one of the forty-four relevant orbits reduce immediately.
If every matrix in the balanced-form space has rank at most one,
Lemma~\ref{B:rank-one-square} kills every square, and rational
polarization kills every mixed product.  This condition is checked
on basis vectors and their pairwise sums, which suffices because
all two-by-two minors are quadratic polynomials.
Otherwise the column values have a common duplicate: two columns
have equal values for every balanced form.  The balancing equation
makes this equivalence relation invariant under every column action,
so the action monoid descends to at most three column values.
Every affine two-cocycle factors through the resulting action map.
If two values remain, apply the two-row/column calculation; if three
remain, apply the relative form in
Remark~\ref{B:three-by-three-relative}.
The image contains the full smaller rectangular band, and all
local-group exclusions are applied to the original finite-type
source.  All products in these cases are zero over $\mathbb Q$.

There are twenty-three remaining orbits.  Eighteen have
$H_4(-;\mathbb F_5)=0$ and therefore zero rational products.
For reference, the full list of independently verified two-sided
bar ranks is as follows.  Each rank sum is
$6(|H|-13)^2$, the degree-two chain rank.
\[
\begin{array}{l|r|r|r}
\text{orbit numbers}&|H|&\operatorname{rank}d_2&\operatorname{rank}d_3\\\hline
22,24,26,29,35&19&31&185\\
23&21&43&341\\
25,33,50&18&25&125\\
30&34&121&2525\\
31,41,42,43&17&20&76\\
32,40&18&26&124\\
49&21&44&340\\
52&15&8&16.
\end{array}
\]
Selected-column witnesses and separate dense finite-field elimination
verify these ranks.  The checker also reconstructs the target tables
from the transformation pairs, so the rank witnesses concern the
specified hulls.

Orbits 27 and 38 have orders 25 and 21 and balanced-form bases
\[
\begin{split}
27:;&(0,-1,2,1,1,0),\\
38:;&(0,-1,1,0,0,0),\ (0,0,0,0,0,1).
\end{split}
\]
In each, two actors have a nonzero local cyclic character of order
two, computed by \eqref{B:3x3-cyclic-character}.  They cannot occur
in the image of a finite-type source, by the same integral
character-inflation and Miller argument.  Directly testing all
$1024$ and $64$ subsets of the remaining optional actors gives
respectively $229$ and $49$ closed submonoids.  In each case they
are all contained in a single maximal allowed submonoid, of orders
23 and 19.  Their bar ranks are respectively $(55,545)$ and
$(32,184)$, proving $H_4(-;\mathbb F_5)=0$ in both cases.

Orbits 34 and 37 each have order 17 and degree-two rank two.
Their displayed balanced bases are
\[
\begin{split}
34:;&(0,-1,1,-1,1,0),\ (0,0,0,1,0,1),\\
37:;&(0,0,0,1,0,0),\ (0,-1,1,0,0,1).
\end{split}
\]
All squares and mixed products are integral coboundaries.
The certificates
\path{code/outputs/a_rational_hull_3x4_orbit34_cups.json} and
\path{code/outputs/a_rational_hull_3x4_orbit37_cups.json}
supply all six integral three-cochains.  Each identity is checked on
all $16^4=65536$ normalized quadruples.  Independent rank bounds and
integral band-cycle pairings also certify that the two forms span
the full rational degree-two cohomology in each hull.

The sole remaining orbit is 28, of degree-two rank one and order
23.  It is precisely the monoid of Proposition~\ref{B:even-hopf23}.
This case cannot be excluded by rational square-vanishing: its
primitive class has square twice a primitive integral four-class.
Suppose a monoid with a three-by-four minimal band had classifying
space $\mathbb CP^2$ and this action hull.  The action map induces
an integral \emph{surjection} on $H_2$: both balanced presentations
start with the same lattice $W\otimes V$, and the larger hull only
adds relations.  Since source and target $H_2$ are both $\mathbb Z$,
this map is an isomorphism.  The primitive degree-two class of the
source is therefore pulled back, up to sign, from the primitive
class of the hull.  Its square would be divisible by two in
$H^4(\mathbb CP^2;\mathbb Z)=\mathbb Z$, a contradiction.
This proves the projective-plane assertion for a minimal band.
The full target-specific coverage, twenty rank certificates, two
local-character prunings, six cup primitives and the exceptional
integral comparison are independently verified together by
\path{code/b_verify_rational_3x4_cups.py}.

For the product target, its rational degree-two rank is two, so
orbit 28 cannot occur.  Every remaining rank-two case has zero
rational products, contrary to the nonzero cross product of
$S^2\times S^2$.  If the minimal group is nontrivial, apply
Lemma~\ref{lem:B-collapse-minimal-group} and pull every exceptional
integral hull character back to the original monoid before using
Miller's theorem, exactly as in the last paragraph of
Theorem~\ref{B:three-by-three-cups}.  Rational cohomology and its
products are preserved.  This proves the product assertion for
arbitrary minimal maximal groups.
\end{proof}

\begin{lemma}[Three-dimensional quotient algebras]\label{B:small-quotient-algebra}
Let $A$ be a unital finite-dimensional $\mathbb Q$-algebra with
$\dim_{\mathbb Q}A\leq3$, let $W$ be a finite-dimensional right
$A$-module, and let $V$ be a finite-dimensional left $A$-module.
If $\operatorname{Tor}_2^A(W,V)\ne0$, then
\[
  \operatorname{Tor}_n^A(W,V)\ne0\qquad\text{for every }n\geq2.
\]
Consequently, if these Tor groups vanish in all sufficiently high
degrees, then $\operatorname{Tor}_2^A(W,V)=0$.
\end{lemma}

\begin{proof}
Write $J=\operatorname{rad}(A)$.  If $J=0$, the algebra is semisimple
and all positive Tor groups vanish.  Suppose $J\ne0$.
The dimension bound gives the following complete list of possibilities:
\[
 \mathbb Q[x]/(x^2),\quad
 \mathbb Q\times\mathbb Q[x]/(x^2),\quad
 \mathbb Q[x]/(x^3),\quad
 \mathbb Q\oplus J\ (\dim J=2, J^2=0),\quad
 T_2(\mathbb Q).
 \tag{\ref{B:small-quotient-algebra}.1}
\]
Here $T_2(\mathbb Q)$ denotes the algebra of upper triangular
$2$ by $2$ matrices.  We first justify the list, including the possible
noncommutative case.

If $A/J\cong\mathbb Q$, the scalar copy of $\mathbb Q$ splits the
quotient and $A=\mathbb Q\oplus J$ as a vector space.  A
one-dimensional nilpotent radical has square zero and gives the dual
numbers.  For a two-dimensional radical, either $J^2=0$, or
$\dim J^2=1$ and $J^3=0$.  In the latter case the nonzero
multiplication
$(J/J^2)\otimes(J/J^2)\longrightarrow J^2$ implies that a lift
$x$ of any nonzero element of $J/J^2$ has $x^2\ne0$.
Thus $1,x,x^2$ is a basis and $x^3=0$, giving
$\mathbb Q[x]/(x^3)$.

The only remaining nonsemisimple case has
$\dim(A/J)=2$ and $\dim J=1$.  The semisimple quotient cannot be
a quadratic field: the nonzero radical would then be a one-dimensional
$\mathbb Q$-vector space carrying a unital action of that field.
Hence $A/J\cong\mathbb Q\times\mathbb Q$ and $J^2=0$.
Lift the two orthogonal idempotents of the quotient to $e$ and $1-e$.
For completeness, if $p$ is any lift of the first idempotent and
$\delta=p^2-p\in J$, then
$e=p+(1-2p)\delta$ is idempotent because $J^2=0$.
The one-dimensional bimodule $J$ has exactly one nonzero component
among $eJ e$, $eJ(1-e)$, $(1-e)Je$, and $(1-e)J(1-e)$.
A diagonal component gives a product of the dual numbers with
$\mathbb Q$; an off-diagonal component gives $T_2(\mathbb Q)$,
up to interchanging the two idempotents.  This proves the list.

The triangular algebra has global dimension at most one.  One can see
this without a classification theorem: its modules are representations
of a quiver consisting of one arrow.  For a representation
$f:V_1\longrightarrow V_2$, let $P_1=(\mathbb Q\to\mathbb Q)$
and $P_2=(0\to\mathbb Q)$ be the two vertex projectives.  The
sequence
\[
 0\longrightarrow P_2\otimes V_1
 \longrightarrow (P_1\otimes V_1)\oplus(P_2\otimes V_2)
 \longrightarrow (V_1\xrightarrow{f}V_2)\longrightarrow0
\]
is exact; at the second vertex its maps are
$v\mapsto(v,-f(v))$ and $(v,w)\mapsto f(v)+w$.
Thus this case cannot have nonzero $\operatorname{Tor}_2$.

Consider next $A=\mathbb Q[x]/(x^m)$ with $m=2$ or $3$.
Every finite-dimensional module is a direct sum of the cyclic modules
$A/(x^a)$, $1\leq a\leq m$, by the Jordan decomposition of its
nilpotent $x$-action.  The summand with $a=m$ is free.  If
$1\leq a<m$, the module $A/(x^a)$ has a free resolution whose
successive differentials alternate between multiplication by
$x^a$ and $x^{m-a}$.  Tensoring with $A/(x^b)$, for
$1\leq b<m$, gives in every positive degree homology of dimension
\[
  \min(a,b)+\min(m-a,b)-b
   =\min\{a,b,m-a,m-b\}>0.
\]
Positive Tor therefore vanishes when either cyclic summand is free,
and is nonzero in every positive degree when both are nonfree.
Direct sums prove the assertion for these algebras.  They also prove
it for $\mathbb Q\times\mathbb Q[x]/(x^2)$, since Tor splits by
the two central idempotents and the $\mathbb Q$ factor contributes
no positive-degree Tor.

Finally suppose $A=\mathbb Q\oplus J$, $\dim J=2$, and $J^2=0$.
Set $k=A/J$.  A minimal free cover of a finite module has kernel
contained in the radical of the free module.  Thus, for a nonprojective
left module $V$, its first syzygy is a nonzero semisimple module
$\Omega V\cong k^r$, $r>0$.  The same holds on the right.
In particular, if a right module $W$ is nonprojective, its minimal
free resolution never terminates: the first syzygy is nonzero
semisimple, and the syzygy of $k^s$ is $J^s\cong k^{2s}$.
All differentials in a minimal resolution become zero after
tensoring with $k$, so
$\operatorname{Tor}_i^A(W,k)\ne0$ for every $i\geq1$.
The hypothesis $\operatorname{Tor}_2^A(W,V)\ne0$ forces both
$W$ and $V$ to be nonprojective.  Dimension shifting now gives
\[
 \operatorname{Tor}_n^A(W,V)
 \cong\operatorname{Tor}_{n-1}^A(W,\Omega V)
 \cong\operatorname{Tor}_{n-1}^A(W,k)^r\ne0
 \qquad(n\geq2).
\]
This settles the last case and proves the lemma.
\end{proof}

\begin{corollary}\label{B:four-outside-minimal-band}
Let $M$ be a finite monoid whose minimal ideal $K$ is a rectangular
band.  If $BM$ has finite-dimensional homotopy type and
$H_4(BM;\mathbb Q)\ne0$, then $|M\setminus K|\geq4$.
\end{corollary}

\begin{proof}
Put $A=\mathbb QM/\mathbb QK$, so
$\dim_{\mathbb Q}A=|M\setminus K|$.
The natural dimension-shifting identification in
Lemma~\ref{lem:B-band-Tor} expresses $H_n(BM;\mathbb Q)$, for
$n\geq2$, as $\operatorname{Tor}_{n-2}^A(W,V)$ for the finite
row and column augmentation modules $V,W$.
If $\dim A\leq3$, nonzero $H_4$ and
Lemma~\ref{B:small-quotient-algebra} would force nonzero
$H_n(BM;\mathbb Q)$ for every $n\geq4$, contradicting finite
homotopy dimension.
\end{proof}

\begin{lemma}[Minimal resolutions over a local algebra]
\label{B:local-minimal-Tor}
Let $A$ be a finite-dimensional local algebra over a field, and put
$D=A/\operatorname{rad}(A)$.  If a finite-dimensional right
$A$-module $W$ is nonprojective, then
\[
 \operatorname{Tor}_n^A(W,D)\ne0\qquad(n\geq1).
\]
The analogous assertion holds for left modules.
\end{lemma}

\begin{proof}
Finite projective modules over $A$ are free.  In a minimal free
resolution all differential entries belong to $J=\operatorname{rad}(A)$.
Such a resolution cannot terminate with a nonzero free module in
positive degree: a nonzero element $a$ with $Ja=0$ gives a nonzero
kernel vector in any chosen free summand of the last differential.
An element $a$ exists by taking a nonzero element in the last nonzero
power of $J$; the semisimple case has no nonprojective modules.
Thus every positive-degree term of the minimal resolution is nonzero.
Tensoring with $D$ makes every differential zero and proves the claim.
\end{proof}

\begin{lemma}[The nonlocal four-dimensional case]
\label{B:nonlocal-four-Tor}
Let $A$ be a nonlocal finite-dimensional $\mathbb Q$-algebra of
dimension at most four.  For finite-dimensional right and left modules
$W,V$, respectively, nonzero $\operatorname{Tor}_2^A(W,V)$ implies
nonzero $\operatorname{Tor}_n^A(W,V)$ in arbitrarily high degrees.
\end{lemma}

\begin{proof}
Semisimple factors contribute no positive Tor.  Write $J=\operatorname{rad}(A)$.
If the semisimple quotient has a field factor of degree greater than
one, the dimension bound leaves room only for a quadratic field
together with a rational factor and a one-dimensional radical.
That radical must be supported on the rational factor: a nonzero
one-dimensional rational vector space cannot carry a unital action
of a quadratic field.  The algebra is therefore a product of a
semisimple factor with the dual numbers, or is semisimple.
These cases satisfy the assertion by
Lemma~\ref{B:small-quotient-algebra}.

It remains to consider $A/J\cong\mathbb Q^r$, where $r\geq2$.
Lift the primitive orthogonal idempotents of this quotient.  The
radical has dimension at most two.  If $J^2\ne0$, necessarily
$r=2$, $\dim J=2$, $\dim J^2=1$, and $J^3=0$.
The one-dimensional bimodule $J/J^2$ must be a loop at one of the
two idempotents: an off-diagonal one-dimensional component has
square zero.  A lift $x$ of that loop satisfies $x^2\ne0$ and
$x^3=0$.  There are no radical components joining the two
idempotents, so
$A\cong\mathbb Q\times\mathbb Q[x]/(x^3)$.
The cyclic-module calculation in
Lemma~\ref{B:small-quotient-algebra} again applies.

Suppose now that $J^2=0$.  The lifted idempotents identify $A$ with
the algebra of a quiver with $r$ vertices and $d=\dim J$ arrows,
modulo all paths of length two.  Here $r+d\leq4$ and $r\geq2$.
Let $\Omega W$ and $\Omega V$ be the kernels of minimal projective
covers.  They are semisimple.  Since $J$ annihilates $\Omega V$,
the map
$\Omega W\otimes_A\Omega V\to P_W\otimes_A\Omega V$
is zero.  Dimension shifting consequently gives
\[
 \operatorname{Tor}_n^A(W,V)
 \cong \operatorname{Tor}_{n-2}^A(\Omega W,\Omega V)
 \qquad(n\geq2),
\]
including degree two, where the right side is the ordinary tensor
product.  If it is nonzero, some vertex $v$ occurs in both
$\Omega W$ and $\Omega V$.  The first occurrence requires an
arrow entering $v$, and the second an arrow leaving $v$.
With at most two arrows and at most $4-d$ vertices, such a vertex
lies on a directed cycle: either an incident arrow is a loop, or
the two arrows form a two-cycle.
For a radical-square-zero quiver algebra, the minimal resolution
of a vertex simple has one generator for every directed path of
the appropriate length.  Repeating the cycle through $v$ therefore
gives nonzero
$\operatorname{Tor}_{m\ell}^A(S_v^{\mathrm{right}},S_v^{\mathrm{left}})$
for every $m\geq0$, where $\ell=1$ or $2$ is its length.
The displayed dimension shift proves the assertion.
\end{proof}

\begin{corollary}[A one-dimensional simple target]
\label{B:four-dimensional-simple-target}
Let $A$ be any finite-dimensional $\mathbb Q$-algebra with
$\dim A\leq4$, let $W$ be a finite-dimensional right module, and
let $V$ be a one-dimensional simple left module.
If $\operatorname{Tor}_2^A(W,V)\ne0$, then
$\operatorname{Tor}_n^A(W,V)\ne0$ for arbitrarily large $n$.
The same conclusion holds with the sides interchanged.
\end{corollary}

\begin{proof}
The nonlocal case is Lemma~\ref{B:nonlocal-four-Tor}.
If $A$ is local, the existence of a one-dimensional simple module
forces $A/\operatorname{rad}(A)\cong\mathbb Q$, and $V$ is that
quotient.  Nonzero positive Tor forces $W$ to be nonprojective,
so Lemma~\ref{B:local-minimal-Tor} applies.
\end{proof}

\begin{lemma}[The invertible binary pencil case]
\label{B:binary-pencil-Tor}
Let $k$ be either $\mathbb Q$ or $\mathbb C$, and let $A_B$ be the
four-dimensional $k$-algebra with basis $1,x,y,c$ and multiplication
\[
 x_i x_j=b_{ij}c,\qquad (x_1,x_2)=(x,y),\qquad
 c^2=cx=xc=cy=yc=0,
\]
where $B=(b_{ij})\in\operatorname{Mat}_2(\{0,1\})$ is invertible.
For finite-dimensional right and left $A_B$-modules $W,V$, respectively,
if $\operatorname{Tor}_2^{A_B}(W,V)\ne0$, then
$\operatorname{Tor}_n^{A_B}(W,V)\ne0$ in arbitrarily high degrees.
\end{lemma}

\begin{proof}
We first work over $k=\mathbb C$ and write $A=A_B$ and
$J=\langle x,y,c\rangle$.  Then $J^2=kc$, $J^3=0$, and $A/J=k$.
The bilinear pairing
\[
 A\times A\longrightarrow k,\qquad
 (a,b)\longmapsto\text{the coefficient of $c$ in $ab$}
\]
is nondegenerate: it pairs $1$ with $c$ in both directions, and its
restriction to $\langle x,y\rangle$ has matrix $B$.
Consequently $A$ is Frobenius and is injective as both a left and a
right module over itself.  More explicitly, the pairing identifies
the appropriate regular module with $\operatorname{Hom}_k(A,k)$,
which is injective because its module Hom functor is the exact
vector-space dual functor.  The left and right socles of $A$ are
both $kc$.  In particular, every nonzero left ideal contains $kc$.

We describe the nonprojective indecomposable left modules.
If $M$ is such a module and $cm\ne0$ for some $m\in M$, then
the homomorphism $A\to M$, $a\mapsto am$, is injective: any
nonzero kernel would contain $c$.  Its image splits because $A$
is injective, contradicting the assumed indecomposability and
nonprojectivity of $M$.  Thus
\[
 cM=0.                                                     \tag{\ref{B:binary-pencil-Tor}.1}\label{eq:B-binary-c-kills}
\]
Moreover, unless $M=k$, one has $\operatorname{soc}M=JM$.
Indeed, $JM\subseteq\operatorname{soc}M$ by
\eqref{eq:B-binary-c-kills}.  A socle
vector outside $JM$ spans a simple direct summand: choose a
vector-space complement containing $JM$, which is automatically
an $A$-submodule since $J$ sends the whole module into $JM$.

Choose a complement $T$ to $S=JM$ in $M$.  The actions of $x,y$
are therefore a pair of linear maps
\[
 X,Y:T\longrightarrow S,
\]
with $\operatorname{im}X+\operatorname{im}Y=S$ and, for $M\ne k$,
$\ker X\cap\ker Y=0$.
Isomorphism of modules is exactly simultaneous equivalence of
these two matrices, and a direct-sum decomposition of the pencil
is exactly a direct-sum decomposition of the module.
The Kronecker canonical form gives the following complete list
of indecomposable possibilities over $\mathbb C$:
\[
\begin{array}{c|c|c}
 \text{module}&(\dim T,\dim S)&(X,Y)\\ \hline
 I_n&(n+1,n)&([I_n\ 0],[0\ I_n]),\quad n\geq1\\
 P_n&(n,n+1)&([I_n\ 0]^{\mathsf T},[0\ I_n]^{\mathsf T}),
                         \quad n\geq1\\
 R_{n,\lambda}&(n,n)&(I_n,J_n(\lambda)),
                         \quad n\geq1,\ \lambda\in\mathbb C\\
 R_{n,\infty}&(n,n)&(J_n(0),I_n),\quad n\geq1.
\end{array}                                                \tag{\ref{B:binary-pencil-Tor}.2}\label{eq:B-binary-pencil-list}
\]
There is also the simple module $k$, with dimensions $(1,0)$.
The other zero-size singular pencil, of dimensions $(0,1)$,
is excluded by $S=JM$.
Here $J_n(\lambda)$ is a single Jordan block.  The exact matrix
pencil theorem used in this list is recorded, with these singular
and regular blocks, in
\cite[Theorem~2.1, p.~221]{deterandmytryshyndopico2024}.

We next determine syzygies, keeping track of indecomposability.
If $a=\dim(M/JM)$, the minimal projective cover is $A^a\to M$.
Its kernel $\Omega M$ lies in $JA^a$, is killed by $c$, and
contains $cA^a=\operatorname{soc}(A^a)$.  Thus $A^a$ is the
injective envelope of $\Omega M$.  If $\Omega M$ were a direct
sum of two nonzero modules, its injective envelope, and hence
the cokernel $M$, would split accordingly.  Neither resulting
cokernel can be zero: a nonzero injective summand is a sum of
copies of $A$, whereas $c$ kills $\Omega M$.
It follows that $\Omega M$ is again indecomposable and
nonprojective.

Put $b=\dim JM$.  The socle of $\Omega M$ has dimension $a$
and $\dim\Omega M=3a-b$.  Unless $\Omega M=k$, its top and
radical dimensions are consequently $(2a-b,a)$.
If $\Omega M=k$, these same dimension equalities force
$(a,b)=(1,2)$, so $M=P_1$.
The list \eqref{eq:B-binary-pencil-list}, and the direct
calculation $\Omega k=J\cong I_1$,
now give
\[
 \Omega k\cong I_1,\qquad
 \Omega I_n\cong I_{n+1},\qquad
 \Omega P_1\cong k,\qquad
 \Omega P_n\cong P_{n-1}\quad(n\geq2).                     \tag{\ref{B:binary-pencil-Tor}.3}\label{eq:B-binary-string-syzygies}
\]
Thus all the singular-pencil modules are positive or negative
syzygies of the simple module; the displayed formulas suffice
without choosing a complete resolution.

For the regular pencils there is an explicit formula that also
proves periodicity.  Identify $T=S=k^n$, and initially take
$(X,Y)=(I_n,J_n(\lambda))$ or $(J_n(0),I_n)$.
More generally the following calculation applies whenever $X,Y$
commute and at least one is invertible.  In the cover
$A\otimes_k T\to M$, the kernel is the direct vector-space sum
of $c\otimes T$ and the vectors
\[
 q(v)=y\otimes Xv-x\otimes Yv,\qquad v\in T.
\]
Indeed, modulo $c\otimes T$ its kernel is that of
$(u,v)\mapsto Xu+Yv$, and the displayed vectors form a basis
of that kernel.  Using $q(v)$ as the top coordinates and
$c\otimes v$ as the radical coordinates, multiplication gives
\[
 \Omega(X,Y)=
 \bigl(b_{12}X-b_{11}Y,\ b_{22}X-b_{21}Y\bigr).             \tag{\ref{B:binary-pencil-Tor}.4}\label{eq:B-binary-regular-syzygy}
\]
In particular, if
\[
 D_B=\begin{pmatrix}b_{12}&-b_{11}\\ b_{22}&-b_{21}\end{pmatrix},
\]
the $r$th syzygy is represented by $D_B^r(X,Y)^{\mathsf T}$.
This iteration is legitimate in the same coordinates.  Both
matrices remain commuting linear combinations of the original
$I_n$ and its single Jordan block.  Since $D_B$ is invertible,
their eigenvalues cannot both vanish, so at least one of the
two matrices remains invertible at every step.

There are six invertible binary $2$ by $2$ matrices.
Four are symmetric; for each of them
$\operatorname{tr}D_B=0$ and $\det D_B=\det B$, whence
\[
 D_B^2=-\det(B)I_2.
\]
The other two, and their powers, are
\[
\begin{split}
 B=\begin{pmatrix}1&1\\0&1\end{pmatrix}:&\quad
 D_B=\begin{pmatrix}1&-1\\1&0\end{pmatrix},\quad D_B^3=-I_2,\\
 B=\begin{pmatrix}1&0\\1&1\end{pmatrix}:&\quad
 D_B=\begin{pmatrix}0&-1\\1&-1\end{pmatrix},\quad D_B^3= I_2.
\end{split}                                               \tag{\ref{B:binary-pencil-Tor}.5}\label{eq:B-binary-periods}
\]
Multiplying both matrices of a pencil by the same nonzero
scalar does not change its module isomorphism type, since one
may rescale the radical coordinates.  Every regular
indecomposable is therefore periodic under $\Omega$, with
period dividing two in the symmetric case and dividing three
in the other two cases.  The right-module calculation is the
same calculation for $A^{\mathrm{op}}=A_{B^{\mathsf T}}$.

Finally decompose $W,V$ into indecomposable summands and select
a pair with nonzero degree-two Tor.  Neither summand is
projective.  If either is a regular-pencil module of period
$q$, dimension shifting gives
\[
 \operatorname{Tor}_{2+rq}^{A}(W,V)
       \cong\operatorname{Tor}_{2}^{A}(W,V)\ne0
       \qquad(r\geq0)
\]
for the selected pair.  Suppose instead that both are
singular-pencil modules or simple modules.  By
\eqref{eq:B-binary-string-syzygies} there exist
nonnegative integers $r_W,s_W,r_V,s_V$ with
\[
 \Omega^{r_W}W\cong\Omega^{s_W}k,
 \qquad
 \Omega^{r_V}V\cong\Omega^{s_V}k.
\]
For all $n>r_W+r_V$, ordinary dimension shifting on the two
sides yields
\[
 \operatorname{Tor}_n^A(W,V)
 \cong
 \operatorname{Tor}_{n-r_W-r_V+s_W+s_V}^A(k,k).
\]
Every positive-degree group on the right is nonzero.  For
example, \eqref{eq:B-binary-string-syzygies} shows directly that the minimal free resolution
of $k$ never terminates, and all its differentials become zero
after tensoring with $k$.  This proves the required recurrence
over $\mathbb C$.

For $k=\mathbb Q$, extension of scalars along
$\mathbb Q\subset\mathbb C$ gives, in every degree,
\[
 \mathbb C\otimes_{\mathbb Q}\operatorname{Tor}_n^{A_B}(W,V)
 \cong
 \operatorname{Tor}_n^{\mathbb C\otimes A_B}
       (\mathbb C\otimes W,\mathbb C\otimes V).
\]
This follows by extending any free resolution, and the extension
is faithfully flat.  Nonzero degree-two Tor therefore survives
over $\mathbb C$, where the result just proved supplies nonzero
groups in arbitrarily high degrees.  Faithfulness brings their
nonvanishing back to $\mathbb Q$.
\end{proof}


\begin{theorem}[Contracted algebras of dimension four]
\label{B:contracted-four-Tor}
Let $N$ be a finite monoid with a zero element, and let
$A=\mathbb QN/\mathbb Q0$ be its contracted algebra.  Suppose
$\dim_{\mathbb Q}A\leq4$.
For every finite-dimensional right module $W$ and left module $V$,
nonzero $\operatorname{Tor}_2^A(W,V)$ implies nonzero
$\operatorname{Tor}_n^A(W,V)$ in arbitrarily high degrees.
\end{theorem}

\begin{proof}
Dimensions at most three are covered by
Lemma~\ref{B:small-quotient-algebra}, and a nonlocal algebra of
dimension four is covered by Lemma~\ref{B:nonlocal-four-Tor}.
Suppose $A$ is local and has dimension four.

The only nonzero unit of the monoid $N$ is its identity.  Otherwise
a nontrivial finite cyclic subgroup would give a nontrivial
idempotent in its rational group algebra, which embeds in $A$,
contradicting locality.  Similarly, the idempotent power of any
other nonzero monoid element must be zero: a different idempotent
basis element would again be a nontrivial idempotent of $A$.
The three nonidentity, nonzero basis elements therefore span an
ideal with quotient $\mathbb Q$, and this ideal is exactly
$J=\operatorname{rad}(A)$.

If $J^2=0$, the local square-zero argument in
Lemma~\ref{B:small-quotient-algebra} works without change with
$\dim J=3$.  If $\dim J^2=2$, the algebra is generated by a
single radical element and is $\mathbb Q[x]/(x^4)$; its cyclic
modules have the same two-periodic resolutions used in that lemma.
Indeed, a lift of a basis of $J/J^2$ generates $J$ modulo $J^2$,
and nilpotence successively shows that its powers span every
radical layer.  The four-dimensional algebra consequently has
basis $1,x,x^2,x^3$ with $x^4=0$.

The remaining case is $\dim J^2=1$, and nilpotence then gives
$J^3=0$.  Since products of monoid basis elements are either zero
or another basis element, $J^2$ is spanned by a single one of the
three basis elements, say $c$.  Write the other two as $x,y$.
Their products have the form
\[
 x_i x_j=b_{ij}c,\qquad x_1=x,\ x_2=y,\qquad
 B=(b_{ij})\in\operatorname{Mat}_2(\{0,1\}),\quad B\ne0.
\]
All products involving $c$ and an element of $J$ are zero.
If $B$ is invertible, Lemma~\ref{B:binary-pencil-Tor} proves the
claim.  It remains to treat the rank-one matrices.

Under a change of the basis of $J/J^2$ and a rescaling of $c$, a
rank-one bilinear form has one of two forms.  If its two linear
factors are proportional it is
$\left[\begin{smallmatrix}1&0\\0&0\end{smallmatrix}\right]$;
otherwise it is
$\left[\begin{smallmatrix}0&1\\0&0\end{smallmatrix}\right]$.
These changes are algebra isomorphisms, so we may use the two
displayed forms.

In the first case $x^2=c$ and $xy=yx=y^2=0$.
Every first left syzygy lies in a direct sum of copies of $J$,
so it is annihilated by $y$ and $c$ and is a module over
$\mathbb Q[x]/(x^2)$.  It is therefore a direct sum of the simple
module $k=A/J$ and $L=Ax=\langle x,c\rangle$.
Directly from multiplication,
\[
 \Omega L=\langle y,c\rangle\cong k^2,\qquad
 \Omega k=J\cong L\oplus k.
\]
Thus the second syzygy of every nonprojective left module contains
a nonzero simple direct summand.  If
$\operatorname{Tor}_2^A(W,V)\ne0$, both modules are
nonprojective; dimension shifting and
Lemma~\ref{B:local-minimal-Tor} then give nonzero Tor in every
degree at least three.

In the second case $xy=c$ and $x^2=yx=y^2=0$.
Again every first left syzygy is a direct sum of $k$ and the
two-dimensional module $L=Ay=\langle y,c\rangle$, on which
$x$ acts as one nilpotent Jordan block.  This time
\[
 \Omega L=\operatorname{ann}_{\mathrm{left}}(y)
          =\langle y,c\rangle=L,\qquad
 \Omega k=J\cong L\oplus k.
\]
If $\Omega V$ contains $k$, dimension shifting and
Lemma~\ref{B:local-minimal-Tor} give nonzero Tor in every degree
at least two.  Otherwise $\Omega V$ is a nonzero direct sum of
copies of the one-periodic module $L$, so
$\operatorname{Tor}_n^A(W,V)$ is periodic from degree two onward.
Its nonzero degree-two group then recurs in every higher degree.
This proves the theorem.
\end{proof}

\begin{corollary}\label{B:five-outside-minimal-band}
If a finite monoid $M$ has a rectangular-band minimal ideal $K$,
$BM$ has finite-dimensional homotopy type, and
$H_4(BM;\mathbb Q)\ne0$, then $|M\setminus K|\geq5$.
\end{corollary}

\begin{proof}
The quotient $\mathbb QM/\mathbb QK$ is the contracted algebra of
the Rees quotient monoid $M/K$ and has dimension
$|M\setminus K|$.  Apply Theorem~\ref{B:contracted-four-Tor}
to the row and column augmentation modules in
Lemma~\ref{lem:B-band-Tor}.
\end{proof}


\begin{proposition}\label{prop:B-band-wedge}
Let $R_{r,s}^1$ be the monoid obtained by adjoining an identity to the
$r$ by $s$ rectangular band.  Then
\[
 BR_{r,s}^1\simeq\bigvee^{(r-1)(s-1)}S^2.
\]
Consequently the minimum order for $S^2\vee S^2$ is $7$, using
the exhaustive lower-order census in Section~\ref{sec:C}.
\end{proposition}
\begin{proof}
Write $R=\mathbb Z[R_{r,s}^1]$ and let
$P_i=\mathbb Z\{x_{i0},\ldots,x_{i,s-1}\}$, a projective right
module because $P_i=x_{i0}R$.  There is an exact resolution
\[
 0\longrightarrow\bigoplus_{j=1}^{s-1}\bigoplus_{i=0}^{r-1}P_i
 \xrightarrow{d_2}R^{s-1}\xrightarrow{d_1}P_0
 \xrightarrow{\eps}\mathbb Z\longrightarrow0.
\]
Here $d_2$ includes each row module into its indicated copy of $R$,
and $d_1$ sends the $j$th free generator to $x_{00}-x_{0j}$.
These $s-1$ differences are an integral basis of $\ker\eps$.
Each annihilates every band element on the right; hence $\ker d_1$
is exactly the direct sum of the band ideals in the $s-1$ copies of
$R$, which is precisely $\im d_2$.  The first map is injective
on the indicated disjoint bases.  Tensoring gives the block-sum map
$\mathbb Z^{r(s-1)}\to\mathbb Z^{s-1}$, followed by the zero
map to $\mathbb Z$.  Its kernel is free of rank $(r-1)(s-1)$
and it is onto.  Thus the only positive homology is in degree two
with that rank.  Every element other than the identity is idempotent,
so the universal group is trivial.  A Hurewicz basis is represented
by a map from the stated wedge, and the simply connected homology
Whitehead theorem makes that map a homotopy equivalence.
Taking $r=2,s=3$ gives order seven.  The independent census rules
out rank-two $H_2$ at every smaller order, giving the lower bound.
\end{proof}

\begin{proposition}\label{prop:B-S4-nine}
The following nine-element monoid has classifying space $S^4$:
\[
\begin{array}{c|ccccccccc}
 \cdot&1&a&b&c&d&u&v&w&t\\\hline
 1&1&a&b&c&d&u&v&w&t\\
 a&a&a&b&a&b&a&a&a&a\\
 b&b&a&b&a&b&a&a&a&a\\
 c&c&c&d&c&d&c&c&c&c\\
 d&d&c&d&c&d&c&c&c&c\\
 u&u&a&b&a&b&a&a&a&u\\
 v&v&a&b&a&b&a&a&v&a\\
 w&w&a&b&a&b&u&a&w&a\\
 t&t&a&b&c&d&a&a&a&t
\end{array}
\]
Together with the order-eight exclusion in the census, this proves
that the minimum order for $S^4$ is $9$.
\end{proposition}
\begin{proof}
This is the table at index $754$ in the order-nine file of
\cite{SweeneyData2026B}; the unit has been relabelled first.
All $729$ associativity equations are checked independently in the
certificate.  Put $R=\mathbb ZM$.  There is a length-four projective
resolution
\[
\begin{split}
0\longrightarrow Ra\oplus R(t-a)
 &\xrightarrow{d_4}R(1-w)\xrightarrow{d_3}R(w-a)\\
 &\xrightarrow{d_2}R(t-a)\xrightarrow{d_1}Ra
       \xrightarrow{\eps}\mathbb Z\longrightarrow0,
\end{split}
\]
where $d_1,d_2,d_3$ are right multiplication by $c-a,u-a,v-a$,
respectively, and $d_4(s,r)=s(b-a)+r$.  All four ring elements
$a,t-a,w-a,1-w$ are idempotent, so the modules are projective.
Bases for the successive terms starting at $P_0$ are
\[
\begin{gathered}
 (a,c),\qquad(u-a,t-a),\qquad(v-a,w-a),\\
 (1-w,b-a,d-c,u-a,t-a),\qquad
 (a,c)\oplus(u-a,t-a).
\end{gathered}
\]
On these bases, $d_1$ sends only the second basis vector to
$c-a$; $d_2$ sends only the second to the first; $d_3$ sends
only the first to the first.  Finally $d_4$ identifies its four
basis vectors with the last four basis vectors of $P_3$.
This proves exactness integrally at every term, including both ends.
Tensoring with the trivial right module kills every difference
idempotent summand.  The resulting complex has one $\mathbb Z$
in degrees zero and four and zero elsewhere.  Thus the only reduced
homology is $H_4=\mathbb Z$.  The group completion is trivial:
$a,b,c,d$ are idempotents and, for every element $m$, $ma$ is
$a$ or $c$, forcing the image of $m$ to be $1$.
The simply connected homology-sphere argument now gives $BM\simeq S^4$.
The lower bound uses only the exact homology exclusion in the
independently enumerated lower-order tables, not a conjectural
description of any higher-degree tail.
\end{proof}

\begin{proposition}\label{prop:B-Moore-models}
For each $n\ge2$ there is an explicit monoid $M_n$ of order
$n^2+4n+1$ with $BM_n\simeq M(\mathbb Z/n,2)$ and a projective
resolution of the trivial integral module of length three.
In particular the orders $13$ and $22$ are certified upper bounds
for the two torsion Moore spaces in Part~B4.
\end{proposition}
\begin{proof}
We check the construction and resolution of Martinian--Steinberg
\cite[Theorem~4.1]{MartinianSteinberg2025B} directly.
Let $K=\{x,y,0,\ldots,n-1\}\times\mathbb Z/n$ be a rectangular
band and $T=(\mathbb Z/n)\times\{s,t\}$, with
$(k,z)(l,w)=(k+l,w)$.  On $K\amalg T$ set
\[
 (i,j)(k,z)=(i,j+k),\qquad
 (k,z)(i,j)=\begin{cases}
 (x,j),&i=x,\\
 (k-1,j),&i=y, z=s,\\
 (k,j),&i=y, z=t,\\
 (k+i,j),&i\in\mathbb Z/n.
 \end{cases}
\]
Adjoin an identity.  The two actions commute; the left action law
follows by applying it separately to $x,y$ and each cyclic row,
and the right action is addition of cyclic indices.  Compatibility
with a product in $K$ holds because such a product retains only its
left row and right column.  These checks cover all eight patterns of
membership in $K,T$ in a triple, and prove associativity.

Put $R=\mathbb ZM_n$, $e=(x,0)$, $t_0=(0,t)$,
$s_1=(1,s)$, and $v=(y,0)-(x,0)$.  There is a resolution
\[
 0\longrightarrow\mathbb ZK\xrightarrow{\xi}
 Rt_0\oplus\mathbb ZK\xrightarrow{\psi}R
 \xrightarrow{\phi}Re\xrightarrow{\eps}\mathbb Z\longrightarrow0,
\]
where $\phi(a)=av$,
$\psi(a,b)=a(t_0-s_1)+b$, and
\[
 \xi(i,j)=\big((i,j),(i,j+1)-(i,j)\big).
\]
All summands are projective: $\mathbb ZK=\bigoplus_jR(x,j)$.
The image of $\phi$ is the augmentation kernel, since
$(k,t)v=(k,0)-(x,0)$ and $v$ supplies the remaining row difference.
The elements of $K$ annihilate $v$, while
$(k,t)$ and $(k+1,s)$ have the same image under $\phi$.
Together with the distinct image of the identity, these observations
show on the displayed bases that
\[
 \ker\phi=\mathbb ZK+
 \sum_k\mathbb Z\big((k,t)-(k+1,s)\big)=\im\psi.
\]
The map $\psi$ is injective on the $\mathbb ZT$ part of $Rt_0$.
Its remaining kernel consists exactly of the vectors $\xi(i,j)$,
whose first coordinates are an independent basis of $\mathbb ZK$.
Thus $\xi$ is injective and the entire sequence is exact.

Tensoring gives ranks $(1,1,n+1,n)$ with
$d_1=0$, $d_2=(0,1,\ldots,1)$ and columns of $d_3$ equal to
$(1,e_{j+1}-e_j)$.  These columns are independent: a linear
combination with zero cyclic-difference coordinates has all
coefficients equal, and its first coordinate is $n$ times that
coefficient.  In $\ker d_2/\im d_3$, all cyclic row differences
eliminate and the sole remaining generator has relation $n=0$.
Hence the reduced homology is exactly $\mathbb Z/n$ in degree two.
The group completion is trivial because the band idempotents become
$1$ and $(x,0)(k,z)=(x,k)$ then forces every element of $T$ to
become $1$.  This identifies the space as the stated simply connected
Moore space.  Independently generated certificates for $n=2,3$
give the shorter abelian-rank lists $(4,9,10,4)$ and $(5,17,18,5)$,
respectively, with tensored last differential multiplication by
$2$ or $3$.
\end{proof}

\begin{proposition}[Computer-assisted Moore-space lower bound]
\label{prop:B-Moore-minimum}
Using the complete multiplication-table census in
\cite{SweeneyEnumeration2026B}, no monoid of order at most twelve
has classifying space $M(\mathbb Z/2,2)$ or $M(\mathbb Z/3,2)$.
Consequently the minimum order for $M(\mathbb Z/2,2)$ is $13$,
and the minimum for $M(\mathbb Z/3,2)$ is at least $13$.
\end{proposition}
\begin{proof}
The precise external input is the complete set of tables in the
\texttt{results1} release of \texttt{semisearch}, restricted only by
the absence of a proper corner satisfying $eMe=eM$ or $eMe=Me$,
and taken up to isomorphism and opposite.  Every other monoid reduces
to one in this list by a classifying-space equivalence.  The list's
completeness is the published enumeration input; we independently
recompute the homological exclusions, and do not describe the tables
as independently regenerated.  The independent census through eight
provides a separate check of its low-order totals.

At order twelve the file contains $1,974,639,821$ tables.
Five are groups.  The $1,973,253,771$ with a $2\times2$
rectangular-band minimal ideal are all excluded by
Proposition~\ref{prop:B-K22-obstruction}.  The remaining families
have the following sizes:
\[
\begin{array}{c|r}
\text{minimal ideal}&\text{number of tables}\\\hline
2\times3\text{ band}&1,376,598\\
2\times4\text{ band}&8,282\\
3\times3\text{ band}&198\\
2\times5\text{ band}&17\\
2\times2\text{ Rees matrix semigroup over }C_2&950.
\end{array}
\]
For a rectangular band the presentation in
Lemma~\ref{lem:B-band-Tor} has at most four generators.  The program
\texttt{b\_filter\_order12.cpp} computes the ranks of its relation
matrix over $\mathbb F_2$, $\mathbb F_3$, and
$\mathbb F_{1000003}$.  The last rank equals its rational rank:
each relation column has sum of absolute values at most four, so
each maximal minor has absolute value at most $4^4=256$.
A presentation of either target cyclic group must have rational rank
equal to the number of generators and rank one less over its prime
field.  This necessary test retains only $1624$ band tables.

The integral Smith form of each retained presentation is then computed.
Two have the wrong integral $H_2$.  Of the others, $1618$ have
nonzero $H_4(-;\mathbb F_2)$, three have nonzero
$H_4(-;\mathbb F_3)$, and one has nonzero
$H_6(-;\mathbb F_3)$.  These are exact ranks of finite matrices
in the two-sided bar complex of $A$, justified by
Lemma~\ref{lem:B-band-Tor}.  They exclude the desired Moore
homotopy types without extrapolating a vanishing tail.
For the $950$ non-band minimal ideals, ordinary normalized-nerve
homology suffices: $535$ fail the necessary mod-two $H_1,H_2$
test, and the remaining $415$ fail the mod-three test.
All retained tables are also checked against every associativity
equation.

For orders below twelve the same argument discards groups and all
$2\times2$ band cases structurally.  It leaves $51,989$ tables in
total, distributed as $1,8,105,2091,49784$ at orders seven through
eleven.  Repeated unitization pads these to order twelve for the same
checker; this preserves the classifying space because the previous
identity is a corner satisfying both reduction equalities.  The
presentation filter retains $68$ band tables and $92$ group-kernel
tables.  The former all have nonzero $H_4$ at the relevant prime;
the latter all fail the necessary $H_1,H_2$ tests.
The complete output records every surviving input index and its
individual obstruction.  The source archive's SHA-256 is
\texttt{6fe085eb2d084630a8c3327796e31a37}\allowbreak\texttt{eed73eab2314109e9a44b39dfb415f70};
per-file hashes, source URLs and exact programs accompany it.
Finally Proposition~\ref{prop:B-Moore-models} supplies the
thirteen-element upper bound for $\mathbb Z/2$.
\end{proof}

\begin{proposition}[Complete source census with independently checked obstructions]
\label{prop:B-fourdim-lower}
Using the complete corner-irreducible table census of
\cite{SweeneyEnumeration2026B}, no finite monoid of order at most
twelve has classifying space $\CP^2$ or $S^2\times S^2$.
\end{proposition}
\begin{proof}
The corner reduction and the source-completeness dependency are exactly
as in Proposition~\ref{prop:B-Moore-minimum}.  A non-group irreducible
minimal ideal has at least two rows and two columns.  Its order is
$r s |H|$ and it omits the global identity.  For total order at most
twelve the possibilities are a rectangular band with one dimension
two, a $3\times3$ rectangular band, or a $2\times2$ Rees matrix
semigroup over $C_2$.  The first family, including both orientations,
is excluded in every order by Proposition~\ref{prop:B-two-row-products}.
Groups are excluded by their fundamental groups.

All remaining rectangular-band tables in the complete source files
are processed independently by \texttt{b\_K33\_small\_check.py}:
there are ten below order twelve and $198$ at order twelve.
The program checks all associativity equations and computes integral
$H_2$ by the balanced presentation in Lemma~\ref{lem:B-band-Tor}.
It then computes mod-two $H_3$ and $H_4$ in the small two-sided bar
complex over $A$.  No table has the necessary values
$H_2=\mathbb Z^2$, $\dim H_3(-;\mathbb F_2)=0$,
$\dim H_4(-;\mathbb F_2)=1$ for $S^2\times S^2$.
For $\CP^2$ exactly thirteen tables survive the analogous test with
$H_2=\mathbb Z$.  The independent normalized-cochain computation
\texttt{a\_order12\_cups.py} excludes each of them: it computes
the unique nonzero class in $H^2(-;\mathbb F_2)$, produces an
explicit three-cochain whose coboundary is its square, and verifies
that equation on all $11^4=14641$ normalized quadruples for each
table.  The square of the degree-two generator of $\CP^2$ is nonzero.

There are $92$ remaining nontrivial-local-group tables below order
twelve and $950$ at order twelve.  The independent ordinary-nerve
bit-vector calculation \texttt{b\_cp2\_group\_filter.py} checks
all $1042$ tables.  Their mod-two triples $(b_1,b_2,b_3)$, with
the last coordinate computed only when needed, have the following
distribution:
\[
\begin{array}{c|r}
(b_1,b_2,b_3)&\text{number}\\\hline
(1,2,\text{not needed})&470\\
(1,1,\text{not needed})&114\\
(0,0,\text{not needed})&7\\
(0,1,1)&167\\
(0,1,2)&159\\
(0,1,3)&94\\
(0,1,4)&31.
\end{array}
\]
Every table fails a necessary low-degree homology condition for both
target spaces.  All these ranks and all thirteen cup-square primitives
are recomputed from the full multiplication tables; no published
homology value is used in this exclusion.  As before, the completeness
of the source enumeration, rather than the homological calculation,
is the external dependency.
\end{proof}

\begin{proposition}[Further complete restricted enumeration]
\label{prop:B-fourdim-lower14}
No finite monoid of order at most fourteen has classifying space
$\CP^2$ or $S^2\times S^2$.
The enumeration-completeness dependency is the bounded-quotient
source dataset of \cite{SweeneyEnumeration2026B}.
\end{proposition}
\begin{proof}
By Corollary~\ref{cor:B-two-row-general}, both target types require
at least three rows and three columns in the minimal ideal of a
corner-irreducible model.  At orders at most fourteen its maximal
group must be trivial, since $3\cdot3\cdot2=18$; thus the ideal
is a rectangular band of order at least nine.  There are at most
four nonunit elements outside that ideal, since the global identity
is outside it.  Therefore every remaining candidate is included in
the source's enumeration with at most four nonunits outside the
minimal ideal.  Although that dataset is not an unrestricted census
at larger orders, it is complete for precisely the families now needed
through order fourteen.

The source files contain $6191$ tables with $3\times3$ minimal band
and one with $3\times4$ minimal band at order thirteen.  At order
fourteen the corresponding numbers are $214707$ and $20$.
The programs \texttt{b\_extract\_bounded\_largebands.py} and
\texttt{b\_filter\_largebands.cpp} read all these tables.
They use the balanced presentation for $H_2$ and the two-sided bar
complex from Lemma~\ref{lem:B-band-Tor}.  The necessary values for
$\CP^2$ leave respectively $135$ and $3007$ tables after checking
$H_2$ and mod-two $H_3,H_4$.  Every one of these has nonzero
$H_5(-;\mathbb F_2)$, so none represents $\CP^2$.
The separate program \texttt{b\_verify\_largeband\_CP2.py}
recomputes the retained tables with independently implemented
integer Smith forms and finite-field matrices, checking the
individual obstruction and all associativity equations.
For $S^2\times S^2$, none of the $6192$ order-thirteen tables
passes the necessary $H_2,H_3,H_4$ test.  At order fourteen this test
leaves $831$ tables.  All have zero mod-two $H_5$, but $530$ have
nonzero mod-two $H_6$.  The separate tuple-indexed bit-vector program
\texttt{b\_verify\_largeband\_products.py} independently checks
these ranks and every differential composite.  For the remaining
$301$ tables, \texttt{a\_order14\_products.py} independently
computes a basis $x,y$ of $H^2(-;\mathbb F_2)$ and verifies
$x^2=xy=y^2=0$.  Each of the $903$ asserted zero products has an
explicit three-cochain primitive, checked on all $13^4=28561$
normalized quadruples.  This excludes $S^2\times S^2$, whose
cross product is nonzero.
Lower orders have already been excluded in
Proposition~\ref{prop:B-fourdim-lower}.

The exact bounded-quotient source archive has SHA-256
\begin{center}\scriptsize\ttfamily
7aad49cacdd3fee42b834f1700597e0d81602e009f61b27a9b489b86bb5ead1b
\end{center}
Its job definition explicitly sets the number of nonunits outside
the minimal ideal to at most four.  Per-file hashes and the full
counts are included in the independent extraction manifest.
\end{proof}

\begin{proposition}[Complete exclusions through order sixteen]
\label{B:fourdim-lower17}
Neither $\mathbb CP^2$ nor $S^2\times S^2$ is the classifying space
of a finite monoid of order at most sixteen.  Thus their minimum
orders are at least seventeen.  The sole census-completeness
input is the bounded-quotient source of
\cite{SweeneyEnumeration2026B}, in the precise families below.
\end{proposition}
\begin{proof}
Orders at most fourteen were excluded in
Proposition~\ref{prop:B-fourdim-lower14}.  Pass to a
corner-irreducible monoid, which can only decrease the order and
preserves the classifying-space type.  Its minimal ideal has at
least two rows and two columns, and
Corollary~\ref{cor:B-two-row-general} excludes either dimension two.
If the maximal group in that ideal were nontrivial, the ideal alone
would have at least $3\cdot3\cdot2=18$ elements.  Hence at orders
fifteen and sixteen it is a rectangular band.  Theorem~\ref{B:three-by-three-cups}
excludes the three-by-three shape.  Up to taking opposites, the
only remaining shapes are three-by-four and three-by-five.

The three-by-four ideal has twelve elements, so there are at most
three nonunits outside it at these orders; the global identity
accounts for one additional outside element.  The three-by-five
ideal can occur only at order sixteen, with no outside nonunit.
The source enumeration with at most four outside nonunits is
therefore complete for every remaining case.  Its counts are
\[
\begin{array}{c|r|r}
\text{order}&3\times4&3\times5\\\hline
15&936&0\\
16&49631&1.
\end{array}
\]
The independent program
\path{code/b_CP_product_lower17_residuals.py}
reads all $50568$ tables, checks every associativity triple, and
computes their integral $H_2$ by Smith form.  Exactly $35912$
have one of the necessary values $\mathbb Z$ or $\mathbb Z^2$.
For each it separately constructs the tuple-indexed two-sided bar
complex over $\mathbb F_2$, verifies consecutive differential
composites, and computes the necessary groups in degrees three
through six until an obstruction is found.  Every table has a
wrong group in those degrees; none survives.  The certificate
\path{code/outputs/b_CP_product_lower17_residuals.json}
records every retained source index and its particular obstruction.
Thus all remaining shapes are excluded, proving the claim.
\end{proof}

\begin{corollary}[Independent lower bounds of nineteen]
\label{B:fourdim-lower19}
Every finite monoid whose classifying space is $\mathbb CP^2$ or
$S^2\times S^2$ has at least nineteen elements.
This lower bound uses the explicit rational action-hull certificates
and does not depend on an external multiplication-table census.
\end{corollary}
\begin{proof}
Suppose a model had at most eighteen elements, and pass to a
corner-irreducible model of no larger order.  A one-sided minimal
ideal permits a further group corner; its classifying space cannot
be either simply connected target.  Thus both minimal-ideal
dimensions are at least two.  Corollary~\ref{cor:B-two-row-general}
excludes either dimension two.  A nontrivial maximal group would
then give a minimal ideal of order at least $3\cdot3\cdot2=18$;
the global identity is outside that nonunital ideal, so the monoid
would have at least nineteen elements.  Hence the remaining minimal
ideal is a rectangular band.

Theorems~\ref{B:three-by-three-cups} and
\ref{B:three-by-four-targets} exclude the three-by-three and
three-by-four shapes and their opposites.  Among bands with both
dimensions at least three and fewer than eighteen elements, the only
remaining possibilities are three-by-five and four-by-four.
They have fifteen and sixteen elements, respectively, leaving at
most three elements outside the minimal ideal.  Both target spaces
have nonzero rational $H_4$ and finite homotopy dimension, contrary
to Corollary~\ref{B:four-outside-minimal-band}.
This exhausts all cases through order eighteen.
\end{proof}

\begin{corollary}[Independent lower bounds of twenty]
\label{B:fourdim-lower20}
Every finite monoid whose classifying space is $\mathbb CP^2$ or
$S^2\times S^2$ has at least twenty elements.
This bound is independent of an external monoid census.
\end{corollary}
\begin{proof}
Suppose the monoid had at most nineteen elements, and reduce to
a corner-irreducible model as before.  A group corner cannot give
either target, and the two-row/column obstruction excludes a
minimal-ideal dimension two.  Theorem~\ref{B:three-by-three-cups}
excludes the three-by-three shape even when its maximal group is
nontrivial.  Any other shape with a nontrivial maximal group has
at least $3\cdot4\cdot2=24$ elements in the minimal ideal alone,
so cannot occur here.  Hence the remaining minimal ideal is a band.

Theorem~\ref{B:three-by-four-targets} excludes the three-by-four
shape and its opposite.  Every remaining band has at least fifteen
elements, leaving at most four elements outside it.  Both target
spaces have nonzero rational fourth homology and finite homotopy
dimension, contrary to Corollary~\ref{B:five-outside-minimal-band}.
This proves the claim.
\end{proof}

\begin{lemma}[The local five-dimensional target obstruction]
\label{B:local-five-targets}
Let $A$ be a local algebra over a field $k$, with
$\dim_k A=5$ and $A/\operatorname{rad}(A)\cong k$.
Let $W$ be a right $A$-module of dimension four and $V$ a left
$A$-module of dimension two.  The conditions
\[
 \dim_k(W\otimes_A V)\in\{1,2\},\qquad
 \operatorname{Tor}_1^A(W,V)=0,\qquad
 \dim_k\operatorname{Tor}_2^A(W,V)=1,\qquad
 \operatorname{Tor}_3^A(W,V)=0
\]
cannot hold simultaneously.
\end{lemma}

\begin{proof}
Put $J=\operatorname{rad}(A)$ and write $\operatorname{top}U$
for the semisimple quotient of a module $U$.
If $\dim\operatorname{top}V=2$, then $V\cong k^2$ and
$\operatorname{Tor}_2^A(W,V)$ has even dimension.  Thus $V$ is
cyclic and has one-dimensional top.
The natural quotient
\[
 W\otimes_A V\twoheadrightarrow
 (W/WJ)\otimes_k(V/JV)
\]
shows that $\dim\operatorname{top}W\leq2$.
If $W$ is cyclic, its minimal cover $A\to W$ has a
one-dimensional kernel, necessarily $k$.  Dimension shifting and
Lemma~\ref{B:local-minimal-Tor} then give nonzero
$\operatorname{Tor}_n^A(W,V)$ for every $n\geq2$, a
contradiction.  We are therefore in the case
\[
 \dim\operatorname{top}W=2,\qquad \dim WJ=2,
 \qquad\dim(W\otimes_A V)=2.
\]

Let $I=\Omega V$.  It has dimension three and lies in $J$.
The assumption $\operatorname{Tor}_1^A(W,V)=0$ makes
$W\otimes_A I\to W$ injective, with image in $WJ$.
On the other hand, $W\otimes_A I$ surjects onto
$\operatorname{top}W\otimes_k\operatorname{top}I$.
It follows that $2\dim\operatorname{top}I\leq2$, so $I$ is
cyclic.  Its first syzygy $L=\Omega^2V$ consequently has
dimension two.  If $L$ were not cyclic, it would be $k^2$, and
\[
 \operatorname{Tor}_3^A(W,V)
   \cong\operatorname{Tor}_1^A(W,L)
   \cong\operatorname{Tor}_1^A(W,k)^2\ne0,
\]
by Lemma~\ref{B:local-minimal-Tor}.
Thus $L$ is cyclic.  Its syzygy $I'=\Omega^3V$ again has
dimension three.  Since
$\operatorname{Tor}_1^A(W,L)=\operatorname{Tor}_3^A(W,V)=0$,
the same injection into $WJ$ proves that $I'$ is cyclic.

The first four projectives $P_0,P_1,P_2,P_3$ in a minimal
resolution of $V$ are therefore all $A$.
After tensoring with $W$, write $d_i$ for the differential from
degree $i$ to degree $i-1$.  Every differential into a single
copy of $W$ has image in $WJ$, hence rank at most two, because
the original resolution is minimal.
Vanishing of degree-one homology gives
$\operatorname{rank}d_1+\operatorname{rank}d_2=4$, so both
ranks are two.  Degree-two homology of dimension one then gives
$\operatorname{rank}d_3=1$.
But $\ker d_3$ has dimension three, whereas
$\operatorname{im}d_4\subseteq WJ$ has dimension at most two.
Thus degree-three homology is nonzero, the final contradiction.
\end{proof}

\begin{proposition}[The two nonlocal five-dimensional exceptions]
\label{B:five-nonlocal-exceptions}
Let $N$ be a finite monoid with zero, and let
$A=\mathbb QN/\mathbb Q0$ have dimension five and be nonlocal.
Then one of the following holds:
\begin{enumerate}
\item For all finite-dimensional right and left modules $W,V$,
nonzero $\operatorname{Tor}_2^A(W,V)$ implies nonzero Tor in
arbitrarily high degrees.
\item $A$ is the radical-square-zero algebra of the line quiver
$1\longrightarrow2\longrightarrow3$.
\item $A$ is the two-vertex algebra with
$a\in e_1Ae_2$, $b\in e_2Ae_1$, $ba=0$, and $ab=c\ne0$,
where all products of $c$ with the radical are zero.
\end{enumerate}
In the second case the composite of the two arrows is zero.
In the third case the projective lengths are three and two.
The alternatives are statements about rational algebras; they do
not prescribe a particular multiplicative monoid basis.
\end{proposition}

\begin{proof}
Write $J=\operatorname{rad}(A)$ and $S=A/J$.
Semisimple algebras contribute no positive Tor.
We first dispose of nonsplit semisimple quotients.
Any nonzero $S$-bimodule component incident with a field factor
of degree $d>1$ has rational dimension at least $d$.
The dimension bound therefore leaves only a quadratic field
together with a rational factor and a radical of dimension at
most two, or a larger semisimple quotient with a radical of
dimension at most one.
In the former case a two-dimensional radical component is either
a single cross component, giving a hereditary triangular algebra,
or a loop at the quadratic factor, giving a local square-zero
extension of that field.  If the radical is supported on the
rational factor, that factor has dimension at most three.
The local square-zero syzygy argument of
Lemma~\ref{B:small-quotient-algebra} works over any division
residue algebra, and the other cases are hereditary or products
with already treated small algebras.
A radical of dimension one is supported on rational factors and
gives either a loop or one cross arrow, with the same conclusion.
A matrix or noncommutative division factor of rational dimension
four cannot act nontrivially on a one-dimensional radical, so it
introduces no additional case.

We may now suppose $S\cong\mathbb Q^r$, $r\geq2$.
Lift its primitive orthogonal idempotents.  If $r\geq4$, the
radical has dimension at most one, and the previous argument
applies.  Suppose $r=3$, so $\dim J=2$.
If $J^2\ne0$, the algebra is a product of
$\mathbb Q[x]/(x^3)$ with two rational factors.
If $J^2=0$, it is a radical-square-zero quiver algebra with
three vertices and two arrows.  The proof of
Lemma~\ref{B:nonlocal-four-Tor} applies whenever every vertex
with an entering and a leaving arrow lies on a directed cycle.
The only exception is the directed line on three distinct
vertices, which is alternative~(2).

It remains to treat $r=2$, $\dim J=3$.
If $J^2=0$, every vertex with both an entering and a leaving arrow
lies on a loop or a two-cycle, so the semisimple-syzygy proof of
Lemma~\ref{B:nonlocal-four-Tor} again applies, regardless of
parallel arrows.
If $\dim J^2=2$, the sole arrow in $J/J^2$ is a loop and the
algebra is $\mathbb Q[x]/(x^4)\times\mathbb Q$.
Suppose $\dim J^2=1$.  Nilpotence gives $J^3=0$ and the quiver
has two arrows.

Two parallel cross arrows cannot generate $J^2$.
One loop at each vertex gives
$\mathbb Q[x]/(x^3)\times\mathbb Q[y]/(y^2)$, after possibly
interchanging the vertices.
Two opposite cross arrows give alternative~(3): their two
composites lie at different vertices, so exactly one can span
the one-dimensional space $J^2$.
The remaining configurations are two loops at one vertex, or a
loop and a cross arrow.

Consider the loop-and-arrow case first.  Up to taking the opposite
algebra, write $x\in e_1Je_1$ and $a\in e_1Je_2$.
The only possible products spanning $J^2$ are $x^2$ and $xa$;
they lie in different idempotent components, so precisely one
is nonzero.  Let $P_i=e_iA$ and let $S_i$ be the right vertex
simples.
If $x^2=c\ne0$ and $xa=0$, then
\[
 \operatorname{rad}P_1=L\oplus P_2,
 \quad L=\langle x,c\rangle,\qquad
 \Omega L=S_1\oplus P_2,\quad
 \Omega S_1=L\oplus P_2.
\]
Here $P_2=S_2$ is projective.  Every submodule of a direct sum
of copies of $\operatorname{rad}P_1$ and
$\operatorname{rad}P_2=0$ is a direct sum of $L,S_1,S_2$:
on the first vertex this is the Jordan decomposition for a
square-zero operator, and the cross arrow acts trivially.
Thus every first syzygy has this form.  The displayed identities
make the positive Tor groups of $L$ and $S_1$ two-periodic, so
nonzero degree-two Tor of any pair persists.

If $xa=c\ne0$ and $x^2=0$, then
\[
 \operatorname{rad}P_1=L\oplus P_2,
 \quad L=\langle x,c\rangle,\qquad \Omega L=L.
\]
On $L$, the cross arrow maps its one-dimensional first-vertex
space isomorphically onto its second-vertex space.
Every submodule of $L^s\oplus P_2^t$ therefore has the same
injectivity property and is a direct sum of copies of $L$ and
$P_2$; no $S_1$ summand occurs.  Hence all first syzygies have
this form and positive Tor is one-periodic after one shift.
This also proves persistence.  Taking opposites exchanges the
roles of the two tensor factors and covers the other orientation.

Finally suppose both loops lie at one vertex.  Then
$A\cong B\times\mathbb Q$ with $B$ a four-dimensional local
algebra.  We must use the contracted-monoid hypothesis here:
an arbitrary such local algebra need not satisfy Tor persistence.
We claim that $B$ itself has a contracted monoid basis of
dimension four.

If $N$ has an idempotent $e\notin\{0,1\}$, then $e$ is one of
the two central primitive idempotents of $B\times\mathbb Q$,
because a local algebra has no other idempotents.
If $\dim Ae=4$, the nonzero elements of the submonoid $Ne$
are a basis of $Ae\cong B$; adjoining its zero gives the required
five-element monoid.
If $\dim Ae=1$, then $Ne=\{0,e\}$ is an ideal and the Rees
quotient $N/\{0,e\}$ gives $A/Ae\cong B$ as a contracted
monoid algebra.

Suppose there is no such idempotent.  Every nonunit is nilpotent
to zero.  The finite semigroup of nonunits is nilpotent: otherwise
its descending powers stabilize at a nonzero set, and repeated
factorization in that finite set gives a nonzero element $z$
with $z=zu$ for a product $u$ of nonunits; a sufficiently high
power of the nilpotent element $u$ gives a contradiction.
Thus the nonunit span is a nilpotent ideal and the semisimple
quotient is the rational group algebra of the unit group.
Its dimension is two, so that group is $C_2=\{1,g\}$.
The three other nonzero monoid elements form a basis of $J$.
In the local factor $B$, a unit of order two is $1$ or $-1$:
writing it as $\pm1+j$ reduces the assertion to
$j(\pm2+j)=0$ with $\pm2+j$ invertible.
Multiplication by $g$ permutes the radical monoid basis, so its
action there cannot be $-1$.  It therefore fixes all three
radical basis elements on both sides.
The element $(1+g)/2$ is the identity of $B$, and together with
those three radical basis elements has products either zero or
another member of the same set.  These four elements, together
with zero, again form a finite monoid whose contracted algebra
is $B$.
Theorem~\ref{B:contracted-four-Tor} now proves persistence in
this last case and completes the classification.
\end{proof}

\subsubsection{The two exceptional five-dimensional module pairs}

The following calculation is invariant under algebra isomorphism; it does
not assume that an isomorphism preserves the distinguished monoid basis.
All modules in this subsection are finite-dimensional over $\Q$.

\begin{lemma}\label{B:five-exceptional-targets}
Let $A$ be one of the two exceptional algebras in
Proposition~\ref{B:five-nonlocal-exceptions}.  Suppose that a right module $W$
and a left module $V$ satisfy
\[
 \dim V=2,\quad \dim W=4,\quad
 \dim\operatorname{Tor}_1^A(W,V)=0,\quad
 \dim\operatorname{Tor}_2^A(W,V)=1.
\]
For the zero-composite line algebra, the additional condition
$\dim(W\otimes_A V)\leq2$ forces
\[
 V=S_1\oplus S_3,\qquad
 W=T_1\oplus T_3\oplus e_2A,
 \qquad \dim(W\otimes_A V)=2.
\]
Here $S_i$ and $T_i$ denote the left and right simple modules at vertex
$i$.  In particular this algebra cannot give the additive homology of
$\CP^2$ through the band--Tor formula.

For the Nakayama algebra, the conditions force
\[
 V=Ae_1/\Q c,\qquad W=T_2\oplus e_1A,
 \qquad \dim(W\otimes_A V)=1.
\]
The unique, up to scalar, balanced form $W\times V\to\Q$ has matrix
rank one.  Thus this algebra cannot give the additive homology of
$S^2\times S^2$, and its degree-two square vanishes in every rectangular
band realization by Lemma~\ref{B:rank-one-square}.
\end{lemma}

\begin{proof}
For the line algebra, put $a\in e_1Ae_2$, $b\in e_2Ae_3$ and $ab=0$.
A left module is a diagram
$V_3\xrightarrow{b}V_2\xrightarrow{a}V_1$ with zero composite.
Choose a basis of $\operatorname{im}b$, extend it to a basis of
$\ker a$, and then extend to $V_2$; compatible bases in $V_1,V_3$
split the representation into the three simples and the two one-arrow
modules.  The same argument applies on the right.  In particular the
right indecomposables are $T_1,T_2,T_3,e_1A,e_2A$, and their only
nontrivial projective resolutions are
\[
 0\longrightarrow e_3A\xrightarrow{b}e_2A
   \xrightarrow{a}e_1A\longrightarrow T_1\longrightarrow0,
 \qquad
 0\longrightarrow e_3A\xrightarrow{b}e_2A
   \longrightarrow T_2\longrightarrow0.
\]
Consequently $\operatorname{Tor}_2(T_1,V)=\ker(b:V_3\to V_2)$,
$\operatorname{Tor}_1(T_1,V)=\ker a/\operatorname{im}b$, and
$\operatorname{Tor}_1(T_2,V)=\ker b$.

The degree-two Tor condition says that exactly one copy of $T_1$
occurs in $W$ and exactly one copy of $S_3$ occurs in $V$.  Since
$\dim V=2$, the second summand of $V$ must be $S_1$ or $S_2$;
the latter makes $\operatorname{Tor}_1(T_1,V)$ nonzero.  Thus
$V=S_1\oplus S_3$.  The degree-one condition excludes $T_2$ from
$W$, so
\[
 W=T_1\oplus T_3^{\oplus s}\oplus(e_1A)^{\oplus p}
                 \oplus(e_2A)^{\oplus q},\qquad
 s+2p+2q=3.
\]
Its tensor product with $V$ has dimension $1+s+p$.
If this is at most two, parity forces $(s,p,q)=(1,0,1)$.

For the Nakayama algebra, take
$a\in e_1Ae_2$, $b\in e_2Ae_1$, $ba=0$, $ab=c$.
Its left indecomposables are
\[
 S_1, S_2, L=Ae_1/\Q c, Ae_1, Ae_2,
\]
of dimensions $1,1,2,3,2$; the right indecomposables are
$T_1,T_2,R=e_1A/\Q c,e_1A,e_2A$ of the same dimensions.
For completeness, this decomposition is elementary: in a representation
with maps $b:V_1\to V_2$ and $a:V_2\to V_1$, choose a basis of
$\operatorname{im}(ab)$ and lift it successively through $a$ and $b$.
These length-three chains split off copies of $Ae_1$.
On the remaining summand both composites vanish; splitting the two
images inside the opposite kernels gives length-two chains and simples.
The right-handed argument is identical.

The relevant right resolutions are
\[
\begin{split}
 &0\longrightarrow e_2A\xrightarrow{a}e_1A
       \longrightarrow T_1\longrightarrow0,\\
 &0\longrightarrow e_2A\xrightarrow{a}e_1A
       \xrightarrow{b}e_2A\longrightarrow T_2\longrightarrow0,\\
 &0\longrightarrow e_2A\xrightarrow{a}e_1A
       \xrightarrow{c}e_1A\longrightarrow R\longrightarrow0.
\end{split}
\]
Each displayed arrow is specified by the image of the generating
idempotent.  These formulas give the following table; each entry is
$(\dim\operatorname{Tor}_0,\dim\operatorname{Tor}_1,
\dim\operatorname{Tor}_2)$.
\[
\begin{array}{c|ccccc}
 &S_1&S_2&L&Ae_1&Ae_2\\\hline
T_1 &(1,0,0)&(0,1,0)&(1,1,0)&(1,0,0)&(0,0,0)\\
T_2 &(0,1,0)&(1,0,1)&(0,0,1)&(0,0,0)&(1,0,0)\\
R   &(1,1,0)&(0,0,1)&(1,1,1)&(1,0,0)&(1,0,0)\\
e_1A&(1,0,0)&(0,0,0)&(1,0,0)&(2,0,0)&(1,0,0)\\
e_2A&(0,0,0)&(1,0,0)&(1,0,0)&(1,0,0)&(1,0,0)
\end{array}
\]
The five possible dimension-two left modules are
$S_1^{\oplus2},S_1\oplus S_2,S_2^{\oplus2},L,Ae_2$.
The first and last have zero degree-two Tor, the third has even
degree-two Tor dimension, and the second has positive degree-one Tor
whenever its degree-two Tor is nonzero.  Therefore $V=L$.
The table then excludes $T_1,R$ from $W$ and requires exactly one
$T_2$.  The remaining dimension is three, so $W=T_2\oplus e_1A$.

Finally, $T_2\otimes_A L=0$ and $e_1A\otimes_A L=e_1L=\Q$.
In the bases $(e_1,a,c)$ of $e_1A$ and $(\bar e_1,\bar b)$ of $L$,
the balanced pairing is supported only on $(e_1,\bar e_1)$.
It therefore has rank one, as asserted.
\end{proof}

\begin{corollary}[An independent lower bound of twenty-one for
the projective plane]\label{B:CP2-lower21}
Every finite monoid whose classifying space is $\CP^2$ has at
least twenty-one elements.
\end{corollary}
\begin{proof}
Corollary~\ref{B:fourdim-lower20} already excludes orders below
twenty.  At order twenty, apply the same corner and minimal-ideal
reductions.  The two-row/column, three-by-three and three-by-four
obstructions leave a rectangular band with at least fifteen
elements; a nontrivial minimal maximal group would require at
least twenty-four ideal elements.  By
Corollary~\ref{B:five-outside-minimal-band}, at least five elements
lie outside the band.  Thus the only remaining shape, up to
opposite, is $3\times5$, with exactly five outside elements.

Put $A=\Q M/\Q K$.  The band--Tor formula gives modules
$\dim V=2$, $\dim W=4$ with Tor dimensions $1,0,1,0,\ldots$
starting in degree zero.  If $A$ is local, its residue division
algebra has dimension dividing five.  The existence of the
two-dimensional module $V$ rules out residue dimension five,
so Lemma~\ref{B:local-five-targets} applies.  If $A$ is nonlocal,
Proposition~\ref{B:five-nonlocal-exceptions} leaves only its two
exceptional algebras.  Lemma~\ref{B:five-exceptional-targets}
excludes the line algebra by additive homology, and the Nakayama
algebra by its zero degree-two square.  This proves the bound.
\end{proof}

\begin{lemma}[The multiplicative bases of the remaining line algebra]
\label{B:five-line-bases}
Suppose that a finite monoid $M$ has a $3\times5$ rectangular
minimal band $K$, exactly five elements outside $K$, and
$\Q M/\Q K$ is the zero-composite line algebra.  If its band--Tor
modules have the product pair in
Lemma~\ref{B:five-exceptional-targets}, all rational products of
degree-two classes on $BM$ vanish.
\end{lemma}
\begin{proof}
Write $a\in e_1Ae_2$, $b\in e_2Ae_3$ for the two radical arrows.
The stated pair $V=S_1\oplus S_3$,
$W=T_1\oplus T_3\oplus e_2A$ is annihilated by $a$.
We check the possible multiplicative bases of $A$, rather than
assuming that the arrows belong to such a basis.

The independently generated complete order-six monoid census in
Proposition~\ref{prop:C-census} supplies the Rees quotient $M/K$.
The exact checker \texttt{b\_five\_line\_monoid\_bases.py}
reads all $1373$ isomorphism-and-opposite representatives and
both orientations.  There are $2325$ distinct resulting tables,
of which $1174$ have a zero.  Contract the zero and enumerate
all functions from the remaining basis to $\{0,1,-1\}$ which
take the identity to one and preserve products.  For a line
algebra these are exactly its three rational characters.
Their common kernel is its two-dimensional radical.  The checker
tests that its square is zero, and computes its simultaneous
left and right character spaces.  The two spaces must be the
arrows $(1,2)$ and $(2,3)$ on three distinct characters.
Every line algebra necessarily passes these tests, so this
filter omits none.

There are exactly $48$ oriented retained tables, coming from
the following zero-based source indices:
\[
\begin{gathered}
83,85,92,205,212,218,221,223,224,293,294,295,\\
474,488,493,576,581,621,642,644,653,664,665,712.
\end{gathered}
\]
In each, the invisible arrow $a$ is a basis element or, up to
nonzero scalar, a difference $q-p$ of two basis elements.
This assertion is checked in exact rational arithmetic; the full
characters, radical vectors and multiplication tables are stored
in the accompanying JSON certificate.

Let $M\to N$ be the image of the simultaneous row and column
action on $K$.  The band embeds in $N$.
If $a$ is a basis element, its row and column augmentation
operators vanish, so both set maps are constant and its image
belongs to the band.  It is therefore zero in $\Q N/\Q K$.

Suppose $a=q-p$.  Two finite-set maps with the same augmentation
operator are either equal or both constant: the differences
$e_{f(i)}-e_{g(i)}$ are independent of $i$, and a nonzero
constant difference of basis vectors forces both maps constant.
If either endpoint character $\chi_1(p),\chi_3(p)$ is nonzero,
both the row and the column operators are nonzero, since both
modules contain the corresponding endpoint simple.  Thus both
set maps agree and $p,q$ have identical profiles.

There are six remaining oriented tables, whose zero-based
retained-record indices are $11,13,14,16,25,34$.
Their row operators vanish.  Each has an explicit product
witness $pr=qr=t\ne0$ in the contracted table.  In its
five-element nonzero-basis numbering the witnesses are
\[
\begin{array}{c|rrrr}
\text{record}&p&q&r&t\\\hline
11&2&3&1&1\\
13&2&3&1&1\\
14&3&4&1&1\\
16&3&4&1&1\\
25&3&4&1&1\\
34&1&2&3&1
\end{array}
\]
The nonzero quotient equality lifts to a literal equality in
$M$.  Since the row maps of $p,q$ are constant,
$\sigma_p\sigma_r=\sigma_t=\sigma_q\sigma_r$ makes those
constant maps equal.  If the common column augmentation
operator is nonzero, the column maps also agree; if it is zero,
both profiles already belong to the band.  In either case
$q-p$ vanishes in $\Q N/\Q K$.

Consequently $\Q N/\Q K$ is a quotient of $A/(a)$.
The latter is the product of a rational field and the path
algebra of a single arrow.  Every quotient of this algebra is
again a product of fields and at most one single-arrow path
algebra, hence hereditary.  The band--Tor formula gives
$H_4(BN;\Q)=0$.
The map $M\to N$ fixes $K$ and retains the same row and column
actions, so the natural balanced-tensor presentation gives an
isomorphism on $H_2(-;\Q)$.  All degree-two classes on $BM$
therefore pull back from $BN$, and all their products vanish.
\end{proof}

\begin{corollary}\label{B:product-lower21}
Every finite monoid whose classifying space is $S^2\times S^2$
has at least twenty-one elements.
\end{corollary}
\begin{proof}
The order-twenty reduction in Corollary~\ref{B:CP2-lower21}
again leaves only a $3\times5$ band with five outside elements,
up to opposite.  Lemma~\ref{B:local-five-targets} excludes the
local algebra case.  Proposition~\ref{B:five-nonlocal-exceptions}
and Lemma~\ref{B:five-exceptional-targets} leave only the line
algebra and its displayed product pair.  Lemma~\ref{B:five-line-bases}
then contradicts the nonzero cross product of $S^2\times S^2$.
\end{proof}


\begin{proposition}[A stronger lower bound for the odd Moore space]
\label{prop:B-Moore3-lower14}
No finite monoid of order at most fourteen has classifying space
$M(\mathbb Z/3,2)$.  Thus its minimum order belongs to $[15,22]$.
The source-census dependency is the same bounded-quotient dataset
as in Proposition~\ref{prop:B-fourdim-lower14}.
\end{proposition}
\begin{proof}
Orders at most twelve are covered by
Proposition~\ref{prop:B-Moore-minimum}.  Apply the same one-sided
corner reductions at orders thirteen and fourteen.  A group cannot
be a model, and a remaining minimal ideal has at least two rows
and two columns.

If the minimal ideal is a rectangular band with one dimension two,
Proposition~\ref{B:two-row-torsion} excludes odd torsion in $H_2$.
For the remaining bands both dimensions are at least three.
All $6192$ order-thirteen and $214727$ order-fourteen tables are
therefore in the complete relevant part of the bounded dataset.
The exact rank filter
\texttt{b\_filter\_moore3\_largebands.cpp} tests the balanced
presentation over $\mathbb F_2$, $\mathbb F_3$, and
$\mathbb F_{1000003}$.  The same minor bound as before proves
that the last rank is the rational rank.  There are no candidates
at thirteen, and exactly two at fourteen, with source indices
$707414$ and $707439$.  Independent integer Smith forms give
$H_2=\mathbb Z/3$ for both.  However, their two-sided bar complexes
both have $\dim H_3(-;\mathbb F_3)=2$, whereas the required Moore
space has dimension one there by the universal coefficient theorem.

It remains to account for nontrivial minimal maximal groups, without
assuming that finite homotopy type forces them to be trivial.
The possible ideal types are $2\times2$ over $C_2$ or $C_3$,
and $2\times3$ over $C_2$, up to the opposite monoid.
The classwise averaging proof of
Lemma~\ref{lem:B-collapse-minimal-group} works over every field
whose characteristic does not divide the order of the minimal
group.  Thus for the $C_2$ cases, collapsing its coordinate gives
a rectangular-band quotient inducing homology isomorphisms over
both $\mathbb Q$ and $\mathbb F_3$.
For a $2\times2$ band, the integral degree-two group is a quotient
of $\mathbb Z$ by relations in $\{0,\pm1,\pm2\}$, so it cannot
have the necessary rational and mod-three dimensions.
For a $2\times3$ ideal over $C_2$, the collapsed quotient has
order at most eight.  The complete low-order band census, checked
by the same integral balanced presentation, has no three-primary
torsion in $H_2$; its complete $281$ table checks at orders five
through eight are recorded in \texttt{b\_small\_band\_no3.json}.
Corner reductions and the group cases account for all smaller models.
This excludes both $C_2$ families.

Finally, a $2\times2$ ideal over $C_3$ has twelve elements, so
at the two orders in question there is at most one additional
nonidentity element.  These cases are also completely contained
in the bounded dataset: two tables at thirteen and twelve at
fourteen.  The independent normalized-nerve calculation
\texttt{b\_Moore3\_order14\_check.py} checks every associativity
equation and computes integral $H_1,H_2$.  Thirteen tables have
$H_2=\mathbb Z$; the remaining table has $H_1=\mathbb Z/3$
and $H_2=\mathbb Z/2$.  Every case is excluded.
The upper bound twenty-two is supplied by
Proposition~\ref{prop:B-Moore-models}.
\end{proof}
\begin{theorem}[The exact minimum for the odd Moore space]
\label{B:Moore3-minimum17}
The least order of a finite monoid whose classifying space is
$M(\Z/3,2)$ is \textbf{17}. The upper-bound monoid is aperiodic.
The lower bound uses the complete bounded-quotient source census
specified in Proposition~\ref{prop:B-fourdim-lower14}, together
with independent exact calculations and the structural exclusions
proved here.
\end{theorem}
\begin{proof}
\textbf{A compact complete multiplication formula.}
Let $I=\{0,1,2\}$ and $J=\{0,1,2,3\}$. Regard a pair of maps
$(\sigma:I\to I,\tau:J\to J)$ as an element of
$T_I\times T_J^{\mathrm{op}}$, with multiplication
\[
 (\sigma,\tau)(\sigma',\tau')
       =(\sigma\sigma',\tau'\tau).
\]
Take all twelve constant pairs $x_{ij}$, the identity, and four
additional pairs $a,b,c,d$ as follows. A tuple lists the successive
values of a function.
\[
\begin{array}{c|c|c}
 &\sigma&\tau\\\hline
1&(0,1,2)&(0,1,2,3)\\
a&(0,0,1)&(0,0,1,1)\\
b&(0,0,1)&(3,3,2,2)\\
c&(0,1,0)&(0,0,2,2)\\
d&(0,1,1)&(1,1,3,3).
\end{array}
\]
This set is closed. The band products are
$x_{ij}x_{kl}=x_{il}$; mixed products are
$m x_{ij}=x_{\sigma_m(i),j}$ and
$x_{ij}m=x_{i,\tau_m(j)}$. The remaining products are
\[
\begin{array}{c|cccc}
 &a&b&c&d\\\hline
a&x_{00}&x_{03}&x_{00}&x_{01}\\
b&x_{01}&x_{02}&x_{02}&x_{03}\\
c&a&b&c&d\\
d&a&b&c&d.
\end{array}
\]
These formulas are the entire seventeen-element table.
Associativity follows from the transformation-pair embedding and
is also checked on all $17^3=4913$ triples. The twelve constants
form its minimal $3\times4$ band. The elements $c,d$ are
idempotent and $a^2=a^3$, $b^2=b^3$, so every cyclic subgroup is
trivial and the monoid is aperiodic.

\textbf{A full integral projective resolution.}
Write $R=\Z M$ and put
\[
 P=Rx_{00},\qquad Q=R(1-c),\qquad E=R(c-x_{00}).
\]
These are projective: their generators are idempotents of $R$,
using $cx_{00}=x_{00}c=x_{00}$. Their integral ranks are
$3,10,4$, respectively. Define
\[
\begin{split}
 w&=x_{20}-x_{00},\\
 s&=x_{01}-x_{00},\\
 u&=x_{01}+a-d+c-2x_{00},\\
 v&=x_{03}+a-b-x_{00},\\
 z&=x_{00}-x_{02}.
\end{split}
\]
There is an augmented sequence
\begin{equation}\label{B:Moore17-resolution}
0\longrightarrow P\xrightarrow{d_3}P\oplus E\oplus E
 \xrightarrow{d_2}Q\xrightarrow{d_1}P
 \xrightarrow{\varepsilon}\Z\longrightarrow0,
\end{equation}
where
\[
 d_1(q)=qw,\qquad
 d_2(p,e_1,e_2)=ps+e_1u+e_2v,\qquad
 d_3(p)=(3p,pz,pz).
\]
The table verifies the required idempotent supports, so these are
$R$-linear maps with the displayed domains and targets.

Here is a direct proof of exactness. A basis of $P$ is
$r_i=x_{i0}$, $0\leq i\leq2$. A basis of $E$ is
$z_i'=x_{i2}-x_{i0}$ for $0\leq i\leq2$, together with
$t=c-x_{00}$. A basis of $Q$ consists of
\[
\begin{gathered}
 q_0=1-c,\qquad
 h_i=x_{i1}-x_{i0},\quad k_i=x_{i3}-x_{i2}\quad(0\leq i\leq2),\\
 A=a-x_{00},\qquad B=b-x_{02},\qquad D=d-c.
\end{gathered}
\]
The map $d_1$ sends $q_0$ to $r_2-r_0$, sends $A,B,D$
to $r_1-r_0$, and kills all $h_i,k_i$. Its image is
$\ker\varepsilon$, and its kernel has basis
$h_i,k_i,A-D,B-D$.
The columns of $d_2$, in the three displayed summands, are
\[
\begin{array}{c|c}
\text{source basis}&\text{image}\\\hline
r_i&h_i\\
z_i'\text{ in the first }E&2h_i-k_i\\
t\text{ in the first }E&h_0+A-D\\
z_i'\text{ in the second }E&h_i+k_i\\
t\text{ in the second }E&k_0+A-B.
\end{array}
\]
They generate the stated kernel integrally. A vector in
$\ker d_2$ has zero coefficient at both copies of $t$, by its
$B,D$ coordinates. The remaining equations say that its first
coordinate is $-3\sum_i n_i r_i$ and both $E$ coordinates are
$\sum_i n_i z_i'$. These are exactly the image of $d_3$, with
an overall sign. The map $d_3$ is injective. This proves
\eqref{B:Moore17-resolution} in every degree, without a truncated
homology inference.

Tensoring with the trivial right module sends $P$ to $\Z$ and
$Q,E$ to zero. The resulting complex has $\Z$ in degrees
$0,2,3$, with differential $\Z\xrightarrow{3}\Z$ from degree
three to degree two. Hence
\[
 H_0(BM;\Z)=\Z,\qquad H_2(BM;\Z)=\Z/3,
 \qquad \widetilde H_i(BM;\Z)=0\quad(i\ne2).
\]
In the universal group every band element is the identity, and
$m x_{00}$ is a band element for every $m$; thus the universal
group is trivial. By Hurewicz, a generator of $\pi_2(BM)$ has
order three. Choose a representative sphere and a nullhomotopy
of its triple. These give a map $M(\Z/3,2)\to BM$ inducing an
isomorphism on $H_2$ and therefore on all homology. Both spaces
are simply connected, so homological Whitehead gives the asserted
homotopy equivalence.

The exact certificate
\path{code/b_Moore3_17_certificate.py} reconstructs the table from
the function pairs and checks every symbolic differential column
against the independently verified resolution. The standalone
verifier checks integral saturation and equivariance, obtaining
projective ranks $(3,10,11,3)$ and differential ranks $(1,2,8,3)$.
The root branch independently reconstructs and verifies the same
table and resolution in \path{code/c_Moore3_17_independent.py}.

\textbf{Excluding all smaller orders.}
Orders through fourteen are covered by
Proposition~\ref{prop:B-Moore3-lower14}. For orders fifteen and
sixteen apply the same one-sided corner reductions. A remaining
non-group minimal ideal has at least two rows and two columns,
and omits the global identity.

For a rectangular-band minimal ideal, one dimension two is
excluded by Proposition~\ref{B:two-row-torsion}; the $3\times3$
case is excluded in every order by
Theorem~\ref{B:three-by-three-no-three}. Every other band has
at least twelve elements, leaving at most four elements outside
it. Thus all these cases occur in the complete relevant portion
of the bounded-quotient census. At fifteen there are $936$
$3\times4$ tables; at sixteen there are $49631$ $3\times4$
tables and one $3\times5$ table, up to opposite. The independent
program \path{code/b_Moore3_largeband_15_16_complete.py} checks
all their associativity triples and computes the exact integral
Smith form of each balanced $H_2$ presentation. None of these
$50568$ groups is $\Z/3$.

The possible nontrivial minimal groups are $C_2$ in a
$2\times2$ or $2\times3$ ideal, and $C_3$ in a $2\times2$
ideal. Larger types would already have at least sixteen elements
before adjoining the global identity. For $2\times2$ over $C_2$,
the classwise-averaging quotient is a homology equivalence over
$\Q$ and $\F_3$ by
Lemma~\ref{lem:B-collapse-minimal-group}. Its minimal band is
$2\times2$, whose balanced $H_2$ is a quotient of $\Z$ by
relations in $\{0,\pm1,\pm2\}$; the required rational and
mod-three dimensions are impossible.

For $2\times3$ over $C_2$, the ideal has twelve elements, so
the bounded dataset again covers every remaining case. There are
$408$ tables at fifteen and $10085$ at sixteen. The program
\path{code/b_Moore3_local_C2_15_16.py} constructs and verifies
each actual minimal-group quotient, then applies the balanced
complex over $\F_3$. Of the $10493$ inputs, $10492$ have the
wrong rational or mod-three $H_2$ dimensions. The remaining
quotient has nonzero $H_4(-;\F_3)$. Averaging transfers this
obstruction to the original monoid. This treats the nontrivial
minimal group directly and does not assume that finite homotopy
type makes it trivial.

Finally, for $2\times2$ over $C_3$, the complete bounded files
contain $95$ and $1261$ tables. The independently implemented
normalized nerve calculation in
\path{code/b_Moore3_local_C3_15_16.py} checks all $1356$.
Exactly $1230$ have nonzero mod-two $H_2$; the remaining $126$
have nonzero mod-three $H_1$. Both contradict the Moore-space
homology. The datasets and programs record every individual
input and exclusion. All cases through sixteen are now covered,
and the seventeen-element construction proves sharpness.
\end{proof}


\begin{proposition}[A smaller explicit quotient]\label{prop:B-CP110}
There is a $110$-element finite aperiodic monoid whose classifying
space is $\CP^2$, with a complete integral projective resolution
of length four.
\end{proposition}
\begin{proof}
Use the $122$-element monoid of Proposition~\ref{prop:B-CP122}.
Write its old elements as $(l,\theta)\in L_0\times O_3$ and
its ideal elements as $(s,j,i)\in A_2\times\{0,1\}\times\{0,1,2\}$.
For every $k,j\in\{0,1\}$ and $i\in\{0,1,2\}$ impose the
twelve identifications
\[
 (x_{kj},\operatorname{const}_i)
       =((k,k,k),j,i).
\]
These are exactly the equalities of the simultaneous left and right
translation actions on the $72$-element minimal ideal.  Indeed each
displayed old element sends every ideal row to the constant tuple
$(k,k,k)$ and every ideal column to $(j,i)$, just as does the
displayed ideal element.  Conversely the six column images distinguish
the ordinal maps and column labels of all other old elements, and
their actions on the twelve row tuples distinguish the remaining
possibilities.  The independent certificate also compares all $122^2$
pairs of translation profiles and finds precisely these twelve
identifications.  Equality of translation profiles is a monoid
congruence, so the quotient $Q$ is associative and has order
$122-12=110$.  Its full table and quotient map are supplied in
\texttt{b\_CP110\_certificate.json}.

The quotient is bijective on the rectangular minimal ideal and induces
exactly the same set of its row and column actions.  The natural
balanced presentation of Lemma~\ref{lem:B-band-Tor} therefore makes
\[
 H_2(BM_{122};\mathbb Z)\longrightarrow H_2(BQ;\mathbb Z)
\]
an isomorphism.  The complete independently verified projective
resolution has the certificate
\begin{center}\small\ttfamily b\_resolution\_CP122\_quotient\_110.json\end{center}
and abelian ranks
\[
 (12,100,164,87,12)
\]
and length four.  After tensoring with the trivial module it has
one generator in degrees zero, two and four and no others.
Thus $H_2(BQ)=H_4(BQ)=\mathbb Z$ and all other positive integral
homology groups vanish.  Aperiodicity descends to the quotient, so
both source and target classifying spaces are simply connected.

Let $u\in H^2(BQ;\mathbb Z)$ be a generator.  Its pullback is
a generator for the already verified $\CP^2$ model, whose square
generates degree four.  Cup naturality implies that $u^2$ pulls back
to a generator.  Since both degree-four groups are infinite cyclic,
the pullback in degree four must be an isomorphism.  The quotient map
is consequently an integral homology equivalence between simply
connected CW complexes, and hence a homotopy equivalence.
The certificate independently checks all $14884$ homomorphism
equations, all $1331000$ target associativity equations, the minimal
ideal comparison, and the exact projective resolution.
\end{proof}

\begin{proposition}[A hundred-element model]\label{prop:B-CP100}
There is a $100$-element finite aperiodic monoid of type $\CP^2$,
with a complete projective resolution and an explicitly verified
homology-equivalence inclusion into the preceding $110$-element model.
\end{proposition}
\begin{proof}
Inside the ordinal monoid $O_3$, retain the identity, the three
constant maps, and the four maps
\[
 (0,0,1),\quad(0,0,2),\quad(0,1,1),\quad(1,1,2).
\]
These eight maps form a submonoid $O'$, as is checked by composing
each ordered pair.  In the $122$-element construction retain
$L_0\times O'$ and the complete $72$-element ideal, and then apply
the twelve identifications of Proposition~\ref{prop:B-CP110}.
This is a submonoid $S$ of the $110$-element quotient, of order
\[
 72+5\cdot8-4\cdot3=100.
\]
The multiplication is therefore completely specified by the preceding
formulas; its full normalized table and its inclusion are also
recorded in the certificates.

The exact projective resolution in
\texttt{b\_resolution\_CP\_sub\_100\_01234679.json}
has length four and abelian ranks
\[
 (12,144,216,95,12).
\]
The independent verifier proves its integral exactness and all module
identities.  Its trivial-coefficient complex has a single generator
in degrees zero, two and four and zero in the other degrees.
To determine the actual map of these groups, the program
\texttt{b\_resolution\_map.py} constructs an augmentation-preserving
chain map into the verified projective resolution of the $110$-element
monoid.  It uses integer Hermite transformations to lift each
source summand's idempotent generator, applies that idempotent to
the lift, and extends by the source monoid action.  Every chain
square and every equivariance equation for all $100$ source elements
is checked.  Thus this is a comparison map for the specified
inclusion, not an abstract isomorphism between two homology groups.
After tensoring, its maps in degrees zero, two and four are
respectively
\[
 (1),\qquad(-1),\qquad(-1).
\]
The inclusion is consequently an integral homology equivalence.
Both monoids are aperiodic, so both classifying spaces are simply
connected, and the homological Whitehead theorem proves the claimed
homotopy equivalence to $\CP^2$.
\end{proof}

\begin{proposition}[An eighty-nine-element model]\label{prop:B-CP89}
There is an $89$-element finite aperiodic monoid with classifying
space $\CP^2$.  Its complete projective resolution has length four,
and its inclusion into the $110$-element model is an explicitly
verified integral homology equivalence.
\end{proposition}
\begin{proof}
In the annulus notation of Proposition~\ref{prop:B-CP122}, let
$O''\subset O_3$ consist of the seven maps
\[
 (0,1,2),\ (0,0,0),\ (0,0,1),\ (0,0,2),\
 (1,1,1),\ (1,1,2),\ (2,2,2),
\]
and let $A'_2=A_2\setminus\{(1,2,2)\}$.  Retain the old
elements $L_0\times O''$ and the ideal
$A'_2\times\{0,1\}\times\{0,1,2\}$, then impose the same
twelve translation-profile identifications as before.
The displayed ordinal maps are closed under composition, and the
specified old actions preserve $A'_2$; these assertions are also
checked directly in the full normalized multiplication table.
Thus this defines a submonoid of the $110$-element model, of order
\[
 11\cdot6+5\cdot7-12=89.
\]

The resolution certificate named \texttt{CP\_sub\_89\_second}
has abelian ranks
\[
 (11,114,204,111,11).
\]
Its trivial-coefficient complex has one generator in degrees zero,
two and four, and none in the other degrees.  The exact comparison
\texttt{CP\_sub\_89\_second\_into\_CP110}, constructed for the
specified inclusion, induces multiplication by $1$ in all three
nonzero degrees.  The standalone verifier reconstructs the actions
from the two multiplication tables, checks all $89$ source actions
in every degree and all chain squares, and recomputes these three
integers.  The separate resolution verifier checks all $89^3$
associativity equations, projectivity, integral exactness and the
zero final kernel.  The inclusion is an integral homology
equivalence between simply connected CW complexes, hence a homotopy
equivalence.

The comparison is necessary: deleting the annulus row $(1,2,3)$
instead gives a companion $89$-element monoid with the same complete
additive homology, but its inclusion induces zero on $H_4$.
Naturality of the square then excludes that companion as a
$\CP^2$ model.  No inference from additive homology alone is used.
\end{proof}

\begin{proposition}[A seventy-seven-element model]\label{prop:B-CP77}
There is a $77$-element finite aperiodic monoid of type $\CP^2$,
with a complete integral projective resolution of length four.
\end{proposition}
\begin{proof}
Continue with the explicit annulus coordinates and let $O''$ be the
seven-element ordinal submonoid in Proposition~\ref{prop:B-CP89}.
Put
\[
 D=\{(0,0,0),(0,0,1),(1,1,1),(2,2,2)\},\qquad
 A''_2=A_2\setminus\{(1,2,2),(1,3,3)\}.
\]
In the $122$-element monoid retain the old elements
\[
 \{1,x_{00},x_{01}\}\times O''
 \quad\cup\quad
 \{x_{10},x_{11}\}\times D
\]
and the ideal
$A''_2\times\{0,1\}\times\{0,1,2\}$.
Apply the twelve old/ideal identifications of
Proposition~\ref{prop:B-CP110}.  The resulting set is closed under
the specified multiplication and contains the identity.  Its order is
\[
 3\cdot7+2\cdot4+10\cdot6-12=77.
\]
The separate reconstruction program
\texttt{b\_CP77\_certificate.py} builds exactly this set from the
annulus labels, checks its complete multiplication table, all
$456533$ associativity equations, all $5929$ equations for its
inclusion into the $110$-element monoid, and stabilization of powers
for all $77$ elements.  Its minimal ideal is the displayed
$10\times6$ rectangular band.

The certificate \texttt{CP\_sub\_77\_first} gives an integral
projective resolution with abelian ranks
\[
 (10,96,176,99,10).
\]
After tensoring with the trivial module, only one copy of $\mathbb Z$
in each of degrees zero, two and four remains.  The exact comparison
\texttt{CP\_sub\_77\_first\_into\_CP110} is induced by the
specified inclusion and has tensor maps $(1),(1),(1)$ in those
degrees.  The independent resolution and comparison verifiers
reconstruct all actions and check integral exactness, projectivity,
every source-equivariance equation and every chain square.
Thus the inclusion is an integral homology equivalence between
simply connected CW complexes, and is a homotopy equivalence.
\end{proof}

\subsubsection{A twenty-three-element projective-plane monoid}

\begin{theorem}\label{B:CP23}
There is a finite aperiodic monoid of order twenty-three whose
classifying space is $\CP^2$.
\end{theorem}
\begin{proof}
Let $I=\{0,1,2\}$ and $J=\{0,1,2,3,4\}$.
Represent a simultaneous action on the rectangular band $I\times J$
by a pair $(\sigma,\tau)$ of transformations, with multiplication
\[
 (\sigma,\tau)(\sigma',\tau')
       =(\sigma\circ\sigma',\tau'\circ\tau).
\]
Take the fifteen constant pairs $x_{ij}=({\rm const}_i,
{\rm const}_j)$, the identity, and the following seven pairs.
The strings specify values on the naturally ordered sets.
\[
\begin{array}{c|cc}
 &\sigma&\tau\\\hline
a&002&02200\\ b&002&41144\\ c&010&04034\\
d&202&03300\\ e&202&12211\\ f&202&40044\\ g&212&00334
\end{array}
\]
This set is closed under the displayed associative multiplication.
Its minimal ideal is the $3\times5$ band.
The elements $a,b,c,g$ are idempotent; $d^2=x_{20}$,
$e^2=x_{22}$, $f^2=x_{24}$, and subsequent powers are unchanged.
Thus the monoid is aperiodic.  The certificate numbering is identity
$0$, $x_{ij}=1+5i+j$, and $a,b,c,d,e,f,g=16,\ldots,22$.

Put $R=\Z M$ and define projective left modules
\[
 P=Rx_{00},\quad E=R(c-x_{00}),\quad
 F=R(1-c),\quad G=R(a-x_{00}).
\]
The four generators are idempotents of $R$.  An exact resolution is
\[
0\longrightarrow P\xrightarrow{d_4}G\oplus G
 \xrightarrow{d_3}P\oplus F\xrightarrow{d_2}E
 \xrightarrow{d_1}P\xrightarrow{\varepsilon}\Z
 \longrightarrow0,
\]
with maps on generating idempotents
\[
\begin{aligned}
d_1(c-x_{00})&=x_{10}-x_{00},\\
d_2(x_{00},0)&=x_{03}-x_{00},\\
d_2(0,1-c)&=c+d-g-x_{00},\\
d_3(a-x_{00},0)&=(x_{00}-x_{20},\ b+e-f-x_{01}),\\
d_3(0,a-x_{00})&=(0,a-x_{00}),\\
d_4(x_{00})&=(x_{02}-x_{00},\ -x_{02}+x_{00}).
\end{aligned}
\]
Each image is fixed by its source idempotent.  For a direct integer
exactness check use the following bases:
\[
\begin{array}{c|l}
P&x_{00},x_{10},x_{20}\\
E&z-zx_{00},\quad
 z\in\{x_{03},x_{04},x_{13},x_{14},x_{23},x_{24},c,d,f,g\}\\
F&z-zc,\quad
 z\in\{1,x_{01},x_{02},x_{11},x_{12},x_{21},x_{22},a,b,e\}\\
G&z-zx_{00},\quad z\in\{x_{02},x_{12},x_{22},a\}.
\end{array}
\]
Their ranks in degrees zero through four are $(3,10,13,8,3)$.
The ranks of $\varepsilon,d_1,d_2,d_3,d_4$ are $(1,2,8,5,3)$;
the composites vanish and every nonzero Smith invariant is one.
These exact integer reductions prove saturation and exactness.
The full matrices are also checked by the independent program
\texttt{b\_verify\_resolutions.py}.
Tensoring with the trivial right module kills $E,F,G$ and sends
$P$ to $\Z$, giving
\[
 H_0(BM;\Z)=H_2(BM;\Z)=H_4(BM;\Z)=\Z
\]
and no other nonzero homology groups.

To determine the square, use the balanced augmentation pairing
\[
 L=\begin{pmatrix}1&1\\0&1\\-1&0\\1&0\end{pmatrix},
\]
with column-set indices first, and extend it by a zero row and
column indexed by $0\in J$ and $0\in I$.
The function $u(m,n)=L_{\tau_m(0),\sigma_n(0)}$ is an integral
normalized two-cocycle: the balance identity for each displayed
actor makes its affine coboundary zero.

The exact bar comparison can be specified recursively.
In the normalized left bar resolution set
$s(z[m_1|\cdots|m_n])=[z|m_1|\cdots|m_n]$ when $z\ne1$,
and zero when $z=1$, extending linearly.
Start with $F_0(x_{00})=x_{00}[\,]$ and, for each generating
projective idempotent $\epsilon$ in positive degree, set
\[
 F_n(\epsilon)=\epsilon\,s\bigl(F_{n-1}d_n(\epsilon)\bigr).
\]
Extend $R$-linearly.  Since $\epsilon$ fixes its boundary,
these formulas give a chain map.
The resulting sparse images contain $3,40,172,448,324$ terms
over the complete abelian bases in degrees zero through four.
Direct integer evaluation gives
\[
 \langle u,F_2(x_{00},0)\rangle=-1,
 \qquad
 \langle u\smile u,F_4(x_{00})\rangle=1.
\]
Every equivariance and chain-map equation is checked in the
comparison certificate associated with
\texttt{b\_orbit164\_sub\_1412\_cup.json}.
Thus $[u]$ and $[u]^2$ generate $H^2$ and $H^4$ respectively.

Finally $BM$ is simply connected.  Every band idempotent becomes
the identity in the universal group, and $mx_{00}\in K$ then
kills every $m\in M$.
Choose a Hurewicz generator $S^2\to BM$.  Its mapping-cylinder
pair is three-connected: it is two-connected by the absolute
Hurewicz theorem, and its zero third relative homology kills the
third relative homotopy group.  Relative Hurewicz now supplies a
single four-cell giving a homology equivalence
$S^2\cup_h e^4\to BM$, hence a homotopy equivalence by
homological Whitehead.  The square calculation gives Hopf
invariant $\pm1$.  Since $\pi_3(S^2)=\Z\eta$ and the Hopf
invariant sends $\eta$ to one, $h=\pm\eta$ and the space is
$\CP^2$.
\end{proof}

\subsubsection{The complete three-by-five action classification}

\begin{theorem}\label{B:three-by-five-targets}
A finite monoid of type $S^2\times S^2$ cannot have a minimal
ideal with three rows and five columns, or the opposite shape,
regardless of its minimal maximal group.
Among finite monoids of type $\CP^2$ with that shape, the minimum
order is twenty-three, attained by Theorem~\ref{B:CP23}.
\end{theorem}

\begin{proof}
We give the complete finite classification and the exact certificates
used to evaluate it.  Collapse the minimal group coordinate rationally
as before.  The row and column augmentation dimensions are two and four,
and a degree-two class is an eight-coordinate balanced form.  There are
$3^3 5^5=84375$ possible transformation pairs.
Their individual relation spaces and all proper sums have the following
exact counts:
\[
\begin{array}{c|rrrrrr}
\text{individual relation rank}&0&2&4&5&6&7\\\hline
\text{number}&1&10&1813&3900&1870&165
\end{array}
\]
There are $68719$ proper sum spaces, in $214$ orbits under
$S_3\times S_5$.  This is a complete relation-space lattice,
not a bounded-coordinate sample.

The enumeration closes under sums one orbit at a time: add every
individual relation space to a representative of each found orbit,
and mark all $720$ transforms of every new space.  Invariance of the
individual relations proves that this suffices for closure everywhere.
Every representative is generated by the relations of its entire
transformation hull, so no unattained extra space is used.
An independent verifier uses double differences of balanced forms,
rather than the constructor's tensor action on relation vectors.
It checks every actor and all $1660426$ representative joins.
For each proper join, a finite-field rank supplies a lower bound and
exact integer annihilator containment supplies the equal upper bound.
Thus the modular calculation is not substituted for rational equality.

Only $164$ orbits have degree-two dimension one or two.  They split as
\[
\begin{array}{c|rrr}
\dim H^2&\text{all forms rank at most one}&
 \text{common repeated columns}&\text{new hulls}\\\hline
1&8&37&26\\
2&6&47&40
\end{array}
\]
The first column is excluded by Lemma~\ref{B:rank-one-square}
and polarization.  In the second, quotient the column set by equality
under every balanced form.  Balance makes this equivalence relation
invariant under every actor, giving a monoid map to a three-by-at-most-four
action monoid.  The column augmentation map is surjective, and all the
original balanced forms descend, so the induced rational $H_2$ map is
an isomorphism.

For clarity, the earlier three-by-three and three-by-four arguments
apply in this relative situation even if the quotient's classifying
space is not finite.  Their zero fourth-homology witnesses and explicit
cup primitives pull back directly.  Each forbidden local character
lifts along the original finite monoid's eventual local group, where
Miller's theorem is applied to the original finite-dimensional target.
In the sole even-square three-by-four case, a band source with
$H_2=\Z$ maps surjectively on integral $H_2$ to the known cyclic
target, so a primitive class pulls back to a primitive class and its
even square still excludes $\CP^2$.
This accounts for every repeated-column case without a finiteness
assumption on an intermediate quotient.

Of the remaining $66$ hulls, $58$ have $H_4(-;\mathbb F_5)=0$.
The other eight, in the fixed zero-based orbit numbering of the
certificate, are as follows:
\[
\begin{array}{c|rrrl}
\text{orbit}&\text{order}&\dim H^2_{\Q}&
 \text{forbidden actors}&\text{remaining calculation}\\\hline
139&28&2&4&\text{allowed hull of order24 has }H_4(\mathbb F_5)=0\\
142&26&2&4&\text{allowed hull of order22 has }H_4(\mathbb F_5)=0\\
164&34&1&4&\text{allowed hull of order30; small images below}\\
194&22&2&0&\text{all three products are integral coboundaries}\\
197&26&2&2&\text{allowed hull of order24 has }H_4(\mathbb F_5)=0\\
198&30&1&4&\text{allowed hull of order26 has }H_4(\mathbb F_5)=0\\
200&30&1&4&\text{allowed hull of order26 has }H_4(\mathbb F_5)=0\\
210&22&2&0&\text{all three products are integral coboundaries}
\end{array}
\]
The forbidden actors have a nonzero restriction of one of the displayed
integral degree-two forms to their eventual cyclic groups.  Their
preimages in a finite source have eventual local groups surjecting onto
those cyclic groups, so the nonzero character remains nonzero on the
source local group.  Miller excludes these actors.  Every allowed
image lies in one of the stated maximal allowed hulls; these maxima
are certified by enumerating all subsets of the remaining actors.
In orbits $194$ and $210$, all six integral three-cochain primitives
are checked on all $21^4$ normalized quadruples each, a total of
$1166886$ exact identities.  Their two forms are independently
certified to span rational $H^2$ by modular rank bounds and integral
grid-cycle pairings.

This already excludes $S^2\times S^2$: every orbit of degree-two
dimension two has zero mixed products on a finite source.
All uses of finite-group restrictions were made on the original
monoid, so the rational minimal-group collapse introduces no gap.

For $\CP^2$, only orbit $164$ remains.  Its four forbidden actors
leave a unique allowed maximal hull of order thirty.  There are exactly
$2186$ closed actor subsets containing its full band and identity.
Among their action monoids of order at most twenty-two, exactly $846$
have integral $H_2=\Z$.  Of these, $842$ have zero
$H_4(-;\mathbb F_5)$.  The other four all have order twenty-two
and Betti numbers
\[
 (\dim H_2,\dim H_3,\dim H_4)(-;\mathbb F_5)=(1,1,2).
\]
No third-homology condition was imposed when enumerating images.
This matters for nonfaithful sources: the profile quotient fixes the
band and its actions, hence preserves integral $H_2$, but need not
preserve higher homology.  If its image has at most twenty-one
elements, its zero rational fourth homology makes the source square
zero.  If its image is one of the four exceptional twenty-two-element
monoids, a source of order at most twenty-two maps bijectively to that
image and is isomorphic to it; its nonzero third homology excludes the
target.  Thus every source of order at most twenty-two is excluded.
A nontrivial minimal group already makes this ideal have at least
thirty elements, so it cannot affect the lower bound.

The independent verification includes $905$ rank certificates for
the $58+5$ ambient hulls and the $842$ small images.  Each certificate
lists independent columns in the two relevant bar differentials;
fresh dense elimination proves their ranks add to the full middle
dimension.  Exact Smith forms are recomputed for every small image,
and the four exceptional Betti profiles are recomputed by separate
dense boundary matrices.  The files
\path{code/outputs/b_rational_orbit_lattice_3x5_verified.json} and
\path{code/outputs/b_rational_3x5_cups_verified.json}
record the independent checks and their complete scopes.
The upper bound is the explicit monoid of Theorem~\ref{B:CP23}.
\end{proof}

\begin{proposition}[Three-by-five obstruction with explicit source hypotheses]
\label{B:three-by-five-finite-target-transfer}
Let $M$ be a finite monoid of order at most twenty-two with a
three-by-five rectangular minimal band, or the opposite shape.
Suppose that
\[
 H_*(BM;\Q)\cong H_*(\CP^2;\Q)
\]
as graded vector spaces, and that every integral degree-two class on
$BM$ restricts to zero along every map $BG\to BM$ from a finite group.
Then the square of a generator of $H^2(BM;\Z)\cong\Z$ is not primitive
in $H^4(BM;\Z)/\operatorname{torsion}$.
In particular the three-by-five lower-bound obstruction applies under
these hypotheses without assuming that $BM$ has finite CW type.
\end{proposition}

\begin{proof}
The minimal band makes the universal group trivial, so $H_1(BM;\Z)=0$.
The integral universal coefficient theorem and the rational homology
hypothesis therefore give $H^2(BM;\Z)=\Z$.  We allow torsion in
$H_2(BM;\Z)$ throughout the argument.

Apply the complete relation-space classification in
Theorem~\ref{B:three-by-five-targets}.  Its hull computations have three
types of output: explicit zero cup products, vanishing rational fourth
homology, and nonzero restrictions of actual integral degree-two
cocycles to eventual local cyclic groups.  The first two types pull
back along the actual monoid maps and require no finiteness assumption
on the source.  For the third type, an eventual local group in the
original source surjects onto the detected cyclic group.  Inflation of
its nonzero degree-two character is nonzero.  The assumed restriction
vanishing excludes exactly the same actors that Miller's theorem
excluded in the finite-CW proof.  Thus no finite-CW assumption on an
intermediate image or quotient is used.

The common repeated-column cases reduce to the preceding
three-by-three or three-by-four hull classifications.  Their zero
products, zero fourth homology and local-character tests apply for the
same reasons.  In the exceptional three-by-four parity case,
Proposition~\ref{B:even-hopf23} supplies a hull whose primitive
degree-two class has square twice an integral degree-four class.
The map to this hull is surjective on integral $H_2$: the row and
column augmentation quotients are surjective, and the balanced tensor
presentation adds relations.  Since the source has free $H_2$ rank
one and the target has $H_2=\Z$, the induced map on free parts is an
isomorphism.  Thus the target's primitive degree-two class pulls back
to a primitive source class, up to sign.  Its square is divisible by
two modulo torsion.  This step does not require torsion-free source
$H_2$.

The only remaining rank-one hull is orbit $164$.  Its forbidden-actor
test leaves the same allowed thirty-element hull and its $2186$
closed action images.  The profile quotient preserves integral $H_2$,
because it fixes the minimal band and the row and column actions.
There are $1385$ such images of order at most twenty-two.  A fresh
exact Smith computation for \emph{every} one of them gives
\[
\begin{array}{c|rrrrr}
 \operatorname{rank}H_2(-;\Z)&1&2&3&4&8\\\hline
 \text{number of images}&846&437&86&15&1.
\end{array}
\]
Every one of the $846$ rational-rank-one images has torsion-free
$H_2=\Z$.  Their indices agree exactly with the previously verified
small-image scope.  Consequently no image with an additional torsion
summand was omitted by the earlier integral filter.

For $842$ of these images, the existing independent rank certificates
give $H_4(-;\mathbb F_5)=0$, hence $H_4(-;\Q)=0$; all source
degree-two squares pulled back from such an image vanish rationally.
The remaining image indices are $1291,1296,1335,1376$, all of order
twenty-two.  Their first three normalized contracted bar chain groups
have dimensions $8,48,288$, and the two exact integer matrix ranks are
\[
 \operatorname{rank}_{\Q}d_1=7,\qquad
 \operatorname{rank}_{\Q}d_2=40.
\]
Thus each has $\dim H_3(-;\Q)=48-7-40=1$.
The profile quotient is surjective on $H_3(-;\Q)$ by
Lemma~\ref{B:band-quotient-H3-surjection}, contradicting
$H_3(BM;\Q)=0$.  This is a rational third-homology obstruction;
positive homology over a finite field alone would not suffice here.

The standalone checker
\path{code/b_orbit164_rational_small_images.py}
imports neither the action-image producer nor its bar implementation.
It reconstructs the actor multiplication, all $593910$ image closure
products, all $1385$ Smith presentations, and both exceptional
differentials.  It checks all $9216$ entries of the four products
$d_1d_2$, and computes their ranks over the integers.
The complete output is
\path{code/outputs/b_orbit164_rational_small_images_verified.json}.
This supplies the additional rational scope and torsion checks needed
for the present statement.
\end{proof}

\begin{remark}\label{B:three-by-five-transfer-scope}
The assumption on rational homology explicitly includes
$H_3(BM;\Q)=0$; bounded rational homology together with one-dimensional
$H^2$ and $H^4$ alone is not being substituted for it.  Only the
restrictions to the eventual local groups appearing in the cited hull
certificates are needed.  The stronger all-finite-group formulation
above is convenient when those restrictions are supplied by a map to
a finite simply connected target.  The proposition concerns band
sources.  For order at most twenty-two a three-by-five minimal ideal
with a nontrivial maximal group is impossible, since that ideal alone
would have at least thirty elements.
\end{remark}

\subsubsection{The local five-dimensional quotient with two three-dimensional modules}

\begin{lemma}\label{B:top-heavy-Tor}
Let $A$ be a finite-dimensional local algebra over a field $k$, with
$A/\operatorname{rad}A=k$.  Let $U$ be a finite-dimensional left $A$-module,
and put
\[
 a=\dim_k(U/JU),\qquad b=\dim_k JU,
 \qquad J=\operatorname{rad}A.
\]
If $a>b$ and $W$ is a finite-dimensional right $A$-module satisfying
$\operatorname{Tor}^A_i(W,U)=0$ for all sufficiently large $i$, then $W$
is projective.  The same assertion holds with left and right interchanged.
\end{lemma}
\begin{proof}
Take a minimal free right resolution of $W$, with term $A^{\beta_i}$ in
degree $i$.  Every differential has entries in $J$.  After tensoring with
$U$, its image in degree $i-1$ lies in $(JU)^{\beta_{i-1}}$.  Consequently
\[
 \dim_k\operatorname{Tor}^A_i(W,U)
 \geq (a+b)\beta_i-b\beta_{i-1}-b\beta_i
 =a\beta_i-b\beta_{i-1}.
\]
Vanishing in all sufficiently large degrees forces
$\beta_i\leq(b/a)\beta_{i-1}$.  Since $b/a<1$, the nonnegative integer
sequence $\beta_i$ must eventually be zero.  A finite-dimensional module
of finite projective dimension over a finite-dimensional local algebra
is free: in a nontrivial finite minimal free resolution, the last
nonzero differential, whose entries lie in $J$, annihilates a nonzero
vector with entries in the appropriate socle of $A$, contradicting its
injectivity.  Thus $W$ is projective.
\end{proof}

\begin{proposition}\label{B:local-quotient-CP2}
Let $M$ be a finite monoid whose minimal ideal $K$ is a rectangular
band.  If $A=\mathbb Q M/\mathbb Q K$ is local and
$\dim_{\mathbb Q}H^2(BM;\mathbb Q)=1$, then every rational
degree-two class has square zero.  In particular $BM$ is not
homotopy equivalent to $\mathbb{CP}^2$.  This assertion has no bound
on the order of $M$ and does not require finite homotopy type.
\end{proposition}
\begin{proof}
Let $W,V$ be the column and row augmentation modules in
Lemma~\ref{lem:B-band-Tor}, and set $J=\operatorname{rad}A$ and
$D=A/J$.  There is a canonical surjection
\[
 W\otimes_A V\longrightarrow
 (W/WJ)\otimes_D(V/JV).
\]
The source has rational dimension one.  Both tops are nonzero, by
Nakayama's lemma, and $D$ is a finite-dimensional division algebra
over $\mathbb Q$.  The target therefore has rational dimension
$\dim_{\mathbb Q}D$ times the product of the two positive $D$-ranks.
Consequently $D=\mathbb Q$, both ranks are one, and the displayed
surjection is an isomorphism.  Every balanced functional on
$W\otimes_A V$ thus factors through the two one-dimensional tops.
Its bilinear matrix has rank at most one.  The square vanishes by
Lemma~\ref{B:rank-one-square}.
\end{proof}

\begin{proposition}\label{B:local-three-by-three-product}
Let $A$ be a five-dimensional local $k$-algebra with residue field $k$.
There are no right and left modules $W,V$, both of dimension three, for
which
\[
 \dim_k\operatorname{Tor}^A_i(W,V)=
 \begin{cases}2&i=0,\\1&i=2,\\0&i=1\text{ or }i\geq3.\end{cases}
\]
More generally, it suffices to assume the displayed dimensions for
$0\leq i\leq3$ and eventual vanishing of the higher Tor groups.
\end{proposition}
\begin{proof}
Neither module is projective, since a nonzero free $A$-module has
dimension divisible by five.  Lemma~\ref{B:top-heavy-Tor}, applied on
either side, shows that both modules have one-dimensional top and
two-dimensional radical.  In particular $V$ is cyclic.  Its first
syzygy has dimension two.  If this syzygy were not cyclic, its top
would have dimension two and its radical would be zero; dimension
shifting and Lemma~\ref{B:top-heavy-Tor} would make $W$ projective.
Thus the first syzygy is cyclic and its minimal projective cover is
$A$.  Its next syzygy has dimension three and must again be cyclic:
otherwise its top has dimension at least two and its radical has
dimension at most one, giving the same contradiction.  Repeating
this argument shows that every term in the minimal resolution of
$V$ is one copy of $A$, while its syzygies alternate in dimension
between three and two.

Tensor this resolution with $W$ and let $r_i$ be the rank of the
differential from degree $i$ to degree $i-1$.  All entries of the
minimal differentials lie in $J$, so $r_i\leq\dim WJ=2$.  The asserted
dimensions in degrees zero through three force, successively,
\[
 r_1=1,\qquad r_2=2,\qquad r_3=0,\qquad r_4=3.
\]
The last equality contradicts $r_4\leq2$.
\end{proof}

\begin{remark}\label{B:local-three-by-three-scope}
For a finite monoid whose minimal ideal is a $4\times4$ rectangular
band and which has five elements outside that ideal, the two
augmentation modules in the band Tor formula both have dimension
three.  Proposition~\ref{B:local-three-by-three-product} therefore
excludes a local contracted quotient for a classifying space with
the rational homology of $S^2\times S^2$.  The proposition makes no
corresponding assertion for the $\mathbb{CP}^2$ profile: in that case
the same rank equations are $2,1,1,2$, which do not by themselves
give a contradiction.
\end{remark}

\begin{proposition}\label{B:four-by-four-five-outside}
A finite monoid with a $4\times4$ minimal band and exactly five
outside elements cannot have type $\CP^2$ or $S^2\times S^2$.
\end{proposition}
\begin{proof}
Here $A=\Q M/\Q K$ has dimension five and its two augmentation
modules $V,W$ both have dimension three.  The local case is
excluded by Proposition~\ref{B:local-three-by-three-product} for the
product and Proposition~\ref{B:local-quotient-CP2} for the projective
plane.  The residue algebra is $\Q$, since its dimension divides
both five and three.  In the nonlocal case,
Proposition~\ref{B:five-nonlocal-exceptions} leaves only the line
and Nakayama algebras.

The Nakayama Tor table in Lemma~\ref{B:five-exceptional-targets}
has no dimension-$(3,3)$ pair with Tor dimensions
$(1\text{ or }2,0,1)$.  Indeed write $v_X,w_X$ for indecomposable
multiplicities, using that lemma's notation.  Degree-two Tor gives
$(w_{T_2}+w_R)(v_{S_2}+v_L)=1$.
If $w_R=1$, degree-one vanishing forces $v_L=v_{S_1}=w_{T_1}=0$;
there is no permitted one-dimensional complement to $R$ in $W$.
Otherwise $w_{T_2}=1$, and degree-one vanishing and dimensions force
$V=S_2\oplus Ae_2$, $W=T_2\oplus e_2A$.  Their degree-zero
Tor has dimension four, again impossible.

For the line algebra write $L_a=Ae_2$, $L_b=Ae_3$ and
$R_a=e_1A$, $R_b=e_2A$.  The degree-one and degree-two conditions
leave the following possible degree-zero dimensions:
\[
\begin{array}{c|ccc}
 &T_1\oplus2T_3&T_1\oplus R_a&T_1\oplus R_b\\\hline
2S_1\oplus S_3&4&4&2\\
S_3\oplus L_a&2&1&1\\
S_3\oplus L_b&4&0&1
\end{array}
\]
These entries follow directly from the two line resolutions in
Lemma~\ref{B:five-exceptional-targets}.
In the two product entries, every balanced form has a common
one-dimensional factor and hence has rank at most one.
Lemma~\ref{B:rank-one-square} and polarization kill all products.
The middle-right projective-plane entry likewise has a rank-one
balanced form.

It remains to exclude $(V,W)=(S_3\oplus L_a,T_1\oplus R_a)$
and its opposite $(S_3\oplus L_b,T_1\oplus R_b)$.
For the first pair the invisible radical arrow is $b$.
The kernel of the action on $V$ is exactly $\Q b$; the kernel
on $W$ is the two-dimensional ideal $Ae_3A$.
Apply the complete forty-eight-basis certificate of
Lemma~\ref{B:five-line-bases}, now to $b$.
In twenty-five cases it is a basis element, so both set maps are
constant and its profile lies in the band.
In the other twenty-three it is a difference $q-p$.
The row maps are nonconstant and equal: a basis element in
$\Q b$ would make $p,q$ linearly dependent.
Fourteen cases also have nonconstant column maps, which are
therefore equal.

The other nine have $p,q\in Ae_3A$, and the following literal
left-product witnesses force their constant column maps equal:
\[
\begin{array}{c|rrrr}
\text{record}&p&q&r&t\quad(rp=rq=t)\\\hline
0&1&2&1&1\\2&1&2&1&1\\4&3&4&3&3\\5&1&2&1&1\\
18&2&3&2&2\\21&2&3&2&2\\23&2&3&2&2\\
30&2&3&1&2\\32&2&3&1&2
\end{array}
\]
Here the multiplication convention gives
$\tau_p\tau_r=\tau_t=\tau_q\tau_r$.
Thus the action-profile quotient kills $b$ modulo the band in
all cases.  Its contracted algebra is a quotient of
$A/(b)\cong\Q\times\Q(1\to2)$; every such quotient is hereditary,
since there is at most one cross arrow and no path of length two.
Its natural $H_2$ is unchanged.  The square
therefore pulls back from zero fourth homology.
The opposite argument excludes the other rank-two pair.

The program \texttt{b\_line\_4x4\_cp2.py} constructs the kernel
ideal intrinsically: choose $x$ with characters $(0,0,1)$, put
$e_3=3x^2-2x^3$, and span $Ae_3A$.  This is idempotent because
the radical squares to zero.  The independent Fraction-arithmetic
checker recomputes the ideal and every witness in both orientations,
verifying all ninety-six table/pair cases.  Hence the multiplicative
basis issue is explicitly resolved in this application as well.
\end{proof}

\begin{corollary}\label{B:fourdim-lower22}
Every finite monoid of type $\CP^2$ or $S^2\times S^2$ has at
least twenty-two elements.
\end{corollary}
\begin{proof}
The previous bounds leave only order twenty-one.
The two-row, three-by-three and three-by-four obstructions, and
the five-outside-elements bound, leave only minimal bands of
shape $3\times5$ or $4\times4$, up to opposite.
Nontrivial minimal groups require too many elements after those
shape exclusions.  Theorem~\ref{B:three-by-five-targets} handles
the first shape and Proposition~\ref{B:four-by-four-five-outside}
handles the second.
\end{proof}

\subsubsection{The complete four-by-four classification for a primitive square}

\begin{proposition}\label{B:four-by-four-CP2-hulls}
Let $M$ be a finite monoid with a $4\times4$ minimal band and
$BM\simeq\CP^2$.  Its transformation image is contained, after
relabeling the rows and columns, in one of the three hulls numbered
$216$, $217$, and $221$ in the complete classification below.
The statement does not assume that this image, or its ambient hull,
has finite-dimensional classifying space.
\end{proposition}

\begin{proof}
Let $V$ and $W$ be the row and column augmentation lattices, both
of rank three.  An actor $(\sigma,\tau)$ contributes the nine
balancing relations
\[
 (w\tau)\otimes v-w\otimes(\sigma v)
 \quad(w\text{ and }v\text{ in the chosen augmentation bases})
\]
in $W\otimes V$.  The rational relation space $R$ of $M$ is the
sum of its actors' relation spaces.  Its maximal transformation
hull consists of all actors whose relations lie in $R$.
This is a monoid, contains every constant pair and the identity,
and has precisely relation space $R$: the actors of $M$ already
span it.  The natural band presentation therefore makes the map
from $M$ to this hull an isomorphism on rational $H_2$.
On integral $H_2$ it is surjective, since it is a quotient of the
same lattice by additional balancing relations.

All $4^4 4^4=65536$ transformation pairs were enumerated.
Their distinct proper individual relation spaces have counts
\[
\begin{array}{c|rrrrrrr}
\text{rank}&0&3&4&5&6&7&8\\\hline
\text{number}&1&12&1269&504&7732&1590&205.
\end{array}
\]
Their proper sums form exactly $169323$ spaces in $577$ orbits
under $S_4\times S_4$.  As in the three-by-five classification,
closure is checked by adjoining every individual relation space to
one representative of every orbit.  The independent dual-form
verifier checks all $6527601$ such joins.  Each proper join has
an exact annihilator containment and an equal modular rank bound,
so the conclusion is an exact rational equality.  It also checks
every actor and that each hull's actors generate its claimed space.
No bounded coordinate grid or partial monoid census enters this
classification.

There are $149$ orbits with one-dimensional rational $H^2$.
They divide as follows:
\[
\begin{array}{l|r}
\text{case}&\text{number of orbits}\\\hline
\text{balanced form of rank at most one}&12\\
\text{a common repeated row or column}&61\\
\text{new hull with }H_4(-;\mathbb F_5)=0&65\\
\text{remaining hulls}&11.
\end{array}
\]
Lemma~\ref{B:rank-one-square} excludes the first case.
For a repeated row or column, quotient the corresponding finite set
by equality under the balanced form.  Balance makes that equivalence
actor-invariant, and the form descends.  The resulting rational
$H_2$ map is an isomorphism, and its integral map is surjective.
The relative three-by-four argument of
Theorem~\ref{B:three-by-four-targets} applies exactly as in
Theorem~\ref{B:three-by-five-targets}: explicit zero products pull
back, and local-character obstructions are tested on the original
finite-dimensional source.  For its even-square case, integral
surjectivity from $H_2(BM)=\Z$ to a rationally nonzero group forces
that target group to be $\Z$ and preserves a primitive generator.
Thus no finite-dimensional hypothesis on the quotient is introduced.

For each of the next sixty-five hulls, the certificate lists
independent columns in the two differentials computing
$\operatorname{Tor}_2^{\Q N/\Q K}(W,V)$.
Fresh dense elimination over $\mathbb F_5$ verifies that their
ranks add to the full middle dimension.  By the band Tor formula,
$H_4(BN;\mathbb F_5)=0$; consequently $H_4(BN;\Q)=0$ and
every rational degree-two square pulled back from $N$ is zero.

The remaining eleven hulls are completely specified by the stored
actor lists and their composition formula.  Their relevant exact
cup calculations are:
\[
\begin{array}{c|rrl}
\text{orbit}&\text{order}&\text{content of the square}&\text{conclusion}\\\hline
210&58&0&H_4(-;\Q)=0\\
216&60&1&\text{bounded images treated below}\\
217&140&1&\text{bounded images treated below}\\
221&82&1&\text{bounded images treated below}\\
223&54&0&\text{square zero rationally}\\
359&56&2&\text{even square}\\
448&29&2&\text{even square}\\
454&39&2&\text{even square}\\
455&29&2&\text{even square}\\
495&29&2&\text{even square}\\
499&29&2&\text{even square}.
\end{array}
\]
Here ``content'' is the greatest common divisor of the square's
values on a saturated integral basis of all degree-four cycles in
an exact projective prefix
$P_4\to P_3\to P_2\to P_1\to P_0\to\Z$;
exactness is required through $P_3$, and the prefix can always be
extended to a resolution.  It is therefore not necessary to assume
or infer anything about an uncomputed tail.

For the eight zero or even cases, an independent checker reconstructs
the projective summands $\Z N(e-f)$, verifies the idempotents,
all differential equivariance equations and integral exactness, and
checks the saved comparison to the normalized bar resolution.
For each summand, checking the generator's idempotent support and
all its basis translates proves that this comparison is an actual
$\Z N$-linear map.  Its chain equation is then checked on every
projective generator.  The two-cocycle
\[
 u(m,n)=L\bigl(\tau_m(0),\sigma_n(0)\bigr),
 \qquad L(0,i)=L(j,0)=0,
\]
is verified on every triple, and its square is evaluated on the
actual bar comparison.  The degree-two cochain evaluates to $1$ or
$-1$ on the sole tensor generator, and both adjacent tensor maps
vanish; hence $H_2(BN;\Z)=\Z$ and this is its primitive dual.
The degree-four cycle basis is checked to have full kernel rank
and saturated integral image.  These checks include all $2557680$
sparse comparison terms for orbit $359$.

For content zero, the source's rational square is zero.
For content two, the primitive degree-two class pulls back
primitively, by the integral $H_2$ surjection above, while its square
evaluates evenly on every integral four-cycle.  Naturality makes
the same true on $M$, contradicting the unit square in $\CP^2$.
This excludes all eight hulls and leaves precisely the three stated.

The complete independent outputs are
\path{code/outputs/b_rational_orbit_lattice_4x4_verified.json},
\path{code/outputs/b_rational_4x4_CP_hulls_verified.json}, and
\path{code/outputs/b_rational_4x4_partial_cups_verified.json}.
Their scope includes every one of the $149$ target-rank orbits.
\end{proof}

\subsubsection{The bounded action images in the three remaining four-by-four hulls}

\begin{proposition}\label{B:four-by-four-bounded22}
No finite monoid of order at most twenty-two, with a
$4\times4$ rectangular-band minimal ideal, has classifying
space homotopy equivalent to $\mathbb{CP}^2$.
\end{proposition}
\begin{proof}
By Proposition~\ref{B:four-by-four-CP2-hulls}, it remains to treat
the three maximal balanced-form hulls numbered $217,216,221$.
We give the complete bounded-image calculation, including the
reason it applies to a possibly nonfaithful action homomorphism.

Let $M$ be a putative source and let $H$ be its image in the
monoid of row and column transformation pairs.  The image
contains all sixteen distinct constant pairs and the identity.
It has at most five further actors.  The natural balanced
presentation of $H_2$ shows that
\[
 H_2(BM;\mathbb Q)\longrightarrow H_2(BH;\mathbb Q)
\]
is an isomorphism: both presentations have the same row and
column augmentation spaces and exactly the same action
relations.  Thus every rational degree-two source class is
pulled back from the image.

An actor whose eventual cyclic group has a nonzero restriction
of the integral hull cocycle cannot occur in $H$.  Indeed its
preimage in $M$ has an eventual finite cyclic group mapping
onto that group.  Inflation of its nonzero integral character
remains nonzero.  Miller's theorem makes the map from this
finite local group into the finite CW type $BM$ null-homotopic,
a contradiction.  This is exactly the local-character test used
in Proposition~\ref{B:four-by-four-CP2-hulls}; the test is imposed
on the original source and does not assume that $BH$ has
finite homotopy type.

The remaining search is finite and completely exhaustive.
Starting with the sixteen constants and the identity, adjoin
one allowed actor and close under literal multiplication.
Reject a closure if it contains a forbidden actor or more than
five additional actors.  Repeat from every new closed set.
Every desired image is reached: adjoining any generating set
of that image produces only intermediate closures contained
in the image.  Products involving a constant remain constants,
so closure only needs the products of the additional actors.
The executable implementation is
\path{code/d_bounded_actor_closure.cpp}; its full inputs are
\path{code/outputs/d_bounded_hull_217.txt},
\path{code/outputs/d_bounded_hull_216.txt}, and
\path{code/outputs/d_bounded_hull_221.txt}.

The rational relation-rank filter is exact, not a test at an
uncontrolled prime.  Each raw balancing row in the nine form
coordinates has integral squared norm at most eight.  Hadamard's
inequality bounds the absolute value of every nonzero minor by
$8^{9/2}<65521$.  Since $65521$ is prime, reduction modulo this
prime preserves the rank of every such relation matrix.  Rank
eight is precisely the required rational $H_2$-rank one.

We quotient the output lists by all row and column permutations
that preserve the hull, also allowing the opposite operation
when it preserves that hull.  Every orbit is explicitly generated
and checked against the original complete list.  The counts are
\[
\begin{array}{c|r|r|r|rrr}
 \text{hull}&\text{closed sets}&H_2\text{-rank-one sets}
 &\text{orbits}&|H|=20&|H|=21&|H|=22\\\hline
 217&55363&37138&2360&52&519&1789\\
 216&8589&4866&2434&68&630&1736\\
 221&16686&10440&2636&76&659&1901
\end{array}
\]
There are no smaller rank-one images in these three hulls.
The complete lists, permutation maps and representatives are
the files with prefixes
\path{code/outputs/d_bounded_hull_217},
\path{code/outputs/d_bounded_hull_216}, and
\path{code/outputs/d_bounded_hull_221}.

For each of these $7430$ orbit representatives, the exact binary
two-sided bar calculation in
\path{code/d_bounded_hull_binary.cpp} computes $H_3$ and $H_4$.
If $q=|H|-17$, its degree-$k$ term before the band shift has
dimension $9q^k$.  Every composite of consecutive boundary
matrices is checked to be zero, and ranks are computed by
finite bit-vector Gaussian elimination.  Every image of order
twenty or twenty-one has $H_4(BH;\mathbb F_2)=0$.
The same is true of all order-twenty-two images except thirty-nine:
fourteen from hull $217$, ten from $216$, and fifteen from
$221$.  Each exception has binary Betti numbers $(1,0,1)$ in
degrees two, three and four.  In particular no image of order
less than twenty-two is discarded on the basis of a third- or
higher-homology mismatch alone.

Vanishing of binary fourth homology implies vanishing of
rational fourth homology, so the square is zero in all the
nonexceptional images.  For the remaining thirty-nine tables,
\path{code/d_bounded_hull_residual_cups.py} constructs an integral
projective resolution prefix and an actual comparison to the
normalized left bar resolution.  The independently implemented
checker \path{code/d_verify_bounded_cups.py} imports neither
construction.  It reconstructs every idempotent projective,
checks the multiplication and transformation actions, and uses
Smith normal forms to check saturation and exactness of
\[
 P_4\longrightarrow P_3\longrightarrow P_2
 \longrightarrow P_1\longrightarrow P_0
 \longrightarrow\mathbb Z\longrightarrow0
\]
at $P_3,P_2,P_1,P_0,\mathbb Z$.  It then checks every chain
square and every $\mathbb ZH$-equivariance equation of the
supplied comparison in degrees zero through four.

In each of the thirty-nine cases the coinvariant prefix has
terms $\mathbb Z,0,\mathbb Z,0,\mathbb Z$ in degrees zero
through four.  Evaluation of the integral hull cocycle on the
degree-two comparison generator is $1$ or $-1$.  Evaluation
of its cup square is zero on \emph{every} degree-four
projective basis vector.  Hence this is a primitive
degree-two class whose square vanishes, independently of
any assertion about the later resolution tail.  A projective
resolution prefix can always be extended, and the checked
comparison extends to the acyclic bar resolution, so this
degree-four cochain calculation proves the claimed
cohomological vanishing.

The full thirty-nine tables, cocycles and comparison images are
in \path{code/outputs/d_bounded_hull_residual_cups.json} and the
thirty-nine accompanying compressed cup certificates.  The
standalone verification record is
\path{code/outputs/d_bounded_hull_cups_independently_verified.json}.
Thus the square is zero in every possible bounded image $H$,
and its pullback to $M$ is zero.  This contradicts the
cohomology ring of $\mathbb{CP}^2$.
\end{proof}

\begin{proposition}[A lower bound on the faithful action image]
\label{B:four-by-four-image24}
If a finite monoid has a $4\times4$ rectangular-band minimal
ideal and classifying space homotopy equivalent to $\mathbb CP^2$,
its image in the monoid of row and column transformation pairs
has at least twenty-four elements.  This concerns the cardinality
of the image, with no upper bound on the source monoid.
\end{proposition}
\begin{proof}
Proposition~\ref{B:four-by-four-CP2-hulls} reduces the image to
one of the three allowed hulls $217,216,221$.  The image fixes
the band, and therefore has the same integral second homology
as the source.  Images of order at most twenty-two have zero
pulled-back rational square by the explicit fourth-homology and
cup certificates in Proposition~\ref{B:four-by-four-bounded22}.

For order twenty-three, run the same closed-subset enumeration
with six optional actors instead of five.  No search rule changes:
every intermediate closure is contained in the final image, so
the bounded closure procedure remains exhaustive.  The exact
counts are
\[
 \begin{array}{c|r|r|r|r}
 \text{hull}&\text{rank-one images}&\text{permutation/opposite orbits}
 &H_4(\mathbb F_2)\ne0&H_3=0,\ H_4\ne0\\\hline
 217&65464&4135&294&3\\
 216&5975&2992&188&0\\
 221&14522&3661&287&1.
 \end{array}
\]
The last two columns concern the orbit representatives and use
$\mathbb F_2$ coefficients.  The literal input tables and
complete lists are the files with prefixes
\path{code/outputs/d_cp23_hull_217},
\path{code/outputs/d_cp23_hull_216}, and
\path{code/outputs/d_cp23_hull_221}.
The binary bar calculation checks every consecutive differential
composite before taking ranks.

Zero fourth binary homology makes every rational square zero.
Nonzero third binary homology is now an obstruction even for
a nonfaithful source: Lemma~\ref{B:band-quotient-H3-surjection}
makes the source third homology surject onto that of its image.
Thus only the four representatives in the last column remain.
Their identifiers and exact integral projective-prefix ranks are
\[
 \begin{array}{c|c|ccccc}
 \text{hull}&\text{case}&P_0&P_1&P_2&P_3&P_4\\\hline
 217&1539&4&14&37&76&150\\
 217&1540&4&14&37&76&150\\
 217&4019&4&14&32&66&140\\
 221&3653&4&12&17&26&60.
 \end{array}
\]
Their coinvariant prefixes have dimensions $(1,0,1,0,2)$.
The integral affine hull cocycle pairs with the second-degree
generator by $-1,-1,1,-1$, respectively.  Its cup square is
identically zero on every degree-four projective basis vector
under the actual integral comparison to the normalized bar.
This proves square-zero cohomology, independently of any
uncomputed resolution tail.

The producer \path{code/a_cp23_image_cups.py} supplies the full
prefixes and sparse comparisons.  The standalone verifier
\path{code/a_verify_cp23_image_cups.py} imports no resolution
or comparison producer.  It reconstructs every projective,
checks integral exactness and saturation through $P_3$, all
equivariance and chain equations, the primitive degree-two
pairing, and the zero square on the full $P_4$ basis.  The
independent results are recorded in
\path{code/outputs/a_cp23_image_cups_independently_verified.json}.
Every possible image of order at most twenty-three is therefore
excluded, proving the proposition.
\end{proof}

\begin{proposition}[Four-by-four images under finite-target hypotheses]
\label{B:four-by-four-finite-target-transfer}
Let $M$ be a finite monoid with a $4\times4$ rectangular-band
minimal ideal. Suppose
\[
 \dim_{\mathbb Q}H^2(BM;\mathbb Q)=1,
 \qquad H_3(BM;\mathbb Q)=0.
\]
Assume $H^2(BM;\mathbb Z)=\mathbb Zv$, that the image of $v^2$
in $H^4(BM;\mathbb Z)$ modulo torsion is nonzero and primitive,
and that every integral degree-two class restricts to zero on
every local finite subgroup of $M$.
Then the faithful row-and-column action image of $M$ has at
least twenty-four elements. In particular, $|M|\leq23$ is
impossible. No finite-CW hypothesis on $BM$ or on its action
image is required, and $H_2(BM;\mathbb Z)$ may have torsion.
\end{proposition}
\begin{proof}
We first explain why the full-hull reduction in
Proposition~\ref{B:four-by-four-CP2-hulls} still applies.
Its maximal hull is defined by the rational balanced relation
space, so the action map is an isomorphism on rational second
homology and a surjection on integral second homology.
Rank-one forms, zero fourth rational homology, and explicit
zero-square comparisons therefore give zero source squares.
Repeated rows or columns reduce to the relative three-by-four
and three-by-three arguments. Their local-character exclusions
use only the restrictions in the hypothesis: the eventual cyclic
group of a source actor surjects onto that of its image, and
inflation preserves the nonzero integral character.

The remaining even-square hulls have $H_2=\mathbb Z$, with the
displayed degree-two class its primitive dual. An integral
surjection from a finitely generated group of free rank one to
$\mathbb Z$ is an isomorphism on free parts. Its primitive dual
thus pulls back to $v$ up to sign, even if the source second
homology has torsion. The even evaluations on a saturated
integral four-cycle basis mean that the square is divisible by
two modulo torsion; the same is then true after pullback.
This contradicts the assumed nonzero primitive free square.
These are all uses of the projective-plane or finite-CW
hypotheses in that hull reduction. The image consequently lies
in one of hulls $217,216,221$.

Suppose its order is at most twenty-three. The local-character
test in Proposition~\ref{B:four-by-four-bounded22} is valid by
the same restriction hypothesis. For images of order at most
twenty-two, all its completed certificates already give zero
pulled-back rational square: every image except the thirty-nine
explicit zero-square cases has zero fourth mod-two homology.
No third-homology prefilter is used for those smaller images.

For order twenty-three, use the complete six-optional-actor
lists in Proposition~\ref{B:four-by-four-image24}. Of their
$769$ representatives with nonzero fourth mod-two homology,
four have zero third mod-two homology and are the exact
integral zero-square cases already checked there. The other
$765$ now have an independently verified rational obstruction:
\[
\begin{array}{c|r|r}
 \text{hull}&\text{representatives}&\dim_{\mathbb Q}H_3\\\hline
 217&291&1\\
 216&188&1\\
 221&286&1.
\end{array}
\]
For clarity, the new calculation is on exactly the representatives
with both $H_3(-;\mathbb F_2)$ and $H_4(-;\mathbb F_2)$ nonzero;
it does not infer rational nonvanishing from binary nonvanishing.

The normalized two-sided bar terms in degrees zero, one and
two have ranks $9,54,324$. The producer
\path{code/d_cp23_rational_h3.py} computes their exact rational
ranks and saves, for every representative, an integral cycle
$z\in\ker d_1$ and an integral functional $\ell$ with
$\ell d_2=0$ and $\ell(z)\ne0$.
The standalone checker
\path{code/d_verify_cp23_rational_h3.py} imports no producer,
classifier, or matrix-rank library. It reconstructs the literal
tables, both band actions, and both differentials, then checks
all these cycle, cocycle, and pairing identities in integers.
It verifies $9307755$ associativity equations, $563040$ band-action
equations, $2230740$ consecutive-boundary equations, and all
$765$ nonzero pairings. The complete record is
\path{code/outputs/d_cp23_rational_h3_independently_verified.json};
the integral witnesses are in
\path{code/outputs/d_cp23_rational_h3_certificates.json.gz}.

The quotient to the action image fixes the band pointwise.
Lemma~\ref{B:band-quotient-H3-surjection}, over $\mathbb Q$,
makes its third rational homology a quotient of that of $M$.
Thus none of these $765$ images can occur. Every remaining
image has zero pulled-back rational square, contrary to the
nonzero free square of $v$. This proves the assertion.
\end{proof}


\begin{theorem}\label{B:CP2-minimum23}
The minimum order of a finite monoid whose classifying space is
homotopy equivalent to $\CP^2$ is exactly twenty-three.
An aperiodic monoid attains this minimum.
\end{theorem}
\begin{proof}
The explicit aperiodic monoid of Theorem~\ref{B:CP23} gives
the upper bound, including its complete finite projective resolution
and primitive cup square.  Corollary~\ref{B:fourdim-lower22}
already excludes orders at most twenty-one, so suppose that
$|M|=22$ and $BM\simeq\CP^2$.
The two-row and three-by-three obstructions exclude those minimal
ideal shapes.  A nontrivial minimal group in any remaining shape
would already require at least $3\cdot4\cdot2=24$ elements
in the ideal.  Thus the minimal ideal is a rectangular band.
The three-by-four obstruction excludes that shape, while
Corollary~\ref{B:five-outside-minimal-band} requires at least
five elements outside the band.  Hence its order is at most
seventeen.  The only shapes remaining, up to opposite, are
$3\times5$ and $4\times4$.
Theorem~\ref{B:three-by-five-targets} excludes the first and
Proposition~\ref{B:four-by-four-bounded22} excludes the second.
This proves the lower bound twenty-three.
\end{proof}

\begin{proposition}\label{B:four-by-four-product25}
A finite monoid of type $S^2\times S^2$ with a $4\times4$
minimal band has at least twenty-five elements.
\end{proposition}
\begin{proof}
Use the complete $577$-orbit relation lattice verified in
Proposition~\ref{B:four-by-four-CP2-hulls}.  There are $249$
orbits with rational degree-two dimension two: eight have every
balanced form of matrix rank at most one, seventy-one have a
common repeated row or column, and $170$ require a new calculation.
The first cases have all products zero by
Lemma~\ref{B:rank-one-square} and polarization.  The repeated-point
cases reduce relatively to the three-by-four obstruction, with all
local-character restrictions tested on the original finite-dimensional
source, just as in the preceding proofs.

Among the $170$ new hulls, $151$ have zero fourth homology over
$\mathbb F_5$.  Every one has a selected-column rank certificate
checked by independent dense elimination, proving rational vanishing
and hence zero products on a source.  The nineteen remaining hulls
divide into the following three classes.

First, nine hulls have a nonzero common radical for their rational
degree-two pairing.  In the stored basis of two balanced forms,
their cup evaluations on the full integral degree-four cycle lattice
are represented by
\[
\begin{array}{c|r|c}
\text{orbit}&\text{order}&\text{pairing, or all pairings zero}\\\hline
137&31&\operatorname{diag}(0,1)\\
161&39&0\\
169&42&\operatorname{diag}(-1,0)\\
172&44&\operatorname{diag}(-1,0)\\
188&31&\operatorname{diag}(0,-1)\\
222&25&\operatorname{diag}(0,2)\\
225&25&\operatorname{diag}(-2,0)\\
237&27&0\\
400&26&0.
\end{array}
\]
The zero rows mean that every product evaluates to zero on every
cycle, even where the cycle lattice has rank greater than one.
These are exact projective-prefix and normalized-bar comparisons,
independently verified by the same checker as the preceding
proposition.  It checks both degree-two classes form an integral
basis, and verifies a saturated basis of the entire degree-four
kernel.  Thus the radical assertion is rational cohomological
vanishing, and requires no statement about the later resolution.

The natural map from a source to its maximal hull is an isomorphism
on rational $H^2$.  A nonzero radical class therefore pulls back to
a nonzero class whose product with every degree-two source class
is zero.  This contradicts the nondegenerate intersection pairing
of $S^2\times S^2$.

Second, the nine hulls
\[
 141,189,476,477,482,518,519,551,552
\]
have all three independent products zero.  Each vanishing is
certified by an explicit integral degree-three cochain whose
coboundary is the product of the two specified affine cocycles.
The identities are checked on every normalized quadruple.
The two cocycles span the full rational $H^2$, both by the complete
balanced-form classification and by independent modular rank bounds
and integral grid-cycle pairings.  These ambient vanishing results
apply to every source, regardless of whether its transformation
image is faithful or finite-dimensional.

Only orbit $346$, of order thirty-six, remains.  Its two classes
have the hyperbolic cup matrix
$\left(\begin{smallmatrix}0&-1\\-1&0\end{smallmatrix}\right)$.
Eleven actors have nonzero restrictions of these integral classes
to their eventual cyclic groups.  As before, Miller's theorem
excludes their occurrence in the image of the original source.
The remaining twenty-five actors form a closed monoid, consisting
of the sixteen constants, the identity and eight optional actors.
All $2^8=256$ subsets of the latter were tested directly.
Exactly one hundred subsets give closed monoids.  Among those of
order at most twenty-four, exactly sixty-six have degree-two rank
two; all have torsion-free integral $H_2$.  Their orders are
\[
\begin{array}{c|rrrr}
\text{order}&20&21&22&23\\\hline
\text{number}&32&18&8&8.
\end{array}
\]
Every one of these sixty-six images has
$H_4(-;\mathbb F_5)=0$, verified by fresh dense elimination
on saved independent columns.  Consequently every possible image
of order at most twenty-four has zero rational products.  Its
natural $H_2$ agrees with that of the source, so pullback excludes
the source as well.  No mismatch in an image's higher homology
is used in place of this cup-product argument.

The complete verification records are
\path{code/outputs/b_rational_4x4_product_ranks_verified.json},
\path{code/outputs/b_rational_4x4_product_cups_verified.json},
the nine files \path{code/outputs/a_rational_hull_4x4_orbit*_cups.json},
and \path{code/outputs/b_4x4_product346_images_verified.json}.
The additional coverage checker verifies that these cases exhaust
all $249$ degree-two-rank-two orbits.
\end{proof}

\begin{corollary}\label{B:product-lower23}
Every finite monoid of type $S^2\times S^2$ has at least
twenty-three elements.
\end{corollary}
\begin{proof}
Corollary~\ref{B:fourdim-lower22} leaves only order twenty-two.
The same shape reduction as in Theorem~\ref{B:CP2-minimum23}
leaves only minimal bands of shape $3\times5$ or $4\times4$.
Theorem~\ref{B:three-by-five-targets} excludes the former and
Proposition~\ref{B:four-by-four-product25} excludes the latter.
\end{proof}

\subsection{Rank-two cup products through twenty-two elements}

\begin{lemma}[All small sources over the eight rank-two images]
\label{B:eight-rank-two-image-inflations}
Let $M$ have a $4\times4$ rectangular-band minimal ideal and at
most twenty-two elements. Suppose its faithful row-and-column
action image is, up to relabeling or opposition, one of the eight
order-twenty images in Lemma~\ref{C:six-rank-two-bounded-hulls}.
Then
\[
 H_3(BM;\mathbb Q)\ne0
 \quad\hbox{or}\quad
 H_5(BM;\mathbb Q)\ne0.
\]
\end{lemma}
\begin{proof}
Fix one of the eight image tables $H$. Label its identity by zero
and its sixteen constant pairs by the row-major labels $1,\ldots,16$.
Its other three profiles must each have an outside preimage in $M$.
These choices, together with the actual identity and band, account
for twenty elements. The remaining zero, one, or two elements may
have \emph{any} image profile, including identity and constant
profiles. Distinct outside elements with identical profiles are
retained as distinct elements.

For each multiset of extra profiles, products with the band are
forced by the actual transformations. A product of two outside
elements may be any element in the fibre of the corresponding
image product. In particular, when that product is a constant
profile, both its actual band element and any outside elements
with that profile are tried. Impose every associativity equation
and quotient only by permutations within the outside profile
fibres. Every possible source with the fixed image identification
is obtained. Opposition is harmless because it preserves the
classifying-space homotopy type.

The complete counts are
\[
\begin{array}{c|r|r|r|r}
 |M|&\text{classes}&\text{labelled tables}
     &H_3\ne0&H_3=0,\ H_5\ne0\\\hline
 20&8&8&0&8\\
 21&336&344&168&168\\
 22&12404&15052&6988&5416.
\end{array}
\]
Both homology columns use rational coefficients. The classes
are taken modulo fibre permutations for a fixed image identification, not a claim
that all $12748$ representatives are pairwise nonisomorphic after
arbitrary image automorphisms.

The producer is \path{code/d_rank_two_small_inflations.cpp}; the
literal inputs and output are
\path{code/outputs/d_rank_two_small_inflation_controls.json} and
\path{code/outputs/d_rank_two_small_inflations.tsv}.
The independent coverage check reconstructs the eight profile
tables, expands the supplied classes into disjoint full fibre
orbits, and separately regenerates all products in a different
cell order. Its $230138$ search nodes give exactly the same
$15404$ labelled tables, with $163523480$ literal associativity
equations checked. All twenty-four image/order scopes agree.
The complete independent record is
\path{code/outputs/b_rank_two_inflation_coverage_verified.json}.

For the homological columns, reconstruct from each literal table
the contracted algebra $A=\mathbb QM/\mathbb QK$, the left row
augmentation $V$, and the right column augmentation $W$.
The certificates are exact projective prefixes
\[
 P_d\longrightarrow\cdots\longrightarrow P_1
 \longrightarrow P_0\longrightarrow V\longrightarrow0,
 \qquad d=2\ \hbox{or}\ 4.
\]
Each term is a displayed sum of literal ideals $A\varepsilon$,
each embedded kernel has full rank and equals the preceding
kernel, and the next map surjects onto it. Tensoring gives
$W\otimes_A A\varepsilon=W\varepsilon$; every differential is
reconstructed from the actual image of its idempotent generator.
Thus the shorter prefixes compute $\operatorname{Tor}_1$ and
the longer prefixes compute $\operatorname{Tor}_3$, including
the incoming differential in each asserted degree.

For the order-twenty and order-twenty-one sources, the common
six-dimensional presentation is obtained by adjoining one or
two new identities to the contracted algebra. Adjoining a new
identity to a unital algebra gives its product with $\mathbb Q$.
Both modules are supported on the old identity, so this changes
none of their Tor groups. The eight nonsplit order-twenty-two
cases have a central cyclic averaging idempotent supporting
both modules. The stored projectives remain ideals in the full
algebra; no complementary summand is discarded without checking
its action.

The standalone checker
\path{code/d_verify_rank_two_inflation_homology.py} imports only
independent integer/Fraction linear algebra, and imports no
producer, classifier, Flint routine, or tensor constructor.
It rechecks all $12748$ literal tables, with $135253488$
associativity equations, and all source-to-module links. It
verifies $1957$ actual row prefixes, $7409$ exact projective
presentations, and $88908$ cover/kernel equivariance equations.
It then reconstructs every right column action and every tensor
differential independently. In all $7156$ cases in the fourth
column, $\operatorname{Tor}_1^A(W,V)$ is nonzero. In the other
$5592$ cases, $\operatorname{Tor}_3^A(W,V)$ is nonzero.
Lemma~\ref{lem:B-band-Tor} gives precisely the asserted third
and fifth homology groups. The complete independent record is
\path{code/outputs/d_rank_two_inflation_homology_verified.json}.
Only exact finite prefixes are used; no uncomputed tail or
finite-type property of the faithful image is assumed.
\end{proof}

\begin{theorem}[The bounded rank-two obstruction]
\label{B:rank-two-products-through22}
Let $M$ be a finite monoid of order at most twenty-two, and
suppose $BM$ has finite CW homotopy type. Assume
\[
 \begin{gathered}
 \dim_{\mathbb Q}H_2(BM;\mathbb Q)=2,\qquad
 \dim_{\mathbb F_2}H_2(BM;\mathbb F_2)=2,
 \\
 H_3(BM;\mathbb F_2)=0,\qquad
 H_5(BM;\mathbb Q)=0.
 \end{gathered}
\]
Then every product of two rational degree-two classes is zero.
In particular, the same conclusion holds under the simpler
integral hypotheses
\[
 H_2(BM;\mathbb Z)=\mathbb Z^2,\qquad
 H_3(BM;\mathbb Z)=0,\qquad H_5(BM;\mathbb Q)=0.
\]
\end{theorem}
\begin{proof}
Write the minimal ideal as a Rees ideal with $r$ rows, $s$
columns, and maximal group $G$. If either $r$ or $s$ is at most
two, every balanced form has matrix rank at most one. The rational
minimal-group collapse of Lemma~\ref{lem:B-collapse-minimal-group},
the rank-one primitive of Lemma~\ref{B:rank-one-square}, and
polarization make all rational products zero. These steps do not
require a finite-type assertion about the collapsed quotient.

Suppose $r,s\ge3$. If $G$ is nontrivial, the order bound forces
$r=s=3$ and $|G|=2$: the next possibility already requires at
least twenty-four elements in the minimal ideal. The finite-CW
three-by-three theorem \ref{B:three-by-three-cups} excludes all
products in this case. Hence assume $G$ is trivial, so the minimal
ideal is a band of order $rs$.

If $rs\ge18$, the contracted algebra has dimension
$q=|M|-rs\le4$. By Lemma~\ref{lem:B-band-Tor}, its Tor groups
are the homology of $BM$ shifted by two. Finite CW type makes
them eventually zero. Theorem~\ref{B:contracted-four-Tor} therefore
forces $\operatorname{Tor}_2=0$, hence $H_4(BM;\mathbb Q)=0$.
This again kills all products.

The remaining band shapes, up to opposition, are
$3\times3$, $3\times4$, $3\times5$, and $4\times4$.
The first is covered by Theorem~\ref{B:three-by-three-cups}.
For the next two, the rank-two parts of the proofs of
Theorems~\ref{B:three-by-four-targets} and
\ref{B:three-by-five-targets} prove that all degree-two products
on a finite-CW source are zero. Their rank-two arguments use
only the degree-two dimension and the original source's local-group
nullity, not the nondegenerate pairing of a product of spheres.

It remains to consider a $4\times4$ band. Use the complete
rank-two action-hull classification in the proof of
Proposition~\ref{B:four-by-four-product25}. The rank-one-form,
repeated-point, zero-fourth-homology, and explicit zero-product
branches all give zero products on the source. The three
common-radical hulls numbered $161,237,400$ also have every
product zero. Orbit $346$ has zero rational fourth homology in
every allowed image of order at most twenty-four, so its source
products vanish as well. The only remaining hulls are
\[
 137,169,172,188,222,225.
\]
We do not discard these merely because their pairings have a
radical line: a degenerate nonzero pairing is allowed by the
present hypotheses.

Let $H$ be the faithful action image of $M$. The quotient fixes
the band pointwise and gives exactly the same balanced-tensor
presentation over every coefficient ring. In particular it is
an isomorphism on second homology over $\mathbb Q$ and
$\mathbb F_2$. It is onto on third homology over $\mathbb F_2$
by Lemma~\ref{B:band-quotient-H3-surjection}. Thus
\[
 b_2(H;\mathbb F_2)=2,\qquad b_3(H;\mathbb F_2)=0.
\]
If $H_4(BH;\mathbb F_2)=0$, then $H_4(BH;\mathbb Q)=0$ and
every source product vanishes by naturality.

The forbidden cyclic-character actors in the six hulls cannot
occur. Each eventual cyclic image group lifts surjectively from
an eventual cyclic group of $M$, and a nonzero integral character
remains nonzero after this inflation. Miller's theorem makes
its restriction zero on the original finite-CW source. Therefore
the complete bounded classification in
Lemma~\ref{C:six-rank-two-bounded-hulls} applies. If a source
product is nonzero, its image has nonzero fourth rational,
and hence nonzero fourth mod-two, homology. It must consequently
be one of the eight specified order-twenty images.

Lemma~\ref{B:eight-rank-two-image-inflations} now gives positive
third or fifth rational homology for $M$. The first alternative
contradicts $H_3(BM;\mathbb F_2)=0$, and the second contradicts
$H_5(BM;\mathbb Q)=0$. This proves the theorem.

For the stated integral special case, a monoid with a rectangular
minimal band has trivial fundamental group: in its universal
group, every band idempotent is one and $me\in K$ forces every
$m$ to be one. Thus $H_1=0$, and the universal coefficient theorem
gives the required mod-two second and third homology. If the
minimal group is nontrivial or one of the other cases above
already applies, its products were excluded without these
mod-two assumptions. This establishes the integral version too.
\end{proof}

\begin{corollary}\label{B:CP2-wedge-S2-lower23}
No finite monoid with at most twenty-two elements has classifying
space homotopy equivalent to $\mathbb CP^2\vee S^2$.
\end{corollary}
\begin{proof}
This space has $H_2=\mathbb Z^2$, zero third and fifth homology,
and a degree-two class with nonzero square. Apply
Theorem~\ref{B:rank-two-products-through22}.
\end{proof}

\begin{remark}\label{B:rank-two-through22-scope}
The theorem does not assert a result for higher second-homology
ranks or for arbitrary sources over a small image. In particular
the unbounded homology of an image was never substituted for a
source obstruction: the finite third/fifth-homology calculations
are performed on every actual source through the stated bound.
For a $4\times4$ band, extending the source order from twenty-two
through twenty-five introduces precisely the three new outside
sizes $q=7,8,9$. It requires new faithful images as well as larger
fibres over the eight images considered here.
\end{remark}

\subsection{A rational and local-character form of the rank-two obstruction}

\begin{lemma}[Rational third homology in the remaining bounded images]
\label{B:rank-two-rational-image-gate}
In the six allowed hulls of Lemma~\ref{C:six-rank-two-bounded-hulls},
every closed image of order at most twenty-two and rational
degree-two rank two has one of the following properties:
\begin{enumerate}
\item its fourth rational homology is zero;
\item its third rational homology is nonzero;
\item it is one of the eight order-twenty images in
Lemma~\ref{B:eight-rank-two-image-inflations}.
\end{enumerate}
\end{lemma}
\begin{proof}
The previously independently verified complete closure and orbit
lists have $868$ rational-rank-two representatives. All $868$
also have $b_2(\mathbb F_2)=2$; none is discarded by a
mod-two second-homology prefilter. Of these, $711$ have zero
fourth mod-two homology, which implies zero fourth rational
homology. There are $149$ representatives with both third and
fourth mod-two homology nonzero, and the other eight are precisely
the stated order-twenty images. The new exact rational calculation
on the $149$ representatives gives
\[
\begin{array}{c|r|r|r}
 \text{hull}&\text{representatives}&b_3(\mathbb Q)=1&b_3(\mathbb Q)=2\\\hline
 137&73&47&26\\
 172&3&3&0\\
 188&73&47&26.
\end{array}
\]
Hulls $169,222,225$ have no such representatives.

The producer \path{code/d_rank_two_rational_h3.py} saves a literal
integral cycle and cocycle with nonzero pairing for each of the
$149$ two-sided bar complexes. The independent checker
\path{code/d_verify_rank_two_rational_h3.py} reconstructs the
tables, actions and differentials without importing a producer
or any matrix-rank library. It checks $1525524$ associativity
equations, $103488$ band-action equations, $269649$ consecutive
boundary equations, and all $149$ nonzero integral pairings.
The complete record is
\path{code/outputs/d_rank_two_rational_h3_independently_verified.json},
with the witnesses in
\path{code/outputs/d_rank_two_rational_h3_certificates.json.gz}.
Thus rational nonvanishing is proved directly in every case
where it is used.
\end{proof}

\begin{theorem}[Rank two under rational boundedness and local nullity]
\label{B:rank-two-rational-through22}
Let $M$ be a finite monoid with at most twenty-two elements.
Assume that
\[
 \dim_{\mathbb Q}H_2(BM;\mathbb Q)=2,\qquad
 H_3(BM;\mathbb Q)=H_5(BM;\mathbb Q)=0,
\]
and that $H_n(BM;\mathbb Q)=0$ for all sufficiently large $n$.
Assume also that every integral degree-two class on $BM$
restricts to zero on every local finite subgroup of $M$.
Then every product of two rational degree-two classes is zero.
No condition on mod-two homology or integral second-homology
torsion is required.
\end{theorem}
\begin{proof}
We retain the shape reduction in
Theorem~\ref{B:rank-two-products-through22}, recording explicitly
how its hypotheses are used. With at most two minimal rows or
columns, the rank-one balanced-form primitive and polarization
give zero products. This remains true with a nontrivial minimal
group, by the rational minimal-group collapse.

If both row and column counts are at least three and the minimal
group is nontrivial, the order bound leaves only a $3\times3$
shape with group of order two. The relative three-by-three
classification excludes its products. Its forbidden integral
characters are pulled back to the original local groups, where
they vanish by hypothesis; no finite-CW assertion about an
image or quotient is needed.

We may therefore assume the minimal ideal is a rectangular band.
If it has at least eighteen elements, the contracted algebra
has dimension at most four. The band Tor formula identifies
its Tor groups with homology shifted by two. Eventual rational
vanishing and Theorem~\ref{B:contracted-four-Tor} force
$\operatorname{Tor}_2=0$, so the fourth rational homology is zero.

The other shapes are $3\times3$, $3\times4$, $3\times5$ and
$4\times4$, up to opposition. In the first three, the rank-two
arguments in the complete hull classifications give zero
products. Each zero-homology witness or explicit cup primitive
pulls back directly. Their only uses of finite source type were
to kill integral characters on source local groups, exactly
the present local-nullity hypothesis.

For a $4\times4$ band, the same complete rank-two classification
leaves only hulls 137, 169, 172, 188, 222 and 225. The rank-one-form,
repeated-point, zero-fourth-homology and explicit zero-product
cases already vanish. The other common-radical hulls
$161,237,400$ have all products zero. Every allowed image of
order at most twenty-four in hull $346$ has zero rational
fourth homology. These exclusions use only naturality and the
stated source local-character restrictions. The six remaining
common-radical hulls are retained; a degenerate nonzero pairing
has not been excluded merely for being degenerate.

Let $H$ be the faithful action image in one of these six hulls.
Its quotient map fixes the band, preserves the balanced-tensor
presentation of rational second homology, and is onto on third
rational homology by Lemma~\ref{B:band-quotient-H3-surjection}.
In particular $H_3(BH;\mathbb Q)=0$. Forbidden actors cannot
occur, since their nonzero integral cyclic characters inflate
to restrictions of integral source classes.
Lemma~\ref{B:rank-two-rational-image-gate} now shows that either
the image has zero fourth rational homology, or it is one of the
eight order-twenty images. In the first case all source products
vanish, because every rational source degree-two class pulls
back from the image. In the second case the complete actual-source
calculation in Lemma~\ref{B:eight-rank-two-image-inflations}
forces $H_3(BM;\mathbb Q)\ne0$ or $H_5(BM;\mathbb Q)\ne0$.
Both contradict the hypotheses. This completes the proof.
\end{proof}

\begin{corollary}\label{B:finite-CW-rank-two-rational22}
If $|M|\leq22$, $BM$ has finite CW homotopy type,
$\dim_{\mathbb Q}H_2(BM;\mathbb Q)=2$, and
$H_3(BM;\mathbb Q)=H_5(BM;\mathbb Q)=0$, then every rational
degree-two product on $BM$ is zero.
\end{corollary}
\begin{proof}
Finite CW type supplies eventual homology vanishing, and
Miller's theorem supplies the local integral-character nullity.
Apply Theorem~\ref{B:rank-two-rational-through22}.
\end{proof}

\begin{remark}\label{B:rank-two-finite-target-scope}
The theorem also applies when a map $BM\to X$ to a finite CW
complex induces a surjection on integral second cohomology,
provided the displayed rational homology hypotheses hold.
Indeed each integral class then pulls back from $X$, and
Miller's theorem kills its restriction to every local finite
group. This formulation asserts neither finite CW type of
$BM$ nor a result for second-homology rank greater than two.
\end{remark}

\begin{proposition}\label{B:local-two-by-five-products}
Let $A$ be a local rational algebra of dimension five or six, and
let $W,V$ be right and left modules of dimensions five and two.
Suppose
\[
 \dim_{\mathbb Q}(W\otimes_A V)=2,\qquad
 \dim_{\mathbb Q}\operatorname{Tor}_2^A(W,V)=1,
\]
and the higher Tor groups eventually vanish.  Then every balanced
bilinear form on $W\times V$ factors through a fixed
one-dimensional quotient of $V$ and has rank at most one.
The analogous statement holds with left and right interchanged.
\end{proposition}
\begin{proof}
Put $J=\operatorname{rad}A$ and $D=A/J$.  The nonzero top of $V$
is a $D$-module, so $\dim_{\mathbb Q}D$ is one or two.  If it is
two, then $V=D$ is simple and $\operatorname{Tor}_2^A(W,V)$ is
a right $D$-module, contradicting its odd rational dimension.
Thus $D=\mathbb Q$.  If $V$ has two-dimensional top, it is
$\mathbb Q^2$ and its degree-two Tor also has even dimension.
Consequently $V$ is cyclic with one-dimensional top.

The module $W$ cannot be cyclic.  For $\dim A=5$, a cyclic
five-dimensional module is $A$ itself, contrary to nonzero Tor.
For $\dim A=6$, the kernel of its cyclic projective cover would
be the one-dimensional simple right module $\mathbb Q$.
Dimension shifting would then make
$\operatorname{Tor}^A_n(\mathbb Q,V)$ vanish eventually.
Minimal resolutions over the local algebra, as in
Lemma~\ref{B:local-minimal-Tor}, force $V$ to be projective,
which is impossible for a two-dimensional module over $A$.

The canonical surjection
\[
 W\otimes_A V\longrightarrow
 (W/WJ)\otimes_{\mathbb Q}(V/JV)
\]
therefore has target dimension at least two.  Since its source
has dimension exactly two, it is an isomorphism and
$\dim(W/WJ)=2$.  Every balanced functional factors through
these two tops, the second of which is one-dimensional.
This is the required rank-one factorization.
\end{proof}

\begin{corollary}\label{B:local-three-by-six-product}
A finite monoid with a $3\times6$ or $6\times3$ rectangular-band
minimal ideal and a local contracted rational quotient of
dimension five or six cannot have classifying space homotopy
equivalent to $S^2\times S^2$.
\end{corollary}
\begin{proof}
The row and column augmentation dimensions are two and five.
The band Tor formula supplies the hypotheses of
Proposition~\ref{B:local-two-by-five-products}.  Every rational
degree-two class has square zero by
Lemma~\ref{B:rank-one-square}; polarization shows that all
products of two such classes vanish.  This contradicts the
cohomology ring of $S^2\times S^2$.
\end{proof}

\subsubsection{Three-by-six bands with five outside elements}

\begin{lemma}\label{B:five-two-by-five-pairs}
In the two exceptional algebras of
Proposition~\ref{B:five-nonlocal-exceptions}, suppose
\[
 \dim V=2,\quad\dim W=5,\qquad
 (\dim\operatorname{Tor}_0,
   \dim\operatorname{Tor}_1,
   \dim\operatorname{Tor}_2)=(2,0,1).
\]
For the zero-composite line algebra the only pair is
\[
 V=S_1\oplus S_3,\qquad
 W=T_1\oplus e_1A\oplus e_2A.
\]
Every balanced form for this pair factors through the same
one-dimensional quotient $S_1$ of $V$.
For the Nakayama algebra the only pair is
\[
 V=L=Ae_1/\mathbb Qc,\qquad
 W=T_2\oplus(e_2A)^{\oplus2}.
\]
Both modules of this latter pair are faithful over
$A/I$, where $I=\operatorname{span}(a,c)$, and
$A/I$ is the path algebra of one arrow.
\end{lemma}
\begin{proof}
Use the decompositions and exact Tor table in
Lemma~\ref{B:five-exceptional-targets}.
For the line algebra the same degree-one and degree-two
conditions force $V=S_1\oplus S_3$ and
\[
 W=T_1\oplus T_3^{\oplus s}\oplus(e_1A)^{\oplus p}
                   \oplus(e_2A)^{\oplus q},\qquad
 s+2p+2q=4.
\]
The tensor dimension is $1+s+p=2$, whose unique solution is
$(s,p,q)=(0,1,1)$.  Both nonzero tensor summands use only
the $S_1$ coordinate of $V$, proving the common factorization.

For the Nakayama algebra the degree-one and degree-two
conditions force $V=L$ and
$W=T_2\oplus(e_1A)^{\oplus p}\oplus(e_2A)^{\oplus q}$.
Now $3p+2q=4$, so $(p,q)=(0,2)$; its tensor dimension is two.
The actions of $1,e_1,b$ on $L$ are linearly independent, and
the same is true on the right module $e_2A$.  Their kernels
as $A$-representations are therefore exactly
$I=\operatorname{span}(a,c)$.  Quotienting by this ideal
leaves two vertex idempotents and one arrow, as claimed.
\end{proof}

\begin{lemma}[The actual Nakayama multiplicative bases]
\label{B:five-Nakayama-common-annihilator}
Suppose a contracted finite-monoid algebra of dimension five
is isomorphic to the Nakayama algebra in
Lemma~\ref{B:five-two-by-five-pairs}.  In its distinguished
monoid basis, the ideal $I=\operatorname{span}(a,c)$ is spanned
by basis vectors lying in $I$ and differences of two basis
vectors that have the same image in $A/I$.
\end{lemma}
\begin{proof}
This is a finite multiplicative-basis assertion, checked against
the complete independently generated order-six census, including
all reducible monoids and all monoids with zero.
The checker \path{code/d_nakayama_five_annihilator.py} examines
all $2325$ distinct literal table/opposite choices supplied by
the $1373$ isomorphism-and-opposite classes in
\path{code/outputs/c_monoids_6.tsv}.

For clarity, the algebra-recognition and ideal test are exact.
After contracting the unique zero, enumerate all rational
characters, whose values on a finite monoid belong to
$\{-1,0,1\}$.  Require exactly two independent characters;
their common kernel $J$ has dimension three.  Require
$\dim J^2=1$ and $J^3=0$, so this kernel is exactly the
radical.  Lift the two orthogonal idempotents and order them
so that the radical corner dimensions are
\[
 (\dim e_1Je_1,\dim e_1Je_2,
   \dim e_2Je_1,\dim e_2Je_2)=(1,1,1,0).
\]
The additional tests
$(e_1Je_2)(e_2Je_1)\ne0$ and
$(e_2Je_1)(e_1Je_2)=0$ identify precisely the stated
Nakayama algebra.  This finds fourteen oriented bases, from
seven source census classes.

In each case the checker forms the intrinsic ideal
$I=e_1Je_2+J^2$, computes the images of every distinguished
basis vector in the three-dimensional quotient, and spans
the zero-image vectors and equal-image differences.  Exact
rational row reduction shows that this span equals $I$
in all fourteen cases.  The full tables, idempotent-independent
ideal bases, quotient coordinates, and spanning relations are
recorded in
\path{code/outputs/d_nakayama_five_annihilator.json}.
There is no assumption that an abstract algebra isomorphism
preserves a particular displayed monoid basis.
An independent implementation,
\path{code/b_verify_nakayama_five_bases.py}, repeats the entire
classification using rational \texttt{Fraction} arithmetic and a
different idempotent lift.  It finds the same fourteen oriented
bases and verifies every ideal-spanning equality; its results are
in \path{code/outputs/b_nakayama_five_bases_verified.json}.
\end{proof}

\begin{proposition}\label{B:three-by-six-five-outside}
A finite monoid with a $3\times6$ or $6\times3$ rectangular-band
minimal ideal and exactly five elements outside that ideal
cannot have classifying space homotopy equivalent to
$S^2\times S^2$.
\end{proposition}
\begin{proof}
Take the opposite if necessary so that the augmentation modules
have dimensions $\dim V=2$, $\dim W=5$, and put
$A=\mathbb QM/\mathbb QK$.  The band Tor formula gives the
Tor dimensions $2,0,1,0,\ldots$.  A local quotient is excluded
by Corollary~\ref{B:local-three-by-six-product}.  Otherwise
Proposition~\ref{B:five-nonlocal-exceptions} and eventual Tor
vanishing leave only the two pairs in
Lemma~\ref{B:five-two-by-five-pairs}.

For the line pair every balanced form has a common
one-dimensional factor.  Lemma~\ref{B:rank-one-square} and
polarization make all degree-two products zero.

For the Nakayama pair, let $H$ be the faithful transformation-pair
image of $M$.  We claim that the induced map of contracted
rational algebras kills $I$.  Each basis vector in $I$ has zero
row and column augmentation operators, so its two set maps
are constant and its profile belongs to the minimal band.
For two basis vectors equal modulo $I$, their augmentation
operators agree on both sides.  Since both $V$ and $W$ are
faithful over $A/I$, the common operators are either both
nonzero or both zero.  In the former case each pair of set
maps is equal: two finite set maps with the same nonzero
augmentation operator must coincide.  To see this directly,
equality of the operators makes
$[f(x)]-[g(x)]$ independent of $x$; a nonzero such difference
forces both maps to be constant.  In the zero case both
profiles already lie in the band.  Thus the image kills every
spanning relation in Lemma~\ref{B:five-Nakayama-common-annihilator},
and hence kills $I$.

The contracted algebra of $H$ is consequently a quotient of
the path algebra of one arrow.  Every quotient of this specific
algebra is hereditary: it either retains the single arrow or
is a product of copies of $\mathbb Q$.  The band Tor formula
gives $H_4(BH;\mathbb Q)=0$.  The natural map
$H_2(BM;\mathbb Q)\to H_2(BH;\mathbb Q)$ is an isomorphism,
because the balanced action presentations are identical.
All degree-two classes of $M$ therefore pull back from $H$,
and all their products vanish.  This again contradicts the
cohomology ring of $S^2\times S^2$.
\end{proof}

\begin{remark}\label{B:Nakayama-two-by-five-rank-correction}
The Nakayama balanced forms themselves have rank two.  On one
$e_2A$ summand, in the bases $(e_2,b)$ and $(\bar e_1,\bar b)$,
the multiplication pairing is
$\left(\begin{smallmatrix}0&1\\1&0\end{smallmatrix}\right)$.
The proof uses the actual multiplicative-basis quotient and
does not infer rank-one factorization from
$e_2A\otimes_A L\cong e_2L$.
As a separate construction check,
\path{code/d_product23_nakayama.py} enumerates all $732$
affine lifts for the standard zero basis: four row-action
types and $183$ column-action types.  Every resulting
twenty-three-element monoid has the required finite additive
homology, but all $2\times2$ cup matrices are zero.  This
last family calculation is not used in the exhaustive
multiplicative-basis argument above.
\end{remark}


\begin{corollary}\label{B:product-lower24}
Every finite monoid of type $S^2\times S^2$ has at least
twenty-four elements.
\end{corollary}
\begin{proof}
Corollary~\ref{B:product-lower23} leaves only order twenty-three.
The previous two-row, three-by-three and three-by-four obstructions
exclude those minimal ideal shapes; a nontrivial minimal group
in any other shape would already require at least twenty-four
ideal elements.  The minimal ideal is therefore a rectangular
band, with at least five elements outside it by
Corollary~\ref{B:five-outside-minimal-band}.
Up to opposite, only $3\times5$, $4\times4$ and $3\times6$
remain.  The first two are excluded by
Theorem~\ref{B:three-by-five-targets} and
Proposition~\ref{B:four-by-four-product25}.
The last has exactly five outside elements, and is excluded by
Proposition~\ref{B:three-by-six-five-outside}.
\end{proof}

\subsubsection{Mixed products of rank-one forms on three rows}

\begin{lemma}\label{B:three-row-rank-one-products}
Suppose the minimal ideal of a finite monoid is a rectangular
band with three rows.  If two rational degree-two classes have
rank-one balanced matrices, their product is zero.  The same
statement holds with rows and columns interchanged.  No finiteness
assumption on the homotopy type is required.
\end{lemma}
\begin{proof}
Write the two balanced forms as $L_i=b_i\otimes a_i$, for $i=1,2$.
Use the affine representatives and notation of
Lemma~\ref{B:rank-one-square}: normalize the row-value functions
$a_i(r)$ to vanish at a distinguished row, and similarly normalize
the column-value functions.  There are rational characters $\chi_i$
such that
\[
 a_i(\sigma_m r)=a_i(m)+\chi_i(m)a_i(r),\qquad
 b_i(\tau_m c)=b_i(m)+\chi_i(m)b_i(c),
\]
and the class of $L_i$ is represented by
$u_i(m,n)=b_i(m)a_i(n)$.  The corresponding multiplication identities
are
\[
 a_i(mn)=a_i(m)+\chi_i(m)a_i(n),\qquad
 b_i(mn)=b_i(n)+\chi_i(n)b_i(m).
\]

If the two row functionals are dependent, every linear combination
of $L_1,L_2$ has rank at most one.  Their individual squares and
the square of their sum vanish by Lemma~\ref{B:rank-one-square};
polarization over $\mathbb Q$ gives $[u_1][u_2]=0$.

Otherwise the three functions $1,a_1,a_2$ form a basis of the
functions on the three-point row set.  Since both $a_i$ vanish at
the distinguished row, there are rational numbers $s,t$ with
\[
 a_1(r)a_2(r)=s a_1(r)+t a_2(r)
 \quad\text{for every row }r.
\]
Apply this identity to $\sigma_m r$ and substitute the affine
formulas.  The coefficient of $a_2(r)$ is
\[
 \chi_2(m)\bigl(a_1(m)+t(\chi_1(m)-1)\bigr)=0.
\]
Set $a'_1(m)=a_1(m)+t(\chi_1(m)-1)$.  This is still an affine
cocycle with character $\chi_1$, and
\[
 u'_1(m,n):=b_1(m)a'_1(n),\qquad
 u'_1-u_1=-t\delta b_1.
\]
Define a rational three-cochain by
\[
 w(m,n,p)=b_1(m)a'_1(n)b_2(n)a_2(p).
\]
Expanding the bar coboundary using the two multiplication
identities gives the exact equality
\begin{align*}
 (\delta w)(m,n,p,q)
 &=u'_1(m,n)u_2(p,q)\\
 &\quad+b_1(m)b_2(n)\chi_1(n)
              \chi_2(p)a'_1(p)a_2(q)\\
 &=u'_1(m,n)u_2(p,q).
\end{align*}
The last equality uses $\chi_2(p)a'_1(p)=0$.  Thus the mixed
product is a coboundary.  Applying the result to the opposite
monoid proves the column version.
\end{proof}

\begin{corollary}\label{B:three-row-semisimple-cups}
Let $M$ have a three-row rectangular-band minimal ideal, and let
$V$ be its two-dimensional rational row augmentation module.
If $V$ is a direct sum of one-dimensional modules, every product
of two rational degree-two classes of $BM$ is zero.
\end{corollary}
\begin{proof}
An equivariant decomposition of $V$ into its two one-dimensional
summands decomposes every balanced map $V\to W^*$ into two
balanced maps of rank at most one.  All pairwise products of
these summands vanish by Lemma~\ref{B:three-row-rank-one-products}.
\end{proof}

\begin{lemma}\label{B:six-dimensional-simple-row}
Let $A$ be a rational algebra of dimension at most six, and let
$V$ be a simple left $A$-module of rational dimension two.
For every right module $W$, either $V$ is projective or
$\dim_{\mathbb Q}\operatorname{Tor}^A_n(W,V)$ is even for every
$n$.  In particular, a one-dimensional degree-two Tor group
cannot occur.
\end{lemma}
\begin{proof}
Write $D=\operatorname{End}_A(V)$.  The division algebra $D$ has
rational dimension one or two.  In the second case every Tor
group is a module over $D$ and has even rational dimension.
In the first case the corresponding simple factor of
$A/J$, where $J=\operatorname{rad}A$, is $M_2(\mathbb Q)$.
Every nonzero bimodule joining this factor to another simple
factor has rational dimension at least two; a bimodule from
this factor to itself has dimension at least four.

If this is the only factor of $A/J$, then $\dim J\leq2$ makes
its component in $J/J^2$ zero, and nilpotence gives $J=0$.
Otherwise $\dim(A/J)\geq5$ and $\dim J\leq1$, so again no
component of $J/J^2$ can meet the $M_2(\mathbb Q)$ factor.
Products of radical components show that no component of any
$J^n/J^{n+1}$ meets it either.  Hence this matrix factor splits
off as a semisimple direct factor of $A$, and its simple module
$V$ is projective.
\end{proof}

\begin{corollary}\label{B:three-by-six-order24-row-reduction}
For a possible twenty-four-element monoid model of
$S^2\times S^2$ whose minimal ideal is a $3\times6$ band,
the two-dimensional row augmentation module over the
six-dimensional contracted rational algebra must be
indecomposable and nonsemisimple.  Its image algebra in
$\operatorname{End}_{\mathbb Q}(V)$ is either
$\mathbb Q[\epsilon]/(\epsilon^2)$ or the algebra of upper
triangular two-by-two rational matrices.
\end{corollary}
\begin{proof}
The band Tor formula gives $\dim\operatorname{Tor}_2(W,V)=1$.
The preceding lemma excludes a two-dimensional simple $V$,
and Corollary~\ref{B:three-row-semisimple-cups} excludes a sum
of one-dimensional simples.  A decomposable two-dimensional
module would have precisely that latter form.  Thus $V$ is a
nonsplit extension of two rational characters.  Choose a basis
adapted to its one-dimensional submodule.  If the characters
are equal, the image algebra is the dual-number algebra; if
they are distinct, it contains the two diagonal idempotents
and a nonzero off-diagonal matrix, and is the full triangular
algebra.
\end{proof}

\begin{proposition}\label{B:three-row-dual-cup-radical}
Suppose the minimal ideal of a finite monoid $M$ is a rectangular
band with three rows, and its rational row augmentation module
$V$ has image algebra $\mathbb Q[\epsilon]/(\epsilon^2)$ acting
nonsemisimply.  Then either all products of rational degree-two
classes vanish, or their cup-product pairing has a nonzero common
radical.  Consequently $BM$ cannot have the cohomology ring of
$S^2\times S^2$.  If $\dim H^2(BM;\mathbb Q)=1$, its generator
has square zero as well.
\end{proposition}
\begin{proof}
First identify the actual row maps, including their affine data.
The image algebra contains a nonzero nilpotent matrix.  Some row
actor must have a nonscalar augmentation operator, which in a
dual-number algebra has the form $\lambda I+\mu N$ with
$\mu\ne0$ and $N^2=0$.  Its powers belong to a finite set.
In characteristic zero this forces $\lambda=0$: a nonzero
$\lambda$ would make the nontrivial Jordan part have infinitely
many powers.  Thus this actor is a three-point chain.  After
relabeling its map is $(0,0,1)$, and in the augmentation basis
$(e_1-e_0,e_2-e_0)$ its matrix is
\[
 N=\begin{pmatrix}0&1\\0&0\end{pmatrix}.
\]
A three-point map whose augmentation operator belongs to
$\operatorname{span}(I,N)$ is one of the following six maps:
\[
 (0,1,2),\quad (0,0,1),\quad(1,1,0),\quad
 (0,0,0),\quad(1,1,1),\quad(2,2,2).
\]
This list follows directly by equating the two columns
$e_{\sigma(1)}-e_{\sigma(0)}$ and
$e_{\sigma(2)}-e_{\sigma(0)}$ with those of $\lambda I+\mu N$.
Their augmentation operators are respectively $I,N,-N,0,0,0$.

Let $a(m)$ be the two-coordinate vector
$e_{\sigma_m(0)}-e_0$, and let $x$ be the row functional with
$x(e_1-e_0)=0$, $x(e_2-e_0)=1$.  Write
$\alpha(m)=x(a(m))$ and let $\chi(m)$ be the common diagonal
character of the augmentation matrix $A_m$.
Inspection of the six maps gives the exact identity
\[
 \alpha(m)A_m=0.
\]
Indeed $\alpha(m)$ is nonzero only for the constant map to row
$2$.  Also
$\alpha(mn)=\alpha(m)+\chi(m)\alpha(n)$.

Every rank-one balanced form has row factor $x$: the nonzero
nilpotent $N$ has a unique invariant line, so a rank-one
equivariant map from $V$ has kernel $\operatorname{im}N$.
Write its affine cocycle as
$v(m,n)=r(m)\alpha(n)$, where
$r(mn)=r(n)+\chi(n)r(m)$.
An arbitrary balanced form has affine cocycle
$u(m,n)=b(m)a(n)$, where $b(m)$ is a row vector and
\[
 a(mn)=a(m)+A_m a(n),\qquad
 b(mn)=b(n)+b(m)A_n.
\]
Define
\[
 w(m,n,p)=r(m)\alpha(n)u(n,p).
\]
Using the cocycle identity for $u$ and the displayed affine
identities, the full coboundary is
\begin{align*}
 (\delta w)(m,n,p,q)
 &=v(m,n)u(p,q)\\
 &\quad+r(m)\chi(n)\alpha(p)b(n)A_p a(q)\\
 &=v(m,n)u(p,q).
\end{align*}
Thus every rank-one balanced class annihilates every
degree-two class.

Finally, $N$ commutes with the action algebra.  If a balanced
map $L:V\to W^*$ has rank two, then $LN$ is a nonzero balanced
map of rank one.  The natural identification of balanced forms
with $H^2$ makes this a nonzero cohomology class in the common
radical.  If no balanced map has rank two, every class has rank
at most one and the preceding coboundary formula makes all
products zero.  These are all rank cases.  A one-dimensional
space of balanced forms cannot contain a rank-two map, since
$L$ and $LN$ would be linearly independent; this proves the
last assertion.
\end{proof}

\begin{corollary}\label{B:three-by-six-order24-triangular}
In the situation of
Corollary~\ref{B:three-by-six-order24-row-reduction}, the row
augmentation image is the full triangular two-by-two rational
matrix algebra.  Thus the possible twenty-four-element product
case has a nonsplit extension of two distinct rational characters
as its row augmentation module.
\end{corollary}
\begin{proof}
The preceding proposition excludes the remaining dual-number
alternative, since the degree-two pairing of $S^2\times S^2$
has zero radical.
\end{proof}

\subsubsection{Two structural reductions for six-dimensional quotients}

\begin{lemma}\label{B:six-nonsplit-triangular-row}
Let $A$ be a rational algebra of dimension six and let $V$ be a
two-dimensional module whose action algebra is the full triangular
two-by-two rational matrix algebra.  If $A/\operatorname{rad}A$
is not a product of copies of $\mathbb Q$, then eventual vanishing
of $\operatorname{Tor}^A_n(W,V)$ implies
$\operatorname{Tor}^A_2(W,V)=0$ for every finite-dimensional right
module $W$.
\end{lemma}
\begin{proof}
The nonsplit extension defining $V$ uses two distinct rational
characters.  Thus $A/J$ contains two copies of $\mathbb Q$, and
$J/J^2$ has a nonzero, at least one-dimensional, component joining
their vertices.  Every other simple factor which is not
$\mathbb Q$ has dimension at least two.  A nonzero radical
bimodule incident to a field factor of degree at least two has
dimension at least two; the same lower bound holds for a matrix
factor.  Such a radical component together with the required
rational cross-arrow would exceed the radical dimension available
in a six-dimensional algebra.  A matrix factor of dimension four
already uses all six semisimple dimensions after the two rational
characters and leaves no radical at all.

Consequently every non-rational simple factor is isolated in
$J/J^2$, and therefore in every $J^n/J^{n+1}$ by multiplication.
It splits off as a semisimple direct factor.  The factor acting
nontrivially on $V$ is a nonlocal split basic rational algebra of
dimension at most four.  The nonlocal part of
Theorem~\ref{B:contracted-four-Tor} applies to every such algebra,
independently of a monoid basis: eventual Tor vanishing forces its
degree-two Tor group to vanish.  Direct-factor decomposition gives
the claimed conclusion for $A$.
\end{proof}

\begin{lemma}[An odd Euler characteristic at a local corner]
\label{B:six-local-corner-odd-Euler}
Let $B$ be a four-dimensional local rational algebra with residue
field $k=\mathbb Q$, and put
\[
 A=\begin{pmatrix}k&k\\0&B\end{pmatrix},
\]
where the right $B$-action on the displayed cross bimodule is its
residue character.  Let $V$ be the two-dimensional left module
with one copy of $k$ at each vertex and with the cross-arrow
acting as the identity.  For every five-dimensional right
$A$-module $W$, the groups $\operatorname{Tor}^A_n(W,V)$ are
nonzero in arbitrarily high degrees.
\end{lemma}
\begin{proof}
Write $f,e$ for the first and second vertex idempotents.
A right $A$-module is a triple
\[
 W=(F,U,t),\qquad F=Wf,\quad U=We,\quad
 t:F\otimes_k k\longrightarrow U,
\]
where $U$ is a right $B$-module and $t$ is $B$-linear when its
domain has the residue action.  Put $d=\dim F$, so
$\dim U+d=5$.  There is an explicit exact sequence of right
$A$-modules
\[
 0\longrightarrow(0,k^d)
 \xrightarrow{\ (\iota,-t)\ }
 (F\otimes_k fA)\oplus(0,U)
 \longrightarrow W\longrightarrow0.
\]
The last map is the identity at the $f$-vertex and sends the
$e$-vertex pair $(z,u)$ to $t(z)+u$.  Its kernel is
$z\mapsto(z,-t(z))$, proving exactness and specifying every map.

The module $fA$ is projective and $fA\otimes_A V=fV=k$.
The induction $(0,U)=U\otimes_B eA$ is exact and sends free
$B$-modules to projective $A$-modules, since $eA=(0,B)$.
Consequently
\[
 (0,U)\otimes_A^{\mathbb L}V
 \simeq U\otimes_B^{\mathbb L}k.
\]
Tensor the displayed short exact sequence with $V$ in the derived
category.  Projection away from the summand $F$ induces a map
\[
 W\otimes_A^{\mathbb L}V
 \longrightarrow
 C\otimes_B^{\mathbb L}k,
 \qquad
 C=\operatorname{Cone}(t:k^d\longrightarrow U).
\]
Its fiber is the vector space $F$ in degree zero.  Equivalently,
the source is the cone of
$k^d\otimes_B^{\mathbb L}k\to
 F\oplus(U\otimes_B^{\mathbb L}k)$ with components the
augmentation and $-t$, whereas the target discards the first
component.  Thus this map is a homology isomorphism in every
degree at least two.

Suppose the Tor groups in the statement vanish eventually.
Then $C\otimes_B^{\mathbb L}k$ has homology in only finitely
many degrees.  A bounded complex of finite-dimensional modules
over a finite-dimensional local algebra has a minimal free
resolution, constructed successively by projective covers.
Its differential entries lie in the radical.  Tensoring with
$k$ therefore makes every differential zero; the assumed
bounded homology forces this minimal resolution to be bounded.
Hence $C$ is a perfect complex of $B$-modules.

Every finite projective $B$-module is free and has rational
dimension divisible by four.  The rational Euler characteristic
of a perfect $B$-complex is therefore divisible by four.  But
the explicit two-term complex $C$ has Euler characteristic
\[
 \chi_{\mathbb Q}(C)=\dim U-d=5-2d,
\]
which is odd.  This contradiction proves the assertion.
\end{proof}

\begin{remark}\label{B:six-local-corner-both-orientations}
For the opposite triangular orientation
$\left(\begin{smallmatrix}B&k\\0&k\end{smallmatrix}\right)$,
the two-dimensional nonsplit cross-arrow module is itself the
projective at the $k$-vertex.  Its higher Tor groups are zero.
Thus the two orientations of a local four-dimensional corner
with a single radical-annihilated cross-arrow cannot produce
Tor dimensions $2,0,1,0,\ldots$ with augmentation dimensions
$(\dim W,\dim V)=(5,2)$.
\end{remark}

\subsubsection{The complete three-point row-module reduction}

\begin{lemma}\label{B:three-point-triangular-target}
Let a finite monoid act on a three-element set, and suppose the
algebra generated on its rational augmentation module is the full
triangular two-by-two matrix algebra.  After relabeling the set as
$\{0,1,2\}$, every acting transformation $s$ satisfies
\[
 s(0)=s(1)\quad\hbox{or}\quad\{s(0),s(1)\}=\{0,1\}.
\]
Collapsing the three constant maps gives a thirteen-element monoid
$S_{13}$.  The induced augmentation representation factors through
this monoid.  The only remaining relabeling that preserves the
unique invariant line is the transposition of $0$ and $1$.
\end{lemma}
\begin{proof}
The full triangular algebra has a unique invariant rational line.
If all nonconstant acting maps were permutations, their finite
multiplicative image would be a group; adjoining the zero
augmentation operator would not change its semisimple rational
representation.  Thus some nonconstant, nonpermutation map acts.
Its augmentation operator has rank one.  Any line invariant under
a rank-one operator is either its kernel or its image.  For a
three-point transformation of rank two, its kernel is the line
spanned by the difference of the two identified points, and its
image is the difference line of its two image points.  The common
invariant line is therefore a difference line of two points.
Relabel them $0,1$.

The displayed condition is exactly preservation of
$\mathbb Q([1]-[0])$.  There are twelve nonconstant transformations
satisfying it: six identify $0,1$ and six permute them.  The three
constant transformations form an ideal, so their contraction gives
$S_{13}$ with twelve nonzero elements and a zero.  The difference
line determines the unordered pair $\{0,1\}$, proving the last
assertion.
\end{proof}

\begin{lemma}\label{B:split-syzygy-persistence}
Let $A$ be a finite-dimensional algebra, $V$ a left module, and
let $K$ be a first syzygy of $V$ after removal of projective direct
summands.  Suppose that for some $d\geq1$, $K$ is a direct summand
of a $d$th syzygy of $K$, up to projective direct summands.
If $\operatorname{Tor}_2^A(W,V)\ne0$, then
$\operatorname{Tor}_n^A(W,V)\ne0$ for arbitrarily large $n$.
\end{lemma}
\begin{proof}
Dimension shifting identifies
$\operatorname{Tor}_2(W,V)$ with $\operatorname{Tor}_1(W,K)$.
The split inclusion and retraction, followed by dimension
shifting, make $\operatorname{Tor}_n(W,K)$ a direct summand of
$\operatorname{Tor}_{n+d}(W,K)$ for every $n\geq1$.
Projective summands contribute zero in those degrees.  Iterating
from $n=1$ proves the claim.
\end{proof}

\begin{proposition}[Complete six-dimensional triangular-row reduction]
\label{B:six-row-census}
Let $A$ be a six-dimensional contracted finite-monoid algebra over
$\mathbb Q$.  Let $V$ be the two-dimensional augmentation module
of an action on three points, with full triangular action algebra,
and let $\dim W=5$.  Suppose
\[
 (\dim\operatorname{Tor}_0^A(W,V),
   \dim\operatorname{Tor}_1^A(W,V),
   \dim\operatorname{Tor}_2^A(W,V))=(2,0,1)
\]
and the Tor groups vanish in all sufficiently large degrees.
Up to isomorphism of the actual contracted monoid and rational
isomorphism of $V$, the pair $(A,V)$ belongs to one of the
following two explicitly enumerated lists:
\[
\begin{array}{c|r}
\text{type}&\text{row-module records}\\\hline
 J^2=0,\ \text{two parallel arrows followed by one arrow}&280\\
 J^2=0,\ \text{oriented triangle with zero length-two products}&178\\
 \text{Nakayama algebra of dimension five times }\mathbb Q&104\\
 x^2=ab=c\text{ with the thin }b\text{-module}&20.
\end{array}
\]
Here the first two entries have three vertices and three arrows;
in the oriented triangle one arrow is the direct bypass of the
other two.  The last two entries are those of
Lemma~\ref{B:six-pd2-row-pairs}.
\end{proposition}
\begin{proof}
We give the complete finite calculation and its independent checks.
Contracting a minimal ideal leaves a monoid with seven elements,
including its zero.  The source census
\path{code/outputs/d_zero_monoids_7.tsv}
contains $10115$ classes up to isomorphism and opposite.
The independent program \path{code/c_verify_zero7.cpp} uses a
different associativity search, with diagonal-first and then
column-major branching.  It counts exactly $1725414$ labeled
tables with prescribed identity and zero.  Expanding every source
class under all $120$ permutations of the other five elements and
opposition gives disjoint orbits with exactly that union; it also
verifies every source associativity equation.
The audit is recorded in
\path{code/outputs/c_zero7_independent_verified.json}.
Thus this step includes reducible monoids, monoids with one-sided
zeros, and nonfaithful action possibilities.

Adding distinct literal opposite tables gives $18344$ oriented
source tables.  The regular trace form, exact rational characters,
and nilpotent radical checks separate them into $227$ semisimple,
$5420$ local, $12475$ nonlocal split-basic, and $222$ remaining
algebras.  Semisimple algebras have zero higher Tor.  A local algebra
cannot have a quotient equal to the triangular algebra, which has
two nonisomorphic simple modules.  This leaves $12697$ source
algebras for the row search.

For each source table, enumerate all zero-preserving unital
homomorphisms to $S_{13}$ from
Lemma~\ref{B:three-point-triangular-target}.  The target matrices
in the basis $[1]-[0],[2]-[0]$ are given explicitly by
\[
 A_s(i,j)=\mathbf1_{s(j)=i}-\mathbf1_{s(0)=i},
 \qquad i,j\in\{1,2\}.
\]
Every target product is literal composition, with a constant
composite replaced by zero.  The search
\path{code/b_three_point_row_homs.cpp}
uses multiplication constraints and checks every completed map.
The independent program
\path{code/b_verify_three_point_row_homs.py}
instead chooses actual generators of each source, assigns their
possible images, and extends by multiplication until either a
collision or a complete homomorphism occurs.  The two algorithms
agree on every one of the $12697$ sources.  Up to the transposition
of $0,1$, there are $57888$ nonsemisimple row types: $24196$ have
dual-number image and $33692$ have triangular image.  Only the
latter are under consideration here.

For precision, write their operators as
$\left(\begin{smallmatrix}\alpha_m&c_m\\0&\beta_m\end{smallmatrix}\right)$.
The characters $\alpha,\beta$ are distinct.  An upper triangular
change of basis changes $c$ by a nonzero scalar and a multiple of
$\alpha-\beta$; the unique invariant line shows that these are
exactly the possible rational module isomorphisms.  Normalizing
this extension vector gives $20730$ rational module records.
The independent verifier also constructs and checks an actual
invertible intertwining matrix for every literal type and checks
that every one of the $33692$ types occurs exactly once in this
partition.

There are $138$ records over nonsplit or nonbasic algebras;
Lemma~\ref{B:six-nonsplit-triangular-row} excludes them.  For the
remaining split-basic records, the program constructs the actual
left projectives $Ae_i$, an epimorphism onto $V$, and its rational
kernel.  Exactly $8545$ records are projective and $5166$ have
projective dimension one, so their degree-two Tor vanishes.
The remaining $6881$ first syzygies are treated as follows:
\[
\begin{array}{c|r}
\text{verified outcome}&\text{records}\\\hline
 J^2=0,\ \operatorname{Tor}_2\ne0\text{ forces unbounded Tor}&4921\\
 J^2=0,\ \text{the two acyclic quivers in the statement}&458\\
 J^2\ne0,\ \text{explicit split syzygy recurrence}&1324\\
 J^2\ne0,\ \text{projective dimension two}&124\\
 J^2\ne0,\ \text{four-dimensional local corner with one arrow}&54.
\end{array}
\]

Here is the exact reason for the first line.  When $J^2=0$, the
first syzygies on both sides are semisimple.  A vertex contributing
to $\operatorname{Tor}_2(W,V)=\Omega W\otimes_A\Omega V$
must occur in $\Omega V$ and have an incoming arrow, since
$\Omega W$ lies in the radical of a right projective.  If every
such possible common vertex lies on a directed cycle, the
path description of the resolution of a simple module gives
nonzero Tor in arbitrarily large degrees.  The exact reachability
check gives $4921$ such records.  Every other square-zero record
has one of the two displayed acyclic three-arrow quivers, with
counts $280$ and $178$.

For the $1502$ records with $J^2\ne0$, the computation uses actual
projective covers and exact rational kernels.  A projective direct
summand is removed only when explicit module maps $F,G$ satisfy
$GF=1$; the next module is the invariant complement $\ker G$.
In $1324$ cases, within at most four further syzygies the initial
stable kernel is a split summand of the new one.  The saved
inclusion and retraction are actual matrices, so
Lemma~\ref{B:split-syzygy-persistence} excludes these cases.
A search failure is never treated as an exclusion.
In $124$ cases the next kernel is projective, giving the two
finite-projective-dimension families in the statement.
The $54$ remaining records all have the actual triangular algebra
$\left(\begin{smallmatrix}\mathbb Q&\mathbb Q\\0&B\end{smallmatrix}\right)$
with $B$ local of dimension four and the thin cross-arrow module.
Their corner multiplication and row action are checked explicitly;
Lemma~\ref{B:six-local-corner-odd-Euler} excludes them.

The complete matrix records are in
\path{code/outputs/b_three_point_row_syzygies.json} and
\path{code/outputs/b_three_point_higher_syzygies.json}.
The standalone checker
\path{code/b_verify_six_row_reduction.py}
imports none of the search, cover, or syzygy implementations.
It reconstructs the target transformations and every oriented
source table, verifies the regular trace ranks, the radical and
orthogonal idempotents, the row-module partition, every exact
kernel and projective cover, every split inclusion and retraction,
and the complete square-zero path filter.  All $20730$ records
pass; the eight outcome counts above sum to $20730$.
The final audit is
\path{code/outputs/b_six_row_reduction_verified.json}.
This proves the stated exhaustive reduction.
\end{proof}

\subsubsection{The finite-projective-dimension triangular row cases}

\begin{lemma}\label{B:six-pd2-row-pairs}
Let $A$ be either the five-dimensional Nakayama algebra of
Lemma~\ref{B:five-exceptional-targets} times $\mathbb Q$, or the
six-dimensional algebra with vertex idempotents $e,f$, radical
basis $x,a,b,c$, and
\[
 x\in eAe,\quad a\in eAf,\quad b\in fAe,\quad c\in eAe,
 \qquad x^2=ab=c,
\]
all other products of radical basis elements being zero.
Let $V$ be the two-dimensional left module on which $b$ is the
nonzero cross-arrow and all other radical elements act by zero.
In the direct-product case the isolated rational factor acts by zero.
Suppose a five-dimensional right module $W$ satisfies
\[
 (\dim\operatorname{Tor}_0^A(W,V),
   \dim\operatorname{Tor}_1^A(W,V),
   \dim\operatorname{Tor}_2^A(W,V))=(2,0,1).
\]
Then $W$ and $V$ have the same annihilator, and are faithful over
the quotient by that annihilator, which is the path algebra of one
arrow.  In the Nakayama product case the annihilator is
$(a,c)\oplus\mathbb Q$; in the six-dimensional case it is $(x,a,c)$.
\end{lemma}
\begin{proof}
In the Nakayama product case write $W=W_N\oplus\mathbb Q^r$.
The Tor table in Lemma~\ref{B:five-exceptional-targets} forces
\[
 W_N=T_2\oplus(e_1A)^{\oplus p}\oplus(e_2A)^{\oplus q}.
\]
The tensor and total dimensions give $p+q=2$ and
$1+3p+2q+r=5$, hence $p=r=0$ and $q=2$.
The faithfulness assertion is the same explicit one-arrow action
as in Lemma~\ref{B:five-two-by-five-pairs}.

For the six-dimensional algebra the following is an exact left
projective resolution, with maps specified by the images of their
vertex generators:
\[
 0\longrightarrow Af\xrightarrow{\ b\ }Ae
   \xrightarrow{\ x\ }Ae\longrightarrow V\longrightarrow0.
\]
Indeed, the first map has image $\operatorname{span}(b,c)$,
which is the kernel of right multiplication by $x$; the latter
has image $\operatorname{span}(x,c)$, the kernel of the quotient
to $V$.  Put $u=\dim We$, $v=\dim Wf$, and write $r,s$ for
the ranks of right multiplication by $x,b$, respectively.
After tensoring the resolution, the three homology dimensions are
\[
 u-r=2,\qquad u-r-s=0,\qquad v-s=1,
 \qquad u+v=5.
\]
Thus $(u,v,r,s)=(2,3,0,2)$.  In particular $x$ acts by zero and
$b:Wf\to We$ is surjective.  The relation $ba=0$ forces $a$ to
act by zero as well, and $c=x^2$ acts by zero.
The surviving $b$-map has rank two; its two vertex projections
and its cross-arrow are linearly independent operators.
Consequently $W$ is faithful over $A/(x,a,c)$, as is $V$.
\end{proof}

\begin{proposition}\label{B:six-pd2-actual-bases}
None of the $124$ finite-projective-dimension triangular row
modules retained by the six-dimensional quotient census can
occur in a monoid with a $3\times6$ rectangular-band minimal
ideal and classifying space $S^2\times S^2$.
\end{proposition}
\begin{proof}
The complete row-module reduction is recorded below in
Proposition~\ref{B:six-row-census}.  Its $124$ cases comprise
$104$ modules over a Nakayama algebra times $\mathbb Q$, and
$20$ modules over the six-dimensional algebra in
Lemma~\ref{B:six-pd2-row-pairs}.  Their row modules are precisely
the modules $V$ in that lemma.  These identifications are checked
by an explicit rational change of algebra basis, followed by all
$36$ multiplication equations, rather than by quiver shape alone.

Here is the additional check on the distinguished monoid bases.
For each basis, form the subspace $L$ spanned by basis elements
whose $V$-operator is zero and by differences of basis elements
having the same $V$-operator.  Direct multiplication verifies
that $L$ is an ideal.  In $98$ Nakayama product cases and in all
$20$ six-dimensional cases, $L$ is the full common annihilator.
In the remaining six Nakayama product cases, $L=(a,c)$; the
isolated rational factor survives.  Thus in every case $A/L$ is
either the path algebra of one arrow or its product with
$\mathbb Q$.

The exact data are in
\path{code/outputs/b_six_pd2_annihilators.json}.
The independent checker
\path{code/b_verify_six_pd2_bases.py} reconstructs all the stated
algebra bases, verifies their complete multiplication, verifies
the row-module actions, and checks each spanning relation and
ideal inclusion.  Its results, including every rational basis
change, are in
\path{code/outputs/b_six_pd2_bases_verified.json}.
It certifies $98$, $20$, and $6$ cases in the three groups above.

To pass from algebra to monoids, let $H$ be the faithful
transformation-pair image of the original monoid.  By the preceding
lemma $W$ and $V$ are faithful over the same algebra quotient.
Thus two distinguished basis elements equal on $V$ are equal on
$W$, and their common augmentation operators are either both
nonzero or both zero.  Equal nonzero augmentation operators imply
equal set maps; zero operators on both sides give constant maps,
whose profile lies in the minimal rectangular band.  It follows
that the contracted algebra of $H$ kills every generator of $L$.
There is no mixed zero/nonzero case here.

Every quotient of the particular algebra $A/L$ is hereditary:
there is at most one arrow, with distinct endpoints, and an
optional isolated vertex.  The band--Tor formula therefore gives
$H_4(BH;\mathbb Q)=0$.  The row and column actions give the same
balanced tensor presentation before and after taking the faithful
image, so $H_2(BM;\mathbb Q)\to H_2(BH;\mathbb Q)$ is an
isomorphism.  All degree-two classes pull back from $BH$, and
all their products vanish.  This excludes $S^2\times S^2$.
\end{proof}

\subsubsection{The two square-zero six-dimensional exceptions}

\begin{lemma}\label{B:six-squarezero-target-modules}
Let $A$ be one of the following radical-square-zero rational
algebras, with orthogonal primitive idempotents $e_1,e_2,e_3$.
Suppose $\dim W=5$, $\dim V=2$, and
\[
 \bigl(\dim\operatorname{Tor}_0^A(W,V),
       \dim\operatorname{Tor}_1^A(W,V),
       \dim\operatorname{Tor}_2^A(W,V)\bigr)=(2,0,1).
\]
The conclusions below concern every such module pair, without
assuming that the displayed algebra basis is a monoid basis.
\begin{enumerate}
\item For the triangle algebra, take
$a\in e_1Ae_2$, $b\in e_2Ae_3$, $c\in e_1Ae_3$ and let $V$
be the two-dimensional module on which $c$ is the nonzero arrow.
Either every balanced form has rank at most one, or $W$ is
faithful over $A/(a)$ or over $A/(b)$.  In the first faithful
case the vertex dimensions of $W$ are $(1,1,3)$; in the second
they are $(2,1,2)$.
\item For the parallel-chain algebra, take
$a\in e_1Ae_2$ and $b,c\in e_2Ae_3$, with $b,c$ independent.
Let $V$ have $b$ acting as the identity between its one-dimensional
vertex spaces and $c$ acting as zero.  Either every balanced form
has rank at most one, or $W$ is faithful over $A/(a)$, with vertex
dimensions $(1,1,3)$.
\end{enumerate}
Every quotient algebra appearing in the faithful alternatives is
hereditary, as is each of its further algebra quotients.
\end{lemma}
\begin{proof}
Write $W_i=We_i$.  In the triangle case the exact projective
resolution of $V$ is
\[
 0\longrightarrow Ae_1\xrightarrow{a}Ae_2
 \xrightarrow{b}Ae_3\longrightarrow V\longrightarrow0.
\]
Here the first arrow sends $e_1$ to $a$, the second sends $e_2$
to $b$, and the last map retains $e_3,c$ and kills $b$.
After tensoring, the complex is
$W_1\xrightarrow{a}W_2\xrightarrow{b}W_3$.
Put $r=\operatorname{rank}a$ and $s=\operatorname{rank}b$.
The three Tor dimensions imply
\[
 \dim W_1=r+1,\qquad \dim W_2=r+s,\qquad
 \dim W_3=s+2.
\]
Their sum is five, so $r+s=1$.

A balanced form is specified by a functional on
$W_3/\operatorname{im}b$.  Its value on the $e_1$ coordinate of
$V$ is obtained by composing that functional with
$c:W_1\to W_3$.  If $c=0$, all forms therefore factor through
the same one-dimensional quotient of $V$ and have rank at most
one.  Otherwise, if $(r,s)=(0,1)$, the action kills $a$ and is
faithful over $A/(a)$: each of its three vertex spaces is nonzero,
and its remaining two arrows $b,c$ act nontrivially in different
Peirce matrix blocks.  If $(r,s)=(1,0)$, the identical argument
gives faithfulness over $A/(b)$.  The displayed vertex dimensions
follow from the three equations above.

In the parallel-chain case, the projective resolution has the
same form with its second arrow $c$ in place of $b$.
Put $r=\operatorname{rank}a$ and $s=\operatorname{rank}c$;
again $r+s=1$ with the same vertex dimensions.  If $r=1$, the
map $a:W_1\to W_2$ is onto, and the relation $ab=0$ forces
$b:W_2\to W_3$ to be zero.  Every balanced form then factors
through the top coordinate of $V$.  If $r=0$, the source of
both $b,c:W_2\to W_3$ has dimension one, and $c$ is nonzero.
When their images are dependent, $b$ is zero after passage to
$W_3/\operatorname{im}c$, so all forms again have the same
one-dimensional row factor.  When they are independent, the
three idempotent actions and the two arrow actions are linearly
independent, and $W$ is faithful over $A/(a)$.

The triangle quotients are the path algebras of two-arrow forks.
The parallel-chain quotient is a two-parallel-arrow algebra
together with an isolated vertex.  These quivers have no
composable pair of arrows.  Every quotient merely deletes
vertices or takes a quotient of an arrow space and is again
a path algebra of this form; hence it is hereditary.
\end{proof}

\begin{lemma}[The literal arrow-ideal gate]
\label{B:six-squarezero-literal-ideals}
Among the $458$ actual row-module cases in the two square-zero
exceptions of Proposition~\ref{B:six-row-census}, the faithful
alternatives of Lemma~\ref{B:six-squarezero-target-modules}
give $636$ one-dimensional arrow ideals to test.  All $280$
parallel-chain ideals and $348$ of the $356$ triangle ideals
are killed by the faithful transformation-pair image of every
monoid realizing the corresponding module pair.  Their images
therefore have zero rational fourth homology and all degree-two
products in the original monoids are zero.

The only eight exceptions to this literal-ideal test have the
source indices and ideals listed in the next proposition.
\end{lemma}
\begin{proof}
The checker \path{code/d_six_square_zero_ideals.py} uses the
complete zero-seven census, the exact primitive idempotents,
and the actual row representations in
\path{code/outputs/b_three_point_row_syzygies.json}.
If $t$ is the top vertex of the row module, $m$ is the vertex
of its one-dimensional first syzygy, and $s$ is the unique
vertex with $e_sJe_m\ne0$, the first tested ideal is
$I=e_sJe_m$.  In a triangle it also tests $I=e_mJe_t$.
Each is checked by exact multiplication to be a one-dimensional
two-sided ideal.

Here is the complete logical content of each successful finite
certificate.  The ideal is spanned by distinguished basis vectors
lying in it and differences $p-q$ of distinguished basis vectors
with the same image modulo it.  A basis vector in $I$ has zero
augmentation operator on both $V$ and $W$, so both of its set
maps are constant and its transformation profile is already in
the minimal band.  For a difference $p-q$, the two row operators
and the two column operators respectively agree.  If the common
row operator is nonzero, the common column operator is also
nonzero: $W$ is faithful over $A/I$ and $V$ factors through that
quotient.  A finite set map is uniquely determined by its nonzero
augmentation operator, so both maps agree literally.

If both operators are zero, both profiles belong to the band.
The remaining case has zero common row operator and nonzero
common column operator.  The certificate then supplies a literal
basis element $z$ with
\[
 pz=qz\ne0
\]
in the contracted monoid table.  In the original monoid the
products therefore equal the same outside element.  Since the
row maps of $p,q$ are constant, their composites on the right
with the row map of $z$ retain their respective constant values.
The displayed equality forces those values to agree.  The column
maps already agree because their augmentation operators are
nonzero.  Thus the profiles are equal in this last case too.
This step does not infer an equality of distinct constant maps
from equality of their zero augmentation operators.

The output
\path{code/outputs/d_six_square_zero_ideals.json} records all
ideal bases, quotient coordinates, spanning relations and actual
nonzero-product witnesses.  It proves the asserted spanning
equalities in $628$ cases.  Precisely eight ideals have no such
spanning list; they are retained below, not discarded.
The independent checker
\path{code/b_verify_six_square_zero_ideals.py} reconstructs all
$636$ corners and verifies their annihilation of $V$, their
two-sided ideal identities, and every literal spanning and
product relation.  It reproduces the same $628$ successes and
eight exceptions: the successful witnesses consist of $317$
basis-zero relations, $208$ differences with both operators
nonzero, and $103$ nonzero right-product repairs.  Its output is
\path{code/outputs/b_six_square_zero_ideals_verified.json}.

In a successful case the contracted algebra of the faithful
transformation-pair image is a quotient of $A/I$, hence hereditary
by Lemma~\ref{B:six-squarezero-target-modules}.  Its fourth rational
homology is zero by the band--Tor formula.  The map from the
original monoid to its transformation-pair image induces an
isomorphism in rational degree-two homology, since the two
balanced action presentations coincide.  Every degree-two class
pulls back from the image, so all its products vanish.
\end{proof}

\begin{proposition}[The eight remaining affine triangle families]
\label{B:six-squarezero-triangle-affine}
None of the eight remaining faithful triangle cases in
Lemma~\ref{B:six-squarezero-literal-ideals} can give the cohomology
ring of $S^2\times S^2$.
\end{proposition}
\begin{proof}
We describe the complete finite enumeration and its certificates.
Source indices refer to
\path{code/outputs/d_six_quotient_algebras.jsonl.gz}; the associated
ordinary zero-seven census indices are included to make the
identification independent of record order.  The coordinates of
the arrow ideal use the six distinguished nonzero basis elements
$q_0=1,q_1,\ldots,q_5$.
\[
\begin{array}{c|c|c|c}
\text{source}&\text{zero-seven class/opposite}&I&
 (\dim W_1,\dim W_2,\dim W_3)\\\hline
763&423/\text{no}&q_2-q_3-q_4+q_5&(1,1,3)\\
764&423/\text{yes}&q_2-q_3-q_4+q_5&(2,1,2)\\
1261&678/\text{no}&q_1-q_2+q_3&(1,1,3)\\
1262&678/\text{yes}&q_1-q_2+q_3&(2,1,2)\\
4187&2246/\text{no}&q_1-q_4+q_5&(2,1,2)\\
4188&2246/\text{yes}&q_1-q_4+q_5&(1,1,3)\\
4513&2415/\text{no}&q_1-q_4+q_5&(2,1,2)\\
4514&2415/\text{yes}&q_1-q_4+q_5&(1,1,3)
\end{array}
\]
The vertex order in the last column is the intrinsic triangle
order of Lemma~\ref{B:six-squarezero-target-modules}, rather than
the arbitrary order of the lifted idempotents in a data file.

The faithful-module statement is essential to the scope of this
enumeration.  None of the six distinguished basis vectors belongs
to $I$, as the quotient coordinates in the ideal certificate
verify.  Since $W$ is faithful over $A/I$, every outside basis
element consequently has a nonzero column augmentation operator:
its six-point column map is nonconstant.  Cases not satisfying this
faithfulness condition have already been excluded in the module
lemma by rank-one forms, or do not have the required Tor dimensions.
More explicitly, for $A/(a)$ the vertex dimensions are $(1,1,3)$:
if $b=0$, the tensor complex has degree-zero dimension three and
degree-one dimension one; if $c=0$, all balanced forms have rank
at most one.  For $A/(b)$ the dimensions are $(2,1,2)$:
if $a=0$, the degree-one and degree-two Tor dimensions are one
and two, respectively; if $c=0$, all balanced forms again have
rank at most one.  The only ways to lose faithfulness are these
vanishings, since all vertex spaces are nonzero and the two
remaining arrow actions occupy different Peirce blocks.

The column search \path{code/d_triangle_column_csp.cpp} starts
with all $6^6=46656$ functions on six points.  For each outside
basis element $q_i$ it retains the nonconstant functions with
the required number of fixed points
\[
 1+\operatorname{tr}(q_i\mid W)
 =1+\sum_j(\dim W_j)\chi_j(q_i),
\]
and all of its literal power relations.  The identity has its
forced map.  Every nonzero source product imposes the literal
composition equality for column maps, in their contravariant
order; every zero source product requires a constant composite.
Finally the displayed arrow relation is imposed exactly on the
five-dimensional augmentation matrices.

There is one symmetry normalization.  Choose an idempotent source
element and retain one representative of each conjugacy class of
its possible six-point idempotent maps.  These classes are exactly
the partitions of six given by the sizes of the fibers over its
fixed points.  The search branches over every remaining value;
each propagation step merely removes a value having no compatible
value in an explicitly imposed equation.  Thus this normalization
and propagation lose no simultaneous action.  The complete output
sets, before further stabilizer reduction, have sizes shown below.
\[
\begin{array}{c|r|r|r|r}
\text{source}&\text{column actions}&\text{column orbits}&
\text{row lifts}&\text{faithful full cases}\\\hline
763&816&39&1&26\\
764&300&62&2&56\\
1261&240&20&1&14\\
1262&85&27&1&10\\
4187&248&32&1&15\\
4188&96&20&1&14\\
4513&624&78&2&156\\
4514&288&52&3&156
\end{array}
\]
The second orbit reduction uses only the stabilizer of the fixed
idempotent map.  For rows, every nonzero augmentation operator
determines its three-point map uniquely.  For each zero operator
on an outside element the program tries all three constant maps
and retains exactly those satisfying all outside product equations.
The last column includes every resulting row lift.

The independent program
\path{code/a_triangle_column_audit.cpp}, with runner
\path{code/a_triangle_column_audit.py}, uses a different
multiplication-closure search.  It reconstructs the pools and
fiber-partition normalization and reproduces all eight complete
solution sets exactly, totaling $2697$ actions.  Its verified
output is \path{code/outputs/a_triangle_column_audit.json}.
Two further independent linkage checks,
\path{code/a_triangle24_coverage.py} and
\path{code/b_verify_triangle24_coverage.py}, verify all $330$
column orbits against these $2697$ actions, every constant row
lift, and the faithful-module rank filter.  They account for all
$93$ discarded column orbits and match the $237$ retained orbits
and their row lifts to exactly $447$ source certificates.  Their
outputs are \path{code/outputs/a_triangle24_coverage.json} and
\path{code/outputs/b_triangle24_coverage_verified.json}.

The cup calculation \path{code/d_triangle24_cups.py} first
retains the faithful $A/I$ branch by checking that the span of
its augmentation matrices has dimension five.  Low-rank tuples
are merely discarded as outside this branch; no new cup claim
is inferred from their rank.  The retained actions form precisely
the $447$ cases in the last column.  For each pair of actions the
full twenty-four-element monoid table is explicit: outside
products use the source table when nonzero; when that source
product is zero, the constant composite row and column maps
specify the unique band element.  Products involving the band
are the two actions or rectangular-band multiplication.  All
associativity equations are checked.
These are the full source tables, with six separate outside
labels and eighteen separate band labels.  Outside labels are
retained even when their transformation profiles agree.
Consequently the $447$ certificates cover nonfaithful source
actions as well; they are not tables obtained by replacing a
source with its smaller faithful image.  No finite-homotopy
assertion about such an image is required anywhere in this
enumeration.

For each monoid the program constructs an exact integral
projective resolution prefix through degree four, two integral
balanced forms forming a rational basis, and their affine normalized
two-cocycles.  It gives a literal module-linear comparison from
the projective prefix to the normalized bar resolution.  For all
four ordered products of the two degree-two classes, evaluation
on the image of \emph{every} degree-four projective generator is
zero.  Thus all four product cochains are identically zero on
this resolution; the conclusion is stronger than testing a few
homology cycles.  All comparisons, tables and cocycles are saved
in the $447$ compressed certificates named in
\path{code/outputs/d_triangle24_cups.json}.

The standalone checker \path{code/b_verify_triangle24_cups.py}
imports neither the resolution producer nor its bar-comparison
routine.  It verifies every table, exactness and integral
saturation of the projective prefix, every comparison equation
and its module translates, the full rational degree-two basis
by its pairing with tensor cycles, and products on a saturated
degree-four cycle basis.  All $447$ certificates pass, with the
same counts as the table, as recorded in
\path{code/outputs/b_triangle24_cups_verified.json}.
No conclusion about an uncomputed resolution tail is used.
Every faithful case has zero degree-two products, proving the
proposition.
\end{proof}

\begin{corollary}\label{B:six-squarezero-product-exclusion}
The $458$ square-zero row-module exceptions in
Proposition~\ref{B:six-row-census} cannot occur for a
twenty-four-element monoid model of $S^2\times S^2$.
\end{corollary}
\begin{proof}
Lemma~\ref{B:six-squarezero-target-modules} either makes all
degree-two forms rank one, in which case
Lemma~\ref{B:three-row-rank-one-products} kills every product,
or gives one of the faithful arrow quotients.  The literal ideal
gate and the eight retained affine families exhaust those
quotients by Lemma~\ref{B:six-squarezero-literal-ideals} and
Proposition~\ref{B:six-squarezero-triangle-affine}.
\end{proof}


\begin{proposition}\label{prop:B-product-model}
There is a $25$-element monoid of type $S^2\times S^2$ with an
explicit projective resolution of length four.
\end{proposition}
\begin{proof}
Take $L_0\times L_0$.  Nerves preserve products and realization of
simplicial sets preserves finite products, so its classifying space is
$BL_0\times BL_0\simeq S^2\times S^2$.
Tensor the two length-two right-module resolutions in
Proposition~\ref{prop:B-band-wedge} over $\mathbb Z$ and take the
total complex.  Each tensor of projective summands is projective over
$\mathbb Z[L_0\times L_0]$, since it is the image of the tensor
of their idempotent projections on free modules.  The augmented
complexes are exact with free abelian groups, and hence split as
complexes of abelian groups; their tensor total complex resolves
$\mathbb Z$.  Its abelian ranks are
$(4,20,41,40,16)$ and its maps are explicitly
$d(x\otimes y)=dx\otimes y+(-1)^{\deg x}x\otimes dy$.
After tensoring with the trivial module, the K\"unneth computation
gives $H_0=\mathbb Z$, $H_2=\mathbb Z^2$, $H_4=\mathbb Z$,
and zero otherwise.  The product identification also verifies the
nonzero cross-product, which additive homology alone would not verify.
\end{proof}

\begin{theorem}\label{B:product-minimum25}
The least order of a finite monoid whose classifying space is
homotopy equivalent to $S^2\times S^2$ is exactly $25$.
\end{theorem}
\begin{proof}
Proposition~\ref{prop:B-product-model} supplies the explicit
$25$-element model $L_0\times L_0$ and its complete projective
resolution.  For the lower bound,
Corollary~\ref{B:product-lower24} leaves only order twenty-four.
The two-row/column and three-by-three obstructions exclude those
minimal ideal shapes.  A nontrivial maximal group in any remaining
minimal ideal would require at least
$3\cdot4\cdot2=24$ ideal elements, in addition to the identity
outside the ideal.  Hence the minimal ideal is a rectangular band.
By Corollary~\ref{B:five-outside-minimal-band} at least five
elements lie outside it, so it has at most nineteen elements.
The three-by-four, three-by-five, and four-by-four product
obstructions leave, up to opposite, only a $3\times6$ band.
Its contracted algebra has dimension six.

Proposition~\ref{B:three-by-six-order24-triangular} reduces the
row augmentation action to the full triangular two-by-two algebra.
The band--Tor formula then gives augmentation dimensions
$(\dim W,\dim V)=(5,2)$ and Tor dimensions $2,0,1,0,\ldots$.
The complete reduction in Proposition~\ref{B:six-row-census}
leaves only the $124$ finite-projective-dimension row cases and
the $458$ square-zero cases.  They are excluded respectively by
Proposition~\ref{B:six-pd2-actual-bases} and
Corollary~\ref{B:six-squarezero-product-exclusion}.
Thus no order-twenty-four model exists, proving the lower bound.
\end{proof}

\begin{remark}\label{rem:B-unresolved-minima}
All six minimum-order questions in Part~B4 are settled.  Their
exact values are $7$ for $S^2\vee S^2$, $9$ for $S^4$,
$13$ for $M(\mathbb Z/2,2)$, $17$ for $M(\mathbb Z/3,2)$,
$23$ for $\CP^2$, and $25$ for $S^2\times S^2$.
Each upper bound has an explicit monoid and a complete homotopy
identification; each matching lower bound records its structural
and exhaustive finite-calculation dependencies above.
\end{remark}

\subsection{Aperiodic endpoints and zigzags of homomorphisms}
\begin{proposition}\label{prop:B-zigzag}
If discrete monoids $M,N$ have homotopy-equivalent classifying spaces,
there is a finite zigzag of unital monoid homomorphisms from $M$ to
$N$, every arrow of which induces a homotopy equivalence of
classifying spaces.  The intermediate discrete monoids need not be
finite.  In particular, a finite monoid whose classifying space has
finite homotopy type admits such a zigzag to a finite aperiodic monoid.
\end{proposition}
\begin{proof}
We give a strict-map construction to distinguish the claim from a
statement merely about equivalent spaces.  Let $F$ be the free-monoid
functor on pointed sets (the basepoint is the unit), with forgetful
functor $U$.  The standard augmented simplicial resolution
$Q_nM=(FU)^{n+1}M$ has faces from the counit and degeneracies from
the unit.  Its underlying augmented simplicial set has the extra
degeneracy furnished by the unit, so $|Q_\bullet M|\to M$ is a
weak equivalence of topological monoids.  Every level is free.

For a set $S$, the map from its free monoid to its free group induces
a classifying-space equivalence.  Indeed, the classifying space of the
free monoid is the realization of the free category on the bouquet of
$S$ loops; collapsing the simplices recording composite paths gives
that graph.  Equivalently, its universal cover is the directed path
category on the Cayley tree of the free group, which contracts to that
tree.  Its fundamental group is the free group on $S$, so this
comparison is precisely its aspherical classifying map.
Consequently degreewise group completion produces a simplicial group
$G_M$ and a map
$|Q_\bullet M|\to|G_M|$ inducing a classifying-space equivalence.
This follows by taking the realization of the levelwise-equivalent
bisimplicial nerves.  The unit inclusions and degeneracies are
cofibrations, so ordinary realization has the required homotopy
invariance.  We have obtained strict topological-monoid maps
\[
 M\ \longleftarrow\ |Q_\bullet M|\ \longrightarrow\ |G_M|
\]
that are equivalences after $B$, and the analogous zigzag for $N$.

The equivalence $BM\simeq BN$ gives an isomorphism between
$\overline W G_M$ and $\overline W G_N$ in the homotopy category
of reduced simplicial sets.  Kan's loop-group theorem, in the precise
form of Goerss--Jardine, Proposition~V.6.3 and Corollary~V.6.4
\cite{GoerssJardine1999B},
therefore gives a finite zigzag of weak equivalences of simplicial
groups between $G_M$ and $G_N$.  To make the last assertion explicit,
replace each group cofibrantly.  All simplicial groups are fibrant,
so the isomorphism in their homotopy category is represented by an
actual map between the two cofibrant replacements.  That map is a
weak equivalence, as are the replacement maps.  Realization gives
strict topological-monoid maps connecting the displayed zigzags,
all equivalences after classifying spaces.

To return to discrete monoids, use Fiedorowicz's functor
$T(A)=F_{\mathrm{Fie}}(A)^\delta$ on topological monoids, where
$\delta$ means that the underlying monoid is given the discrete
topology.  Theorem~3.4 of his 1984 paper \cite{Fiedorowicz1984B} provides natural maps
\[
 A\longrightarrow F_{\mathrm{Fie}}(A)
       \longleftarrow F_{\mathrm{Fie}}(A)^\delta
\]
which both induce classifying-space equivalences.  Naturality and
two-out-of-three imply that $T$ sends every map in the preceding
zigzag to a homomorphism inducing a classifying-space equivalence.
For a discrete endpoint $M$, the natural monoid map $M\to T(M)$
itself induces an equivalence: its composite with the change-of-topology
map has this property, and so does that change-of-topology map.
The same applies to $N$.  Appending these endpoint maps proves the
first assertion.  All the resulting classifying spaces are CW
realizations, so their weak equivalences are homotopy equivalences.

For the final assertion, Theorem~\ref{A:finite-dimension} makes the
universal group trivial, so $BM$ is simply connected.
A finite CW model for it is therefore covered
by Proposition~\ref{prop:B-finite-realization}, which supplies a
finite aperiodic $N$ with $BN\simeq BM$.  Apply the first assertion.
\end{proof}

\begin{remark}\label{rem:B-finite-intermediate}
The literal endpoint requirement in Part~B5 is met by this proposition.  Requiring
every intermediate monoid to be finite is a stronger assertion, and
the free resolution, group completion and discretization construction
above does not prove it.  It must not be inferred by replacing the
words ``classifying spaces are equivalent'' with ``finite monoids are
connected by equivalences''.
\end{remark}

